% Pratiques pour éviter la contamination et l’infection — SVTEEHB 3ème — Cours signé OnBuch+
% Programme officiel MINESEC. Compiler avec : tectonic N05-pratiques-pour-eviter-contamination-infection.tex
\newcommand{\DOCMATIERE}{SVTEEHB}
\newcommand{\DOCNIVEAU}{3ème}
\newcommand{\DOCMODULE}{SVTEEHB – classe de 3ème}
\newcommand{\DOCLECON}{N05}
\newcommand{\DOCTITRE}{Pratiques pour éviter la contamination et l’infection}
% OnBuch+ — préambule « cours signés » ENRICHI (compilé avec tectonic / XeLaTeX)
% Système visuel « Pop » d'OnBuch+ : fond crème (jamais blanc), bordures encre
% épaisses, ombres « dures » décalées, titres Archivo Black en capitales.
% Version MATHÉMATIQUES : graphes (pgfplots/tikz), boîtes pédagogiques
% (définition, propriété, méthode, attention, le savais-tu, exercice résolu,
% exercice, TP, bilan), page de garde et signature OnBuch+ ancrée en bas.
%
% Macros attendues dans main.tex AVANT \input{preamble} :
%   \DOCMATIERE  \DOCNIVEAU  \DOCMODULE  \DOCLECON (numéro)  \DOCTITRE
\documentclass[11pt]{article}

\usepackage{fontspec}
\usepackage[french]{babel}
\usepackage[a4paper,margin=2cm,top=3.4cm,bottom=2.8cm,headheight=1.4cm,footskip=1.3cm]{geometry}
\usepackage{amsmath,amssymb,amsfonts}
\usepackage{mathtools} % \xmapsto, \coloneqq... souvent utilisés par les modèles
\usepackage{cancel} % simplifications barrées (bilans de forces)
% \abs{z} (module, valeur absolue) et \norm{u} (norme), utilisés par les modèles
\providecommand{\abs}{}\let\abs\relax\DeclarePairedDelimiter{\abs}{\lvert}{\rvert}
\providecommand{\norm}{}\let\norm\relax\DeclarePairedDelimiter{\norm}{\lVert}{\rVert}
\usepackage[table]{xcolor} % [table] : fournit aussi \rowcolors (alternance de lignes)
\usepackage{pagecolor}
\usepackage{tikz}
\usetikzlibrary{arrows.meta,positioning,calc,shapes.geometric,shapes.misc,shapes.arrows,decorations.pathmorphing,decorations.pathreplacing,decorations.markings,patterns,shadows,mindmap,trees,fit,backgrounds,matrix,intersections,angles,quotes}
\usepackage{pgfplots}
\usepackage[version=4]{mhchem} % équations nucléaires \ce{^{235}_{92}U}
\usepackage[european,siunitx]{circuitikz} % schémas électriques
\pgfplotsset{compat=1.18}
\usepgfplotslibrary{fillbetween,groupplots} % aires entre courbes (calcul intégral)
\usepackage{enumitem}
\usepackage{fancyhdr}
% [most] charge toutes les bibliothèques tcolorbox (listings, minted, raster,
% documentation...) et gonfle inutilement la mémoire TeX sur un document long
% (~300 boîtes) : on ne charge que ce qui est réellement utilisé.
\usepackage[skins,breakable]{tcolorbox}
\usepackage{graphicx}
\usepackage{colortbl}
\usepackage{booktabs}
\usepackage{tabularx}
\usepackage{multirow}
\usepackage{diagbox} % case d'en-tête diagonale des tableaux à double entrée
% Colonnes extensibles centrée / gauche / droite (C, L, R), souvent utilisées par les modèles
\newcolumntype{C}{>{\centering\arraybackslash}X}
\newcolumntype{L}{>{\raggedright\arraybackslash}X}
\newcolumntype{R}{>{\raggedleft\arraybackslash}X}
\usepackage{array}
\usepackage{multicol}
\usepackage{siunitx}
% parse-numbers=false : accepte aussi \SI{2,5 \times 10^{-3}}{...}
\sisetup{locale=FR,output-decimal-marker={,},group-minimum-digits=4,parse-numbers=false}
\DeclareSIUnit\ohm{\ensuremath{\symup{\Omega}}} % U+2126 absent de la police
\DeclareSIUnit\division{div} % réglages d'oscilloscope (ms/div, V/div)
\DeclareSIUnit\tr{tr} % tours (vitesse de rotation, ex. \SI{1500}{\tr\per\minute})
\DeclareSIUnit\molar{M} % concentration molaire (ex. \SI{3}{\molar}, \SI{1}{\micro\molar})
\DeclareSIUnit\an{an} % année (le modèle confond parfois avec \year, un compteur TeX) — ex. \SI{5}{\kilo\meter\per\an}
\DeclareSIUnit\FCFA{FCFA} % franc CFA (ex. \SI{500}{\FCFA\per\kilo\gram})
\DeclareSIUnit\degreeBrix{\textdegree Brix} % teneur en sucre d'un sirop/jus (réfractométrie)
\DeclareSIUnit\month{mois} % le modèle utilise parfois \month, un compteur TeX, comme unité de durée
\DeclareSIUnit\microliter{\micro\liter} % le modèle écrit parfois \microliter en un seul token
\DeclareSIUnit\million{millions} % dénombrement (ex. \SI{15}{\million\per\mL} spermatozoïdes)
\usepackage{qrcode}
% \degree : commande fréquemment utilisée par les modèles pour un angle ou une
% température en dehors de \SI{}{} (non fournie par les paquets chargés).
\providecommand{\degree}{^{\circ}}
% \qed (fin de démonstration), utilisé par les modèles sans amsthm
\providecommand{\qed}{\hfill$\square$}
% \barre : double barre des valeurs interdites dans un tableau de variations (tkz-tab)
\providecommand{\barre}{\ensuremath{\|}}
% environnement proof (amsthm non chargé) utilisé par les modèles
\ifdefined\proof\else\newenvironment{proof}[1][Démonstration]{\par\noindent\textit{#1.}\ }{\qed\par}\fi
\usepackage{lastpage}
\usepackage{hyperref}
\hypersetup{hidelinks}

% ============================================================
% Palette « Pop » OnBuch+ — mêmes hex que la classe P (plus_ui.dart)
% ============================================================
\definecolor{poporange}{HTML}{F59321}
\definecolor{popdark}{HTML}{E07A0C}
\definecolor{popcream}{HTML}{FFF6E8}
\definecolor{popink}{HTML}{1C1714}
\definecolor{poppurple}{HTML}{7A5AE0}
\definecolor{popgreen}{HTML}{2E9E6A}
\definecolor{poppink}{HTML}{E0457B}
\definecolor{popblue}{HTML}{2F7DE1}
\definecolor{popgold}{HTML}{C98A12}
\definecolor{popmuted}{HTML}{8A7F76}
\definecolor{popline}{HTML}{EADFD2}
\definecolor{popgray}{HTML}{8A7F76}
\colorlet{popgrey}{popgray}
% Alias tolérants (le modèle nomme parfois les couleurs différemment)
\colorlet{popwhite}{white}
\colorlet{popblack}{popink}
\colorlet{popred}{poppink}
\colorlet{popyellow}{popgold}
% Teintes claires (fonds de tableaux/encadrés) — le modèle nomme parfois
% les couleurs avec un suffixe L (ex. popblueL) sans les avoir définies.
\colorlet{poporangeL}{poporange!20}
\colorlet{popdarkL}{popdark!20}
\colorlet{poppurpleL}{poppurple!20}
\colorlet{popgreenL}{popgreen!20}
\colorlet{poppinkL}{poppink!20}
\colorlet{popblueL}{popblue!20}
\colorlet{popgoldL}{popgold!20}
\colorlet{popredL}{popred!20}
\colorlet{popgrayL}{popgray!20}
% Teintes claires pour fonds de tableaux / zones
\colorlet{popblueL}{popblue!10!white}
\colorlet{poporangeL}{poporange!14!white}
\colorlet{popgreenL}{popgreen!12!white}
\colorlet{poppurpleL}{poppurple!12!white}
\colorlet{poppinkL}{poppink!12!white}
\colorlet{popgoldL}{popgold!14!white}

\pagecolor{popcream}

% ============================================================
% Typographie
% ============================================================
\defaultfontfeatures{Ligatures=TeX}
\setmainfont{PlusJakartaSans-Medium.ttf}[Path=../../../fonts/,BoldFont=PlusJakartaSans-Bold.ttf]
% Maths en sans-serif (Fira Math) avec chiffres et lettres droites de Plus
% Jakarta, pour que les formules se fondent dans le texte.
\usepackage[math-style=ISO,bold-style=ISO]{unicode-math}
\setmathfont{FiraMath-Regular.otf}[Path=../../../fonts/]
\setmathfont{PlusJakartaSans-Medium.ttf}[Path=../../../fonts/,range={up/num,up/{Latin,latin},\mathup}]
\usepackage{icomma} % virgule décimale sans espace en mode math
% Caractères mathématiques Unicode absents des polices de texte (Plus Jakarta,
% Archivo Black) : rendus par leur équivalent mathématique, en texte comme en
% formule — sinon ils s'affichent en carré vide (ex. « Arithmétique dans ℤ »).
\usepackage{newunicodechar}
\newunicodechar{ℝ}{\ensuremath{\mathbb{R}}}
\newunicodechar{ℤ}{\ensuremath{\mathbb{Z}}}
\newunicodechar{ℂ}{\ensuremath{\mathbb{C}}}
\newunicodechar{ℕ}{\ensuremath{\mathbb{N}}}
\newunicodechar{ℚ}{\ensuremath{\mathbb{Q}}}
\newunicodechar{𝔻}{\ensuremath{\mathbb{D}}}
\newunicodechar{α}{\ensuremath{\alpha}}
\newunicodechar{β}{\ensuremath{\beta}}
\newunicodechar{γ}{\ensuremath{\gamma}}
\newunicodechar{δ}{\ensuremath{\delta}}
\newunicodechar{ε}{\ensuremath{\varepsilon}}
\newunicodechar{θ}{\ensuremath{\theta}}
\newunicodechar{λ}{\ensuremath{\lambda}}
\newunicodechar{μ}{\ensuremath{\mu}}
\newunicodechar{σ}{\ensuremath{\sigma}}
\newunicodechar{τ}{\ensuremath{\tau}}
\newunicodechar{φ}{\ensuremath{\varphi}}
\newunicodechar{ψ}{\ensuremath{\psi}}
\newunicodechar{ω}{\ensuremath{\omega}}
\newunicodechar{Σ}{\ensuremath{\Sigma}}
% majuscules produites par \MakeUppercase dans les titres (α-aminés -> Α)
\newunicodechar{Α}{\ensuremath{\alpha}}
\newunicodechar{Β}{\ensuremath{\beta}}
\newunicodechar{Γ}{\ensuremath{\Gamma}}
\newunicodechar{Δ}{\ensuremath{\Delta}}
\newunicodechar{Ε}{\ensuremath{\varepsilon}}
\newunicodechar{Θ}{\ensuremath{\Theta}}
\newunicodechar{Λ}{\ensuremath{\Lambda}}
\newunicodechar{Μ}{\ensuremath{\mu}}
\newunicodechar{Τ}{\ensuremath{\tau}}
\newunicodechar{Φ}{\ensuremath{\Phi}}
\newunicodechar{Ψ}{\ensuremath{\Psi}}
\newunicodechar{Ω}{\ensuremath{\Omega}}
\newunicodechar{↦}{\ensuremath{\mapsto}}
\newunicodechar{∈}{\ensuremath{\in}}
\newunicodechar{∉}{\ensuremath{\notin}}
\newunicodechar{∧}{\ensuremath{\wedge}}
\newunicodechar{⇔}{\ensuremath{\Leftrightarrow}}
\newunicodechar{⟺}{\ensuremath{\Longleftrightarrow}}
\newunicodechar{⇒}{\ensuremath{\Rightarrow}}
\newunicodechar{⟹}{\ensuremath{\Longrightarrow}}
\newunicodechar{∘}{\ensuremath{\circ}}
\newunicodechar{∪}{\ensuremath{\cup}}
\newunicodechar{∩}{\ensuremath{\cap}}
\newunicodechar{≡}{\ensuremath{\equiv}}
\newunicodechar{⊂}{\ensuremath{\subset}}
\newunicodechar{⊥}{\ensuremath{\perp}}
\newunicodechar{∃}{\ensuremath{\exists}}
\newunicodechar{∀}{\ensuremath{\forall}}
\newunicodechar{∛}{\ensuremath{\sqrt[3]{\,}}}
\newunicodechar{∜}{\ensuremath{\sqrt[4]{\,}}}
\newunicodechar{⃗}{\ensuremath{{}^{\to}}}
\newunicodechar{ⁿ}{\ifmmode^{n}\else\textsuperscript{\ensuremath{n}}\fi}
\newunicodechar{ˣ}{\ifmmode^{x}\else\textsuperscript{\ensuremath{x}}\fi}
\newunicodechar{ᵇ}{\ifmmode^{b}\else\textsuperscript{\ensuremath{b}}\fi}
\newunicodechar{ʳ}{\ifmmode^{r}\else\textsuperscript{\ensuremath{r}}\fi}
\newunicodechar{ᵘ}{\ifmmode^{u}\else\textsuperscript{\ensuremath{u}}\fi}
\newunicodechar{ʸ}{\ifmmode^{y}\else\textsuperscript{\ensuremath{y}}\fi}
\newunicodechar{ᵃ}{\ifmmode^{a}\else\textsuperscript{\ensuremath{a}}\fi}
\newunicodechar{ᵉ}{\ifmmode^{e}\else\textsuperscript{\ensuremath{e}}\fi}
\newunicodechar{ᵏ}{\ifmmode^{k}\else\textsuperscript{\ensuremath{k}}\fi}
\newunicodechar{ᵖ}{\ifmmode^{p}\else\textsuperscript{\ensuremath{p}}\fi}
\newunicodechar{ᴺ}{\ifmmode^{N}\else\textsuperscript{\ensuremath{N}}\fi}
\newunicodechar{ᶜ}{\ifmmode^{c}\else\textsuperscript{\ensuremath{c}}\fi}
\newunicodechar{ʰ}{\ifmmode^{h}\else\textsuperscript{\ensuremath{h}}\fi}
\newunicodechar{ᵠ}{\ifmmode^{q}\else\textsuperscript{\ensuremath{q}}\fi}
\newunicodechar{ₙ}{\ifmmode_{n}\else\textsubscript{\ensuremath{n}}\fi}
\newunicodechar{ₐ}{\ifmmode_{a}\else\textsubscript{\ensuremath{a}}\fi}
\newunicodechar{ᵢ}{\ifmmode_{i}\else\textsubscript{\ensuremath{i}}\fi}
\newunicodechar{ᵣ}{\ifmmode_{r}\else\textsubscript{\ensuremath{r}}\fi}
\newunicodechar{ₖ}{\ifmmode_{k}\else\textsubscript{\ensuremath{k}}\fi}
\newunicodechar{ᵦ}{\ifmmode_{\beta}\else\textsubscript{\ensuremath{\beta}}\fi}
\newunicodechar{ₓ}{\ifmmode_{x}\else\textsubscript{\ensuremath{x}}\fi}
\newfontfamily\popdisplay{ArchivoBlack-Regular.ttf}[Path=../../../fonts/]
\newfontfamily\poplogo{SpaceGrotesk-Bold.ttf}[Path=../../../fonts/]
\newfontfamily\popsemi{PlusJakartaSans-SemiBold.ttf}[Path=../../../fonts/]
\newfontfamily\popxbold{PlusJakartaSans-ExtraBold.ttf}[Path=../../../fonts/]
\newfontfamily\popxboldls{PlusJakartaSans-ExtraBold.ttf}[Path=../../../fonts/,LetterSpace=3.0]

\setlist[enumerate]{leftmargin=1.7em,itemsep=5pt,topsep=4pt}
\setlist[itemize]{leftmargin=1.7em,itemsep=4pt,topsep=4pt,label={\color{poporange}\textbullet}}
\setlength{\parindent}{0pt}
\setlength{\parskip}{5pt}
\renewcommand{\arraystretch}{1.25}

% pgfplots : style Pop par défaut
\pgfplotsset{
  every axis/.append style={axis line style={popink,line width=1pt},tick style={popink},
    grid style={popline,line width=0.5pt},label style={font=\small},tick label style={font=\footnotesize}},
  popcurve/.style={poporange,line width=1.8pt,no markers,smooth},
}
% Même style disponible dans un \draw TikZ simple (hors pgfplots)
\tikzset{popcurve/.style={poporange,line width=1.8pt}}
\tikzset{popgrid/.style={popline,very thin}}
\colorlet{popcurve}{poporange} % le nom du style est parfois utilisé comme couleur

% ============================================================
% Pill / squiggle
% ============================================================
\newcommand{\poppill}[3]{%
  \tikz[baseline=-0.55ex]\node[draw=popink,line width=1.2pt,rounded corners=2.6pt,%
    fill=#2,inner xsep=6.5pt,inner ysep=3pt]{%
    \popxboldls\fontsize{7}{7}\selectfont\color{#3}\MakeUppercase{#1}};%
}
\newcommand{\popsquiggle}[1]{%
  \tikz[baseline=-0.4ex]{
    \draw[#1,line width=2.6pt,line cap=round]
      plot [smooth] coordinates {(0,0.09) (0.34,0.01) (0.68,0.21) (1.02,0.01) (1.36,0.10)};
  }%
}
% Mot-clé surligné (marqueur orange sous le texte)
\newcommand{\cle}[1]{{\popxbold\color{popdark}#1}}

% ============================================================
% En-tête / pied
% ============================================================
\pagestyle{fancy}
\fancyhf{}
\renewcommand{\headrulewidth}{0pt}
\renewcommand{\footrulewidth}{0pt}
\lhead{\raisebox{-0.15ex}{\poplogo\fontsize{13}{13}\selectfont\color{popink}OnBuch\color{poporange}+}}
\rhead{\poppill{\DOCMATIERE\ \textbullet\ \DOCNIVEAU\ \textbullet\ Leçon \DOCLECON}{white}{popink}}
\lfoot{\popxbold\fontsize{7.6}{7.6}\selectfont\color{popmuted}\MakeUppercase{Cours signé OnBuch+}\ \textbullet\ {\popsemi\DOCTITRE}}
\rfoot{\tikz[baseline=-0.6ex]\node[rounded corners=8.5pt,fill=popink,minimum height=17pt,minimum width=17pt,inner xsep=5pt,inner sep=0]{\popxbold\fontsize{8}{8}\selectfont\color{popcream}\thepage\,/\,\pageref*{LastPage}};}

% ============================================================
% Titres
% ============================================================
\newcommand{\chaptitre}[2]{%
  \begin{tcolorbox}[enhanced,colback=poporange,colframe=popink,boxrule=2.6pt,arc=11pt,
      drop shadow=popink,left=15pt,right=15pt,top=12pt,bottom=11pt,before skip=3pt,after skip=18pt]
    \poppill{#2}{popink}{popcream}\par\vspace{9pt}
    {\raggedright\hyphenpenalty=10000\exhyphenpenalty=10000\popdisplay\fontsize{22}{24}\selectfont\color{popcream}\MakeUppercase{#1}\par}
  \end{tcolorbox}
}
% Section numérotée : pastille orange avec le numéro + titre Archivo
\newcounter{popsec}
\newcommand{\coursec}[1]{%
  \stepcounter{popsec}\setcounter{popsub}{0}%
  \par\vspace{16pt}\noindent
  \begin{minipage}{\linewidth}
  \tikz[baseline=-0.7ex]\node[draw=popink,line width=1.6pt,fill=poporange,rounded corners=4pt,minimum size=22pt,inner sep=0]{\popdisplay\fontsize{12}{12}\selectfont\color{popcream}\thepopsec};%
  \hspace{8pt}{\popdisplay\fontsize{15}{17}\selectfont\color{popink}\MakeUppercase{#1}}\par
  \vspace{4pt}\noindent{\color{poporange}\rule{36pt}{3.6pt}}
  \end{minipage}\par\nopagebreak\vspace{8pt}
}
\newcounter{popsub}
\newcommand{\courssub}[1]{%
  \stepcounter{popsub}%
  \par\vspace{9pt}\noindent{\popxbold\fontsize{11.5}{13}\selectfont\color{popink}{\color{poporange}\thepopsec.\thepopsub}\hspace{6pt}#1}\par\nopagebreak\vspace{4pt}
}

% ============================================================
% Boîtes pédagogiques — même recette : fond blanc, bordure épaisse colorée,
% ombre dure encre, badge plein. Titre optionnel [#1].
% ============================================================
\tcbset{popbase/.style={breakable,enhanced,colback=white,boxrule=2.3pt,arc=9pt,
  shadow={3pt}{-3pt}{0pt}{popink},left=14pt,right=14pt,top=11pt,bottom=11pt,before skip=11pt,after skip=15pt,
  fontupper=\popsemi\fontsize{10.6}{14.5}\selectfont\color{popink},
  fontlower=\popsemi\fontsize{10.6}{14.5}\selectfont\color{popink}}}
% Coche dessinée (le glyphe ✓ n'existe pas dans les polices du document)
\newcommand{\popcheck}[1]{\tikz[baseline=-0.5ex]\draw[#1,line width=1.6pt,line cap=round,line join=round](-0.13,0.0)--(-0.04,-0.09)--(0.13,0.1);}
\providecommand{\coche}{\popcheck{popgreen}}
\providecommand{\resultat}[1]{\par\smallskip\noindent\fcolorbox{popgreen}{popgreenL}{\parbox{\dimexpr\linewidth-2\fboxsep-2\fboxrule\relax}{\textbf{Résultat.} #1}}\par\smallskip}
\newcommand{\popboxhead}[3]{\poppill{#1}{#2}{white}\ifx\relax#3\relax\else\hspace{6pt}{\popxbold\fontsize{10.6}{12}\selectfont\color{popink}#3}\fi\par\vspace{7pt}}

\newtcolorbox{aretenir}[1][]{popbase,colframe=popgreen,before upper={\popboxhead{À retenir}{popgreen}{#1}}}
\newtcolorbox{exemplebox}[1][]{popbase,colframe=poporange,before upper={\popboxhead{Exemple}{poporange}{#1}}}
\newtcolorbox{definition}[1][]{popbase,colframe=popblue,before upper={\popboxhead{Définition}{popblue}{#1}}}
\newtcolorbox{propriete}[1][]{popbase,colframe=poppurple,before upper={\popboxhead{Propriété}{poppurple}{#1}}}
\newtcolorbox{methode}[1][]{popbase,colframe=popgold,before upper={\popboxhead{Méthode}{popgold}{#1}}}
\newtcolorbox{attention}[1][]{popbase,colframe=poppink,before upper={\popboxhead{Attention}{poppink}{#1}}}
\newtcolorbox{experience}[1][]{popbase,colframe=popblue,colback=popblueL,before upper={\popboxhead{Expérience}{popblue}{#1}}}
% « Le savais-tu ? » : carte violette pleine (contexte, culture, Cameroun)
\newtcolorbox{savaistu}[1][]{popbase,colback=poppurple,colframe=popink,
  fontupper=\popsemi\fontsize{10.4}{14.2}\selectfont\color{white},
  before upper={\poppill{Le savais-tu\,?}{white}{poppurple}\ifx\relax#1\relax\else\hspace{6pt}{\popxbold\color{white}#1}\fi\par\vspace{7pt}}}
% Exercice résolu : énoncé en haut, \tcblower puis la solution sur fond crème
\newcounter{popexo}\newcounter{popexr}
\newtcolorbox{exoresolu}[1][]{popbase,colframe=poporange,colbacklower=popcream,
  segmentation style={popink,line width=1.2pt,dashed},
  before upper={\stepcounter{popexr}\popboxhead{Exercice résolu \thepopexr}{poporange}{#1}},
  before lower={\poppill{Solution}{popink}{popcream}\par\vspace{6pt}}}
% Exercice (non résolu en place) — la correction est dans la partie Corrigés
\newtcolorbox{exercice}[1][]{popbase,colframe=popink,
  before upper={\stepcounter{popexo}\popboxhead{Exercice \thepopexo}{popink}{#1}}}
% Corrigé
\newtcolorbox{corrige}[1][]{popbase,colframe=popgreen,colback=popgreenL,
  before upper={\popboxhead{Corrigé}{popgreen}{#1}}}
% Objectifs / prérequis / bilan
\newtcolorbox{objectifs}{popbase,colframe=popink,colback=poporangeL,
  before upper={\popboxhead{Objectifs de la leçon}{popink}{}}}
\newtcolorbox{prerequis}{popbase,colframe=popink,colback=popblueL,
  before upper={\popboxhead{Prérequis}{popblue}{}}}
\newtcolorbox{bilan}{popbase,colframe=popink,colback=white,boxrule=2.8pt,
  before upper={\popboxhead{Fiche bilan}{poporange}{L'essentiel à connaître pour l'examen}}}
% Figure encadrée avec légende
\newtcolorbox{popfigure}[1][]{enhanced,breakable=false,colback=white,colframe=popline,boxrule=1.4pt,arc=8pt,
  left=10pt,right=10pt,top=10pt,bottom=8pt,before skip=10pt,after skip=12pt,halign=center,
  lower separated=false,halign lower=center,
  fontlower=\popsemi\fontsize{9}{11.5}\selectfont\color{popmuted},
  #1}
\newcommand{\legende}[1]{\par\vspace{6pt}{\popsemi\fontsize{9}{11.5}\selectfont\color{popmuted}#1}}

% ============================================================
% Signature OnBuch+
% ============================================================
\newcommand{\poplogoinline}[1]{{\poplogo\fontsize{#1}{#1}\selectfont\color{popink}OnBuch\color{poporange}+}}
\newcommand{\signatureonbuch}{%
  \begin{tcolorbox}[enhanced,colback=popink,colframe=popink,boxrule=2.2pt,arc=10pt,
      drop shadow=poporange,left=14pt,right=14pt,top=10pt,bottom=10pt,before skip=0pt,after skip=0pt]
    \tikz[baseline=-0.7ex]\node[circle,fill=popgreen,draw=popcream,line width=1.3pt,minimum size=22pt,inner sep=0]{\popcheck{white}};%
    \hspace{9pt}\begin{minipage}[c]{0.62\linewidth}
      {\popxbold\fontsize{12}{13}\selectfont\color{popcream}Signé\ \textbullet\ {\poplogo OnBuch\color{poporange}+}}\par\vspace{2pt}
      {\raggedright\popsemi\fontsize{8}{10.5}\selectfont\color{popline}\MakeUppercase{Équipe pédagogique OnBuch+}\\\MakeUppercase{Conforme au programme officiel MINESEC}\par}
    \end{minipage}\hfill
    {\poplogo\fontsize{22}{22}\selectfont\color{popcream}OnBuch\color{poporange}+}
  \end{tcolorbox}%
}

% ============================================================
% Page de garde
% #1 titre · #2 matière · #3 niveau · #4 module · #5 n° leçon · #6 durée ·
% #7 description / objectifs (texte ou itemize)
% ============================================================
\newcommand{\pagegarde}[7]{%
  \thispagestyle{empty}
  \newgeometry{a4paper,margin=1.7cm,top=1.6cm,bottom=1.4cm}%
  \begin{tikzpicture}[remember picture,overlay]
    \fill[poporange!18] ([xshift=-3.2cm,yshift=-3.6cm]current page.north east) circle (4.6cm);
    \fill[poppurple!14] ([xshift=2.4cm,yshift=7.6cm]current page.south west) circle (3.8cm);
  \end{tikzpicture}%
  {\poplogo\fontsize{30}{30}\selectfont\color{popink}OnBuch\color{poporange}+}\hfill
  \raisebox{0.6ex}{\poppill{Cours signé \textbullet\ Premium}{popink}{popcream}}\par
  \vspace{1pt}\popsquiggle{poppurple}\par
  \vspace{8pt}
  \begin{tcolorbox}[enhanced,colback=poporange,colframe=popink,boxrule=2.8pt,arc=13pt,
      drop shadow=popink,left=18pt,right=18pt,top=12pt,bottom=13pt,after skip=10pt]
    \poppill{#2 \textbullet\ #3}{popink}{popcream}\hspace{5pt}\poppill{#4}{white}{popink}\par\vspace{8pt}
    {\popdisplay\fontsize{14}{14}\selectfont\color{popink}LEÇON #5}\par\vspace{4pt}
    {\raggedright\hyphenpenalty=10000\exhyphenpenalty=10000\popdisplay\fontsize{25}{27}\selectfont\color{popcream}\MakeUppercase{#1}\par}
    \vspace{8pt}
    {\popxbold\fontsize{9.5}{9.5}\selectfont\color{popink}\MakeUppercase{Durée conseillée : #6}}
  \end{tcolorbox}
  \begin{tcolorbox}[enhanced,colback=white,colframe=popink,boxrule=2.2pt,arc=10pt,
      drop shadow=popink,left=16pt,right=16pt,top=10pt,bottom=10pt,after skip=8pt]
    \poppill{Dans cette leçon}{poppurple}{white}\par\vspace{5pt}
    \popsemi\fontsize{9.3}{12.6}\selectfont\color{popink}#7
  \end{tcolorbox}
  \noindent\begin{minipage}[t]{0.487\linewidth}
    \begin{tcolorbox}[enhanced,colback=popgreen,colframe=popink,boxrule=2.2pt,arc=10pt,
        drop shadow=popink,left=11pt,right=11pt,top=10pt,bottom=10pt,height=3.1cm,valign=center]
      \tcbox[colback=white,colframe=popink,boxrule=1.3pt,arc=5pt,boxsep=0pt,nobeforeafter,
        left=3pt,right=3pt,top=3pt,bottom=3pt,valign=center]{\qrcode[height=1.9cm]{https://play.google.com/store/apps/details?id=com.luvvix.onbuch&hl=fr}}%
      \hspace{8pt}\begin{minipage}[c]{0.5\linewidth}
        {\popdisplay\fontsize{11}{12.5}\selectfont\color{white}\MakeUppercase{Télécharge OnBuch}}\par\vspace{4pt}
        {\raggedright\popsemi\fontsize{8.4}{11}\selectfont\color{white}Gratuit \textbullet\ Google Play\\Tuteur IA, annales, concours, résultats\par}
      \end{minipage}
    \end{tcolorbox}
  \end{minipage}\hfill
  \begin{minipage}[t]{0.487\linewidth}
    \begin{tcolorbox}[enhanced,colback=poppurple,colframe=popink,boxrule=2.2pt,arc=10pt,
        drop shadow=popink,left=11pt,right=11pt,top=10pt,bottom=10pt,height=3.1cm,valign=center]
      \tikz[baseline=-0.7ex]\node[circle,fill=white,draw=popink,line width=1.3pt,minimum size=1.6cm,inner sep=0]{\popdisplay\fontsize{24}{24}\selectfont\color{poppurple}+};%
      \hspace{8pt}\begin{minipage}[c]{0.6\linewidth}
        {\popdisplay\fontsize{11}{12.5}\selectfont\color{white}\MakeUppercase{Passe à OnBuch+}}\par\vspace{4pt}
        {\raggedright\popsemi\fontsize{8.4}{11}\selectfont\color{white}Tous les cours signés, Tuteur IA illimité, fiches PDF — onglet OnBuch+ de l'app\par}
      \end{minipage}
    \end{tcolorbox}
  \end{minipage}
  \vspace{8pt}
  \popcontenu
  \vfill
  \signatureonbuch
  \restoregeometry
  \newpage
}

% Rangée « Ce cours contient » (page de garde)
\newcommand{\poptile}[3]{% #1 couleur #2 gros texte #3 légende
  \begin{tcolorbox}[enhanced,colback=#1,colframe=popink,boxrule=2pt,arc=9pt,shadow={2.5pt}{-2.5pt}{0pt}{popink},
      left=6pt,right=6pt,top=9pt,bottom=9pt,halign=center,valign=center,height=1.75cm,width=\linewidth,nobeforeafter]
    {\popdisplay\fontsize{12.5}{13}\selectfont\color{white}\MakeUppercase{#2}}\par\vspace{3pt}
    {\centering\popsemi\fontsize{7.8}{9.5}\selectfont\color{white}#3\par}
  \end{tcolorbox}}
\newcommand{\popcontenu}{%
  \par\noindent{\popxboldls\fontsize{7.5}{7.5}\selectfont\color{popmuted}CE COURS SIGNÉ CONTIENT}\par\vspace{7pt}
  \noindent\begin{minipage}[t]{0.235\linewidth}\poptile{popblue}{Cours}{complet, illustré, pas à pas}\end{minipage}\hfill
  \begin{minipage}[t]{0.235\linewidth}\poptile{poporange}{Méthodes}{et exercices résolus}\end{minipage}\hfill
  \begin{minipage}[t]{0.235\linewidth}\poptile{poppink}{Exercices}{type examen + corrigés}\end{minipage}\hfill
  \begin{minipage}[t]{0.235\linewidth}\poptile{popgreen}{Fiche bilan}{l'essentiel à retenir}\end{minipage}\par}

% Sommaire
\newtcolorbox{sommaire}{enhanced,colback=white,colframe=popink,boxrule=2.2pt,arc=10pt,
  drop shadow=popink,left=18pt,right=18pt,top=13pt,bottom=13pt,before skip=6pt,after skip=20pt,
  before upper={\poppill{Sommaire}{poppurple}{white}\par\vspace{9pt}\popsemi\fontsize{10.6}{14.5}\selectfont\color{popink}}}

% Signature de fin de cours — poussée tout en bas de la dernière page
\newcommand{\findecours}{%
  \par\vfill\nopagebreak
  \begin{center}\popsquiggle{poporange}\end{center}\vspace{4pt}
  \signatureonbuch
}
\AtBeginDocument{\def\llbracket{\mathopen{[\mkern-3.5mu[}}\def\rrbracket{\mathclose{]\mkern-3.5mu]}}} % intervalles d'entiers (glyphe ⟦ absent de la police)
% Glyphes absents de Fira Math : redéfinis à partir de glyphes disponibles
\AtBeginDocument{%
  \def\setminus{\mathbin{\backslash}}\let\smallsetminus\setminus
  \def\vdots{\mathord{\vbox{\baselineskip3.2pt\lineskiplimit0pt\kern4pt\hbox{.}\hbox{.}\hbox{.}}}}%
  \def\ddots{\mathinner{\mkern1mu\raise6.5pt\hbox{.}\mkern2mu\raise3.6pt\hbox{.}\mkern2mu\raise0.7pt\hbox{.}\mkern1mu}}%
  \def\triangle{\mathord{\scalebox{0.9}{$\symup{\Delta}$}}}%
  \def\nabla{\mathord{\raisebox{\depth}{\scalebox{1}[-1]{$\symup{\Delta}$}}}}%
  \def\checkmark{\popcheck{popdark}}%
  \def\star{\mathbin{\ast}}\def\bigstar{\mathord{\ast}}%
  \def\diamond{\mathbin{\scalebox{0.6}{$\Diamond$}}}%
  \def\bigcup{\mathop{\mathchoice{\raisebox{-0.3em}{\scalebox{1.7}{$\cup$}}}{\scalebox{1.25}{$\cup$}}{\cup}{\cup}}\displaylimits}%
  \def\bigcap{\mathop{\mathchoice{\raisebox{-0.3em}{\scalebox{1.7}{$\cap$}}}{\scalebox{1.25}{$\cap$}}{\cap}{\cap}}\displaylimits}%
  \def\lesseqqgtr{\mathrel{\lessgtr}}\def\bigoplus{\mathop{\oplus}\displaylimits}%
}
\providecommand{\ssi}{si et seulement si} % abréviation fréquente
\AtBeginDocument{\providecommand{\wideparen}{\overparen}} % arc de cercle

%% FIGFIX-BEGIN v5 — correctifs globaux des figures (docs/onbuch-generation/tools/figfix)
\usepackage{adjustbox}\usepackage{etoolbox}\tcbuselibrary{raster}
% toute figure TikZ/pgfplots plus large que la ligne est réduite à la largeur disponible
\BeforeBeginEnvironment{tikzpicture}{\begin{adjustbox}{max width=\linewidth}}
\AfterEndEnvironment{tikzpicture}{\end{adjustbox}}
% légendes pgfplots lisibles par-dessus les courbes
\pgfplotsset{every axis/.append style={legend style={fill=white,fill opacity=0.92,text opacity=1,draw=black!15,inner sep=3pt}}}
% étiquettes d'arêtes (syntaxe edge["..."]) avec fond blanc
\tikzset{every edge quotes/.append style={fill=white,inner sep=1.2pt}}
%% FIGFIX-END
\begin{document}

\pagegarde{Pratiques pour éviter la contamination et l’infection}{SVTEEHB}{3ème}{SVTEEHB – classe de 3ème}{N05}{2 séances (fiche de progression harmonisée)}{Cette leçon étudie les pratiques qui empêchent l'entrée des microbes dans l'organisme : l'asepsie, l'antisepsie et l'utilisation des préservatifs. À partir d'expériences historiques et de situations concrètes de la vie courante, l'élève apprend à distinguer ces pratiques et à les appliquer pour prévenir contaminations et infections.
\vspace{4pt}
\begin{multicols}{2}\small
\begin{itemize}
\item À la fin de cette leçon, je sais définir les notions de contamination, d'infection, d'asepsie et d'antisepsie.
\item À la fin de cette leçon, je sais interpréter les résultats des expériences historiques de Semmelweis et de Lister pour établir l'efficacité des pratiques d'hygiène.
\item À la fin de cette leçon, je sais citer et décrire les principales pratiques d'asepsie (lavage des mains, stérilisation, cuisson des aliments, hygiène corporelle).
\item À la fin de cette leçon, je sais citer et décrire les pratiques d'antisepsie et les antiseptiques d'usage courant.
\item À la fin de cette leçon, je sais distinguer asepsie et antisepsie dans des situations concrètes de la vie courante.
\item À la fin de cette leçon, je sais expliquer le rôle du préservatif dans la prévention des infections sexuellement transmissibles et des grossesses non désirées.
\end{itemize}\end{multicols}}

\begin{sommaire}
\begin{multicols}{2}
\begin{enumerate}
\item Rappels : contamination et infection
\item L'asepsie : empêcher l'arrivée des microbes
\item L'antisepsie : détruire les microbes
\item TP N°6 : pratique de l'asepsie et de l'antisepsie
\item L'utilisation des préservatifs
\item Synthèse et application à la vie courante
\item Activité d'intégration
\item Méthodes et exercices résolus
\item Exercices
\item Corrigés des exercices
\item Fiche bilan
\end{enumerate}
\end{multicols}
\end{sommaire}

\begin{objectifs}
\begin{itemize}
\item À la fin de cette leçon, je sais définir les notions de contamination, d'infection, d'asepsie et d'antisepsie.
\item À la fin de cette leçon, je sais interpréter les résultats des expériences historiques de Semmelweis et de Lister pour établir l'efficacité des pratiques d'hygiène.
\item À la fin de cette leçon, je sais citer et décrire les principales pratiques d'asepsie (lavage des mains, stérilisation, cuisson des aliments, hygiène corporelle).
\item À la fin de cette leçon, je sais citer et décrire les pratiques d'antisepsie et les antiseptiques d'usage courant.
\item À la fin de cette leçon, je sais distinguer asepsie et antisepsie dans des situations concrètes de la vie courante.
\item À la fin de cette leçon, je sais expliquer le rôle du préservatif dans la prévention des infections sexuellement transmissibles et des grossesses non désirées.
\item À la fin de cette leçon, je sais réaliser et justifier les étapes d'une pratique d'asepsie et d'antisepsie (TP N°6).
\item À la fin de cette leçon, je sais rédiger et justifier une démarche, une réponse ou une conclusion à partir de documents.
\end{itemize}
\end{objectifs}

\begin{prerequis}
\begin{itemize}
\item Notion de microbe (bactéries, virus, champignons) vue en classe de 4ème
\item Notions de contamination et d'infection (leçons précédentes de l'unité)
\item Voies de pénétration des microbes dans l'organisme (plaies, voies respiratoires, voies sexuelles)
\item Notion d'infection sexuellement transmissible (IST) et exemples (VIH/SIDA)
\end{itemize}
\end{prerequis}

\newpage

\coursec{Pratiques pour éviter la contamination et l'infection}

\courssub{Situation d'accroche et problématique}

Imagine deux scènes qui se déroulent le même jour à Bafoussam. À l'hôpital de district, une infirmière prépare son matériel : elle se lave les mains minutieusement avec une solution hydroalcoolique, met des gants stériles, désinfecte la peau du patient avant de suturer une plaie. Au marché central, à quelques centaines de mètres, un vendeur de viande découpe sa marchandise sur une table en bois jamais nettoyée, les mains sales, les mouches voltigeant autour. Quelques jours plus tard, deux personnes présentent une plaie similaire : l'une, soignée à l'hôpital, guérit sans encombre ; l'autre, qui s'est fait soigner « à la maison » dans des conditions d'hygiène précaires, voit sa plaie devenir rouge, chaude, douloureuse, suintant du pus, avec de la fièvre.

\begin{aretenir}[Problématique]
Pourquoi certains gestes simples suffisent-ils à empêcher les microbes de pénétrer dans l'organisme, tandis que leur absence entraîne des complications graves ? Comment se protéger efficacement des contaminations, y compris lors des rapports sexuels ?
\end{aretenir}

Pour répondre, nous devons d'abord distinguer deux notions souvent confondues : la \cle{contamination} et l'\cle{infection}.

\courssub{Contamination et infection : définitions et distinction}

\begin{experience}[Observation clinique : une plaie qui se complique]
\textbf{Mise en situation.} Un enfant tombe et s'écorche le genou sur une terre latéritique (sol rouge, riche en microbes). Deux cas de figure :
\begin{itemize}
    \item \textbf{Cas A :} La plaie est lavée à l'eau et au savon, désinfectée, couverte d'un pansement propre. Elle guérit en quelques jours.
    \item \textbf{Cas B :} La plaie n'est pas lavée, on y met de la terre « pour sécher le sang ». Au bout de 48 h, le genou est gonflé, rouge, chaud, douloureux, du pus s'écoule, l'enfant a de la fièvre (38,5 °C).
\end{itemize}
\textbf{Question.} Que s'est-il passé dans le cas B qui ne s'est pas passé dans le cas A ?
\end{experience}

\begin{definition}[Contamination]
La \textbf{contamination} est la \textbf{pénétration des micro-organismes pathogènes (bactéries, virus, champignons microscopiques) à l'intérieur de l'organisme} par une porte d'entrée.
\end{definition}

\begin{definition}[Infection]
L'\textbf{infection} est le \textbf{développement (multiplication) de ces micro-organismes dans l'organisme} et l'\textbf{apparition de troubles cliniques} (signes locaux : rougeur, chaleur, douleur, œdème, pus ; signes généraux : fièvre, fatigue, maux de tête).
\end{definition}

\begin{attention}[Piège fréquent : confondre contamination et infection]
\emph{Contamination} $\neq$ \emph{Infection}.
\begin{itemize}
    \item On peut être \textbf{contaminé sans être encore infecté} : les microbes viennent d'entrer, mais ne se sont pas encore multipliés en nombre suffisant pour provoquer des symptômes (période d'incubation, ou élimination rapide par les défenses naturelles).
    \item L'infection est la \textbf{conséquence} d'une contamination qui n'a pas été bloquée ou contrôlée à temps.
    \item Au BEPC, on attend que vous précisiez : « La contamination est l'entrée ; l'infection est le développement et les troubles. »
\end{itemize}
\end{attention}

\courssub{Les agents de la contamination}

Les agents responsables sont des \cle{micro-organismes} (microbes) invisibles à l'œil nu. Trois grands groupes sont à retenir pour le programme de 3ème :

\begin{center}
\begin{tabularx}{\linewidth}{|l|X|X|}\hline
\textbf{Groupe} & \textbf{Exemples d'agents} & \textbf{Maladies illustratives (Cameroun)} \\ \hline
\textbf{Bactéries} & \textit{Staphylococcus aureus}, \textit{Streptocoques}, \textit{Salmonella}, \textit{Vibrio cholerae} & Infections des plaies, angines, fièvre typhoïde, choléra \\ \hline
\textbf{Virus} & VIH, Virus de l'hépatite B (VHB), Virus de la grippe, Rougeole, COVID-19 & Sida, Hépatite B, Grippe, Rougeole, Covid-19 \\ \hline
\textbf{Champignons microscopiques (Levures, Moisissures)} & \textit{Candida albicans}, \textit{Aspergillus} & Mycoses cutanées, candidose buccale/vaginale, aspergillose \\ \hline
\end{tabularx}
\end{center}

\begin{savaistu}[Une poignée de porte : réservoir à microbes]
Des études menées dans des lieux publics (hôpitaux, écoles, marchés) montrent qu'une poignée de porte peut héberger \textbf{plusieurs milliers de bactéries par cm²} déposées par les mains (Staphylocoques, Colibacilles fécaux). Se laver les mains après les toilettes et avant de manger n'est pas une simple politesse : c'est la première barrière contre la contamination digestive.
\end{savaistu}

\courssub{Les portes d'entrée des microbes dans l'organisme}

L'organisme humain est une forteresse. Les microbes ne peuvent pas traverser une peau intacte et saine. Ils doivent trouver une \textbf{porte d'entrée} (voie de pénétration). En voici les quatre principales :

\begin{popfigure}
\begin{tikzpicture}[scale=0.9, every node/.style={font=\small}]
    % Corps humain simplifié
    \draw[popdark, line width=1.5pt] 
        (0,0) ellipse (2.5 and 4.5) % tronc
        (0,4.5) circle (1.2)       % tête
        (-2.5,1) .. controls (-3.5,2) and (-3.5,-1) .. (-2.5,-1) % bras G
        (2.5,1) .. controls (3.5,2) and (3.5,-1) .. (2.5,-1)   % bras D
        (-1.2,-4.5) .. controls (-1.5,-6.5) and (-0.5,-6.5) .. (-0.8,-4.5) % jambe G
        (1.2,-4.5) .. controls (1.5,-6.5) and (0.5,-6.5) .. (0.8,-4.5);   % jambe D

    % Portes d'entrée avec flèches et étiquettes
    % 1. Plaie (bras)
    \draw[popred, very thick, ->] (-4.2, 1.5) -- (-2.8, 1.2);
    \node[popred, align=left, anchor=east] at (-4.4, 1.5) {\textbf{1. Plaies, blessures}\\(coupure, égratignure,\\brûlure, piqûre d'insecte)};

    % 2. Voies respiratoires (nez/bouche)
    \draw[popblue, very thick, ->] (-1.5, 5.2) -- (-0.5, 5.0);
    \node[popblue, align=left, anchor=east] at (-2.0, 5.5) {\textbf{2. Voies respiratoires}\\(nez, bouche)\\Air poussiéreux, postillons,\\gouttelettes (grippe, TB, COVID)};

    % 3. Voies digestives (bouche)
    \draw[popgreen, very thick, ->] (1.5, 5.2) -- (0.5, 5.0);
    \node[popgreen, align=right, anchor=west] at (2.0, 5.5) {\textbf{3. Voies digestives}\\(bouche, intestin)\\Aliments/eaux souillés\\(choléra, typhoïde, amibiase)};

    % 4. Voies urinaires et génitales
    \draw[poporange, very thick, ->] (0.5, -3.5) -- (1.5, -3.5);
    \node[poporange, align=left, anchor=west] at (1.7, -3.5) {\textbf{4. Voies urinaires}\\\& \textbf{génitales}\\Rapports sexuels non protégés\\(VIH, Hépatite B, IST),\\(cathéter, accouchement)};

    % Légende interne
    \node[align=center, fill=popcream, draw=popline, rounded corners, inner sep=4pt] at (0, -5.5) {
        \textbf{Portes d'entrée principales des microbes}
    };
\end{tikzpicture}
\legende{Schéma stylisé du corps humain montrant les quatre principales portes d'entrée des micro-organismes pathogènes : 1. Peau lésée (plaies), 2. Voies respiratoires, 3. Voies digestives, 4. Voies urinaires et génitales.}
\end{popfigure}

\begin{propriete}[Règle des portes d'entrée]
\begin{itemize}
    \item \textbf{Peau et muqueuses intactes} = barrière physique et chimique efficace (pH acide de la peau, cils vibratiles, mucus, flore commensale).
    \item \textbf{Peau lésée ou muqueuses exposées} = porte ouverte pour les microbes.
    \item La contamination nécessite \textbf{la rencontre} entre un agent pathogène, une porte d'entrée et un hôte sensible.
\end{itemize}
\end{propriete}

\courssub{De la contamination à l'infection : conséquence logique}

\begin{methode}[Raisonnement : pourquoi empêcher la contamination ?]
\begin{enumerate}
    \item \textbf{Observation :} Une plaie non désinfectée s'infecte (pus, fièvre) ; une plaie désinfectée guérit.
    \item \textbf{Hypothèse :} Les microbes entrés par la plaie (contamination) se multiplient et causent les troubles (infection).
    \item \textbf{Expérience (modèle) :} On met en culture deux milieux nutritifs identiques. Milieu A : ensemencé avec des microbes (simule contamination). Milieu B : stérile (témoin). Incubation 24 h à 37 °C.
    \item \textbf{Résultat :} Milieu A : trouble, colonies visibles (multiplication). Milieu B : clair, stérile.
    \item \textbf{Interprétation :} Une fois à l'intérieur (contamination), les microbes trouvent nutriments, température, humidité $\rightarrow$ multiplication exponentielle $\rightarrow$ infection.
    \item \textbf{Conclusion :} \textbf{Pour éviter l'infection, il faut empêcher la contamination}, c'est-à-dire \textbf{bloquer l'entrée des microbes} aux portes d'entrée.
\end{enumerate}
\end{methode}

C'est toute la logique de la \textbf{prophylaxie} (prévention) : agir \emph{avant} l'installation du microbe.

\courssub{Deux grandes catégories de pratiques préventives}

Face aux portes d'entrée, l'humanité a développé deux stratégies complémentaires, que nous allons détailler dans les sections suivantes :

\begin{center}
\begin{tabularx}{\linewidth}{|l|X|X|}\hline
\textbf{Stratégie} & \textbf{Principe} & \textbf{Exemples courants} \\ \hline
\textbf{ASEPSIE} & \textbf{Empêcher l'arrivée des microbes} (créer/maintenir un milieu stérile ou aseptique). & Stérilisation du matériel chirurgical, lavage des mains chirurgical, champs opératoires stériles, port de gants/masque/charte, conservation des aliments (appertisation), \textbf{préservatifs} (barrière physique). \\ \hline
\textbf{ANTISEPSIE} & \textbf{Détruire ou inhiber les microbes déjà présents} sur une surface vivante (peau, muqueuse) ou inerte. & Désinfection d'une plaie (Bétadine®, Dakin®, eau oxygénée), antisepsie cutanée avant injection/prise de sang, lavage des mains antiseptique (solution hydroalcoolique), traitement de l'eau (chlore, Aquatabs®). \\ \hline
\end{tabularx}
\end{center}

\begin{attention}[Distinction cruciale Asepsie / Antisepsie]
\begin{itemize}
    \item \textbf{Asepsie} = « sans septicémie » = \textbf{empêcher l'entrée}. C'est une démarche de \textbf{propreté stérile} (matériel, champs, air). \emph{Ex :} Stériliser une aiguille à 121 °C pendant 20 min (autoclave).
    \item \textbf{Antisepsie} = « contre la putréfaction » = \textbf{tuer sur place}. C'est l'usage d'un \textbf{produit chimique (antiseptique)} sur du \textbf{vivant} (peau, muqueuse) ou du matériel non stérilisable à chaud. \emph{Ex :} Badigeonner la peau avec de la povidone-iodée (Bétadine®) avant une piqûre.
    \item \textbf{Ne pas confondre avec :} \emph{Désinfection} (sur matériel inerte / surfaces / sols / instruments) et \emph{Stérilisation} (destruction de TOUTES les formes microbiennes, y compris spores, sur du matériel inerte).
\end{itemize}
\end{attention}

\coursec{L'asepsie : empêcher l'arrivée des microbes}

\courssub{Principes et méthodes de l'asepsie}

L'asepsie vise à créer et maintenir un environnement \textbf{exempt de micro-organismes} (milieu stérile ou aseptique) pour empêcher toute contamination. Elle repose sur des barrières physiques et des procédés de stérilisation.

\begin{definition}[Stérilisation]
La \textbf{stérilisation} est l'opération qui permet d'éliminer ou de détruire \textbf{tous les micro-organismes} (formes végétatives \textbf{et} spores) portés par un objet ou un milieu. Un objet stérile ne contient \textbf{aucun germe viable}.
\end{definition}

\begin{propriete}[Principaux procédés de stérilisation (matériel inerte)]
\begin{itemize}
    \item \textbf{Chaleur humide (Autoclave) :} \textbf{121 °C pendant 20 minutes} (ou 134 °C pendant 3 min) sous pression (1,1 bar). \emph{Référence au Cameroun :} autoclaves des blocs opératoires, CSI, laboratoires. C'est la méthode de choix pour le matériel métallique, textiles, verrerie, milieux de culture.
    \item \textbf{Chaleur sèche (Étuve / Poupinel) :} \textbf{180 °C pendant 30 à 60 minutes}. Pour le matériel métallique non rouillable, poudres, huiles, verrerie grasse.
    \item \textbf{Stérilisation par filtration :} Filtres à pores de \textbf{0,22 µm} (retiennent les bactéries). Pour les solutions thermosensibles (sérums, antibiotiques, vaccins).
    \item \textbf{Rayonnements ionisants (Gamma, UV) :} Industriel (gants jetables, seringues, pansements stériles du commerce). UV : désinfection d'air/surfaces (pas stérilisation totale).
    \item \textbf{Ébullition (100 °C, 15-20 min) :} \textbf{Désinfection} seulement (ne tue pas les spores). Utilisée en l'absence d'autoclave (petit matériel, eau de boisson).
\end{itemize}
\end{propriete}

\begin{experience}[TP N°6 (Partie A) : Vérification de l'efficacité de la stérilisation à l'autoclave]
\textbf{Objectif.} Vérifier que l'autoclave détruit bien les microbes (y compris spores).
\textbf{Matériel.} Autoclave, boîtes de Pétri avec gélose nutritive, bandes témoins (spores de \textit{Geobacillus stearothermophilus}), matériel à stériliser (pinces, compresses), incubateur à 37 °C et 55-60 °C.
\textbf{Protocole.}
\begin{enumerate}
    \item Placer le matériel et une bande témoin dans l'autoclave. Lancer un cycle 121 °C / 20 min.
    \item À la sortie : 
    \begin{itemize}
        \item Ouvrir aseptiquement la bande témoin, l'ensemencer sur gélose, incuber à 55-60 °C (température de croissance de \textit{G. stearothermophilus}).
        \item Prélever un échantillon du matériel stérilisé (ex. compresses), l'écraser sur gélose, incuber à 37 °C.
    \end{itemize}
    \item Témoin positif : bande témoin non stérilisée ensemencée (doit pousser). Témoin négatif : gélose non ensemencée (doit rester stérile).
\end{enumerate}
\textbf{Résultats attendus.}
\begin{itemize}
    \item Bande témoin stérilisée : \textbf{pas de croissance} à 55-60 °C $\rightarrow$ spores détruites $\rightarrow$ stérilisation efficace.
    \item Matériel stérilisé : \textbf{pas de colonie} à 37 °C après 24-48 h.
    \item Témoin positif : croissance. Témoin négatif : stérile.
\end{itemize}
\textbf{Interprétation.} L'autoclave (chaleur humide sous pression) assure la stérilisation. La bande témoin à spores est le test biologique de référence.
\end{experience}

\courssub{Asepsie chirurgicale et soins : les gestes barrières}

Au-delà de la stérilisation du matériel, l'asepsie impose une \textbf{conduite opératoire} rigoureuse pour ne pas récontaminer le champ stérile.

\begin{methode}[Règles d'or de l'asepsie au bloc opératoire / soins invasifs]
\begin{enumerate}
    \item \textbf{Préparation de l'opérateur :} Lavage chirurgical des mains (solution hydroalcoolique, 3-5 min, ongles courts, pas de bijoux), séchage stérile, gantage stérile (technique « fermée »), port de la charlotte, masque, surblouse stérile.
    \item \textbf{Préparation du champ opératoire :} Désinfection large de la peau du patient (antisepsie \emph{avant} l'acte asepsique), pose de champs stériles (drapage) ne laissant découvert que le site opératoire.
    \item \textbf{Règle du « stérile ne touche que le stérile » :} Un objet/ gant stérile ne doit toucher que des objets stériles. Tout contact avec un objet non stérile (peau non préparée, table non drapée, face non stérile du gant) \textbf{contamine} le gant/objet $\rightarrow$ il faut le changer.
    \item \textbf{Flux unidirectionnel :} Propre $\rightarrow$ Sale. Ne pas croiser les circuits (matériel sale ne repasse pas par le circuit propre).
    \item \textbf{Durée :} Le champ opératoire reste stérile tant qu'il n'est pas mouillé, percé, ou contacté par du non stérile. Au-delà de 3-4 h, risque de contamination aérienne.
\end{enumerate}
\end{methode}

\begin{savaistu}[Asepsie alimentaire : l'appertisation (Nicolas Appert, 1795)]
Au Cameroun, la conservation des tomates, mangues, jus de fruits, sauces (pâte d'arachide) par \textbf{appertisation} (mise en conserve : remplissage $\rightarrow$ épuisement air $\rightarrow$ fermeture hermétique $\rightarrow$ traitement thermique 100-121 °C $\rightarrow$ refroidissement) est une application industrielle de l'asepsie. Le bocal hermétique \textbf{empêche l'arrivée} des microbes (barrière physique) après destruction de ceux présents à l'intérieur par la chaleur. C'est la base de l'industrie agroalimentaire locale (ex. conserves de la SOCAPALM, CHANANA, ou productions artisanales améliorées).
\end{savaistu}

\coursec{L'antisepsie : détruire les microbes sur le vivant}

\courssub{Principes et agents antiseptiques}

L'antisepsie utilise des \textbf{produits chimiques (antiseptiques)} appliqués sur des \textbf{surfaces vivantes} (peau saine, muqueuses, plaies) pour détruire ou inhiber les micro-organismes. Elle ne vise pas la stérilité totale (impossible sur vivant sans nécrose), mais une \textbf{réduction drastique de la charge microbienne}.

\begin{definition}[Antiseptique]
Substance chimique capable de \textbf{tuer (bactéricide, virucide, fongicide) ou d'inhiber (bactériostatique)} les micro-organismes, \textbf{utilisable sur les tissus vivants} (peau, muqueuses) sans toxicité excessive aux concentrations d'emploi.
\end{definition}

\begin{center}
\begin{tabularx}{\linewidth}{|l|X|X|X|}\hline
\textbf{Famille / Molécule} & \textbf{Mécanisme / Spectre} & \textbf{Indications courantes (3ème)} & \textbf{Précautions / Contre-indications} \\ \hline
\textbf{Povidone-iodée (Bétadine®)} & Iode libéré : large spectre (bactéries, virus, champignons, spores). Action rapide. & \textbf{Référence :} Antisepsie cutanée pré-opératoire, soins de plaies, brûlures. Dérivé : Bétadine® dermique (peau), Bétadine® gynécologique (muqueuses). & \textbf{Contre-indiqué :} Allergie iode, femme enceinte/allaitante (risque hypothyroïdie fœtal), nouveau-né, thyroïdique. Ne pas mélanger avec Dakin/eau oxygénée (inactivation). \\ \hline
\textbf{Chlorhexidine (Hibitane®, Plurexid®)} & Cationique : bactéricide Gram+ \textgreater Gram-, virus enveloppés. Action rémanente (persistance 4-6 h). & Antisepsie cutanée (mains chirurgicales, site opératoire), désinfection matériel (sondes), bain antiseptique. & Irritation possible muqueuses/yeux/oreilles (ototoxicité). Inactivée par savons anioniques, pus, sang. \\ \hline
\textbf{Solutions hydroalcooliques (SHA)} & Alcool (éthanol/iso-propanol 60-70\%) + émollient : dénaturation protéines. Large spectre, séchage rapide. & \textbf{Lavage des mains hygiénique/chirurgical} (friction 30 s / 3-5 min). Antisepsie cutanée rapide (piqûre, prise de sang). & Inflammable. Irritant sur plaies/muqueuses (brûlure). Inefficace sur mains visiblement souillées (laver au savon d'abord). \\ \hline
\textbf{Dakin® (Hypochlorite de sodium 0,1-0,5\%)} & Chlore actif : oxydant puissant, large spectre, bon marché. & Soins de \textbf{plaies sales, nécrotiques, brûlures} (nettoyage + antisepsie). Bain de bouche dilué. & Instable à la lumière/air (préparer frais). Irritant pur (diluer 1/10 pour plaies). Inactivé par matière organique (pus, sang). \\ \hline
\textbf{Eau oxygénée (H₂O₂ 3-10 vol / 1-3\%)} & Peroxyde : oxydation, effet mécanique (mousse = détersion). Bactéricide, virucide limité. & \textbf{Nettoyage plaies} (effet bulles = débridement mécanique), hémostase superficielle. & Instable (flacon opaque). cytotoxique pour fibroblastes (retard cicatrisation si usage prolongé). Ne pas utiliser en profondeur (embolie gazeuse rare). \\ \hline
\textbf{Hexamidine (Hexomédine®)} & Diamidine : bactériostatique Gram+, fongistatique. & Soins cordon ombilical nouveau-né, plaies superficielles, dermatoses infectées. & Spectre étroit (peu actif Gram-). Usage local uniquement. \\ \hline
\end{tabularx}
\end{center}

\begin{attention}[Règles d'or de l'antisepsie (BEPC)]
\begin{itemize}
    \item \textbf{Nettoyer avant de désinfecter :} La matière organique (sang, pus, saletés) inactive la plupart des antiseptiques. \emph{Protocole :} Nettoyage eau/savon ou sérum physiologique $\rightarrow$ Rinçage $\rightarrow$ Antiseptique.
    \item \textbf{Respecter la concentration et le temps de contact :} Une Bétadine® diluée « à l'œil » ou essuyée trop vite ne marche pas. SHA : friction jusqu'à séchage complet (30 s).
    \item \textbf{Ne pas mélanger les antiseptiques :} Risque d'inactivation (ex. Iode + Dakin/Eau oxygénée = inefficacité) ou de toxicité.
    \item \textbf{Adapter au site :} Peau intacte (SHA, Chlorhexidine, Bétadine dermique), Muqueuses (Bétadine gynéco, Chlorhexidine 0,05\%, Hexamidine), Plaie (Dakin dilué, Bétadine dermique, Eau oxygénée pour nettoyage initial), Nouveau-né (Hexamidine, Chlorhexidine 0,05\% - PAS Bétadine, PAS Alcool pur).
    \item \textbf{Antisepsie $\neq$ Antibiotique :} L'antiseptique est local, large spectre, non sélectif, pas de résistance classique. L'antibiotique est systémique (ou local), sélectif, risque de résistance.
\end{itemize}
\end{attention}

\courssub{TP N°6 (Partie B) : Comparaison de l'efficacité de deux antiseptiques (Méthode de diffusion sur gélose / Disques)}

\begin{experience}[TP N°6 : Pratique de l'asepsie et de l'antisepsie (Protocole type BEPC)]
\textbf{Objectif.} Comparer l'activité antiseptique de la Bétadine® (povidone-iodée) et de la Chlorhexidine sur \textit{Staphylococcus aureus} (souche de référence ATCC 25923 ou isolat clinique sensible).
\textbf{Matériel.} Boîtes de Pétri + gélose Mueller-Hinton (milieu standard antibiogramme/antiseptigramme), suspension bactérienne calibrée (0,5 McFarland $\approx$ 1,5 $\times$ 10\textsuperscript{8} UFC/mL), étaleur stérile, disques de papier filtre stériles (6 mm), pinces stérilisées à la flamme, solutions mères Bétadine® 10\% (1\% iode) et Chlorhexidine 4\%, dilutions sérielles (ex. 1/10, 1/100, 1/1000), incubateur 37 °C, règle.
\textbf{Protocole (asepsie stricte : flamme, manipulation rapide, couvercle entrouvert).}
\begin{enumerate}
    \item \textbf{Ensemençage :} Étaler 100 µL de suspension bactérienne sur toute la surface de 3 boîtes par antiseptique testé + 1 boîte témoin (sans disque). Laisser sécher 5-10 min.
    \item \textbf{Dépôt des disques :} Tremper disques stériles dans chaque dilution (et eau stérile pour témoin disque). Déposer 1 disque par dilution sur la gélose (flamme pince entre chaque). Appuyer légèrement pour adhérence.
    \item \textbf{Incubation :} 18-24 h à 37 °C.
    \item \textbf{Lecture :} Mesurer le \textbf{diamètre des zones d'inhibition (ZI)} en mm (disque + halo clair). Noter si ZI $\ge$ 12-15 mm (seuil souvent utilisé pour « sensible » en antiseptigramme, selon recommandations SF2H/CA-SFM).
\end{enumerate}
\textbf{Résultats types attendus (exemple).}
\begin{center}
\begin{tabular}{|c|c|c|c|c|}\hline
\textbf{Produit} & \textbf{Dilution} & \textbf{Concentration finale approx.} & \textbf{ZI moyen (mm)} & \textbf{Interprétation} \\ \hline
Bétadine® & 1/10 & 0,1\% iode & 25-30 & +++ (Très actif) \\ \hline
Bétadine® & 1/100 & 0,01\% iode & 15-20 & ++ (Actif) \\ \hline
Bétadine® & 1/1000 & 0,001\% iode & 0-8 & - (Inactif) \\ \hline
Chlorhexidine & 1/10 & 0,4\% & 20-25 & +++ \\ \hline
Chlorhexidine & 1/100 & 0,04\% & 12-16 & + \\ \hline
Chlorhexidine & 1/1000 & 0,004\% & 0 & - \\ \hline
Témoin (Eau) & - & - & 0 & - \\ \hline
\end{tabular}
\end{center}
\textbf{Interprétation.}
\begin{itemize}
    \item Les deux antiseptiques sont efficaces aux concentrations d'usage (Bétadine dermique $\approx$ 1\% iode = dilution 1/10 ; Chlorhexidine cutanée $\approx$ 0,5-4\%).
    \item La Bétadine® garde une activité à dilution plus forte (1/100) grâce à l'effet rémanent de l'iode libre.
    \item La Chlorhexidine a une action rémanente sur peau (persistance) absente pour l'iode (inactivé par pus/sang).
    \item \textbf{Conclusion :} Le choix dépend de l'indication (site, matière organique, allergie, rémanence souhaitée).
\end{itemize}
\end{experience}

\coursec{L'utilisation des préservatifs : une asepsie « barrière »}

\courssub{Le préservatif, barrière physique contre la contamination sexuelle}

Les rapports sexuels non protégés constituent une porte d'entrée majeure pour les agents des \textbf{Infections Sexuellement Transmissibles (IST)} : VIH, VHB, VHC, \textit{Neisseria gonorrhoeae} (blennorragie), \textit{Treponema pallidum} (syphilis), \textit{Chlamydia trachomatis}, HPV (cancer col utérus), Herpès, etc. Le préservatif (masculin ou féminin) est la \textbf{seule méthode} offrant une protection simultanée contre la grossesse non désirée \textbf{et} les IST (double protection).

\begin{definition}[Préservatif masculin (externe)]
Gaine fine en \textbf{latex de caoutchouc naturel} (ou polyuréthane/polyisoprène pour allergiques au latex), non poreuse, élastique, lubrifiée, à usage unique, se déroulant sur le pénis en érection \textbf{avant tout contact génital}.
\end{definition}

\begin{definition}[Préservatif féminin (interne)]
Gaine en \textbf{nitrile} (polymère synthétique, non allergisant, plus résistant que le latex), plus large, avec deux anneaux souples (un à l'extrémité fermée pour insertion vaginale, un à l'extrémité ouverte restant à l'extérieur). Inséré dans le vagin \textbf{avant} le rapport.
\end{definition}

\begin{propriete}[Efficacité du préservatif (Données OMS / ONUSIDA)]
\begin{itemize}
    \item \textbf{Utilisation parfaite (théorique) :} Efficacité \textbf{$\ge$ 98\%} pour prévenir la grossesse ; \textbf{$\ge$ 90-95\%} pour réduire la transmission du VIH (études sérodiscordantes).
    \item \textbf{Utilisation réelle (typique) :} Efficacité \textbf{$\approx$ 85-87\%} (grossesse) ; \textbf{$\approx$ 80-90\%} (VIH/IST). L'écart est dû aux \textbf{erreurs d'utilisation} (voir ci-dessous).
    \item \textbf{Protection IST :} Très efficace pour IST par fluides (VIH, gonocoque, chlamydia, trichomonas, hépatites B/C). Moins efficace pour IST par contact cutané-muqueux direct (Herpès, HPV, Syphilis chancre) si la lésion est en dehors de la zone couverte.
\end{itemize}
\end{propriete}

\begin{methode}[Bon usage du préservatif masculin (Étapes clés BEPC)]
\begin{enumerate}
    \item \textbf{Vérifier :} Date de péremption, marque de conformité (NF, CE, ISO 4074), aspect de l'emballage (intact, gonflé = air présent = non percé).
    \item \textbf{Ouvrir :} Déchirer l'emballage sur le côté (pas les dents, pas ciseaux, pas ongles longs).
    \item \textbf{Sens :} Vérifier le sens de déroulement (bord replié vers l'extérieur). Le poser sur le gland \textbf{avant tout contact} (pré-éjaculat = risque VIH/grossesse).
    \item \textbf{Pincer le réservoir :} \textbf{Étape cruciale.} Pincer le bout (réservoir) entre pouce et index pour chasser l'air (évite l'éclatement à l'éjaculation).
    \item \textbf{Dérouler :} De l'autre main, dérouler jusqu'à la base du pénis en érection.
    \item \textbf{Pendant le rapport :} Vérifier qu'il reste en place. Lubrifiant \textbf{uniquement à base d'eau ou silicone} (jamais gras : vaseline, huile, beurre, crème $\rightarrow$ rupture latex en minutes).
    \item \textbf{Après éjaculation :} Retirer \textbf{avant la fin de l'érection}, en tenant le préservatif à la base pour éviter fuite. Nœud, poubelle (pas toilettes = bouchage).
    \item \textbf{Un rapport = Un préservatif.} Jamais de réutilisation. Jamais deux préservatifs superposés (frottement = rupture).
\end{enumerate}
\end{methode}

\begin{attention}[Erreurs fréquentes entraînant l'échec (Rupture / Glissement / Fuite)]
\begin{itemize}
    \item \textbf{Oubli du réservoir (air piégé) :} Cause n°1 de rupture.
    \item \textbf{Lubrifiant gras (vaseline, huile, beurre de karité, crème hydratante) :} Dégrade le latex en \textbf{moins de 5 minutes}.
    \item \textbf{Mise en place trop tardive :} Contact pré-éjaculatoire (risque VIH, grossesse).
    \item \textbf{Retrait trop tardif :} Perte d'érection $\rightarrow$ glissement $\rightarrow$ fuite sperme.
    \item \textbf{Conservation inadaptée :} Portefeuille (chaleur, frottement), boîte à gants voiture (chaleur extrême), poche près du téléphone (chaleur). $\rightarrow$ Vieillissement prématuré, micro-fissures. \textbf{Conserver au frais, sec, à l'abri de la lumière.}
    \item \textbf{Taille inadaptée :} Trop petit $\rightarrow$ rupture ; Trop grand $\rightarrow$ glissement. Existent tailles 49-56 mm (largeur à plat).
\end{itemize}
\end{attention}

\begin{savaistu}[Contexte camerounais : Stratégie « ABC » et accès aux préservatifs]
Au Cameroun, la lutte contre le VIH/Sida et les IST repose sur la stratégie \textbf{ABC} (\textbf{A}bstinence / \textbf{B}e faithful - Fidélité / \textbf{C}ondom - Préservatif) complétée par le dépistage (TDR, autotest) et la PrEP/PEP. Les préservatifs masculins (Marques : \textbf{Prudence}, \textbf{Manix}, \textbf{Generic} gratuits dans les structures publiques) et féminins (\textbf{FC2}) sont distribués gratuitement dans les \textbf{Centre de Dépistage Volontaire (CDV)}, \textbf{Centre de Santé Intégré (CSI)}, \textbf{Hôpitaux}, \textbf{ONG} (CAMNAFAW, ACMS, Care Cameroun) et points de vente sociaux (boutiques, pharmacies, stations-service). La loi protège l'accès aux préservatifs pour les mineurs (pas d'âge minimum légal pour l'achat/obtention gratuite). La circoncision médicale volontaire (CMV) réduit aussi le risque d'acquisition du VIH chez l'homme hétérosexuel d'environ 60\% (recommandée OMS/ONUSIDA dans les pays à forte prévalence comme le Cameroun).
\end{savaistu}

\coursec{Synthèse et application à la vie courante}

\courssub{Tableau récapitulatif : Asepsie vs Antisepsie vs Désinfection vs Stérilisation}

\begin{center}
\begin{tabularx}{\linewidth}{|l|X|X|X|}\hline
\textbf{Concept} & \textbf{Cible (Support)} & \textbf{Objectif} & \textbf{Exemple type 3ème} \\ \hline
\textbf{ASEPSIE} & Environnement, matériel, mains (gants), champs, \textbf{préservatif} & \textbf{Empêcher l'arrivée} (Barrière / Stérilité) & Gants stériles, champ opératoire, autoclave, préservatif, appertisation \\ \hline
\textbf{ANTISEPSIE} & \textbf{Vivant} (Peau saine, muqueuse, plaie) & \textbf{Détruire / Inhiber} sur place & Bétadine® sur peau, SHA mains, Dakin plaie, Chlorhexidine cordon \\ \hline
\textbf{DÉSINFECTION} & \textbf{Inerte} (Sols, surfaces, instruments, eau, linge) & \textbf{Réduire charge microbienne} (pas stérilité) & Sol salle d'op (détergent-désinfectant), eau (chlore/Aquatabs), instruments (bain chimique avant stérilisation) \\ \hline
\textbf{STÉRILISATION} & \textbf{Inerte} (Matériel médical, milieux, conserves) & \textbf{Détruire TOUTES formes} (y compris spores) & Autoclave 121°C/20min, Étuve 180°C/30min, Filtration 0,22µm, Rayons $\gamma$ \\ \hline
\end{tabularx}
\end{center}

\courssub{Application à la vie courante : Check-list prévention}

\begin{itemize}
    \item \textbf{Portes d'entrée cutanées :} Nettoyer toute plaie (eau + savon) $\rightarrow$ Antiseptique (Bétadine/Dakin) $\rightarrow$ Pansement propre (asepsie). Vaccin antitétanique à jour (rappel tous les 10-20 ans, ou 5 ans si plaie sale).
    \item \textbf{Portes d'entrée respiratoires :} Masque si malade / lieu bondé / épidémie (grippe, COVID, TB). Aérer les pièces (renouvellement air). Hygiène respiratoire (toux dans coude, mouchoir jetable).
    \item \textbf{Portes d'entrée digestives :} \textbf{Lavage mains} (savon + eau / SHA) : après toilettes, avant manger/nourrir enfant, après manipuler argent/animaux/déchets. Eau potable (bouillie, filtrée, chlorée, bouteille scellée). Aliments cuits/peelés, conservation au frais, séparation cru/cuit.
    \item \textbf{Portes d'entrée uro-génitales :} \textbf{Préservatif à chaque rapport} (sauf couple fidèle, tous deux testés négatifs, projet grossesse). Dépistage VIH/IST régulier. Hygiène intime douce (eau, pas de douche vaginale). Vaccination HPV (filles 9-14 ans, programme EPI Cameroun), Hépatite B (naissance + 6-10-14 sem).
    \item \textbf{Environnement :} Gestion déchets (tri, poubelles fermées), assainissement (latrines hygiéniques, pas déjection à l'air libre), lutte anti-vectorielle (moustiquaires imprégnées MILDA contre paludisme, destruction gîtes larvaires).
\end{itemize}

\begin{aretenir}[Synthèse BEPC : Ce qu'il faut maîtriser pour l'examen]
\begin{itemize}
    \item \textbf{Contamination} (entrée microbes) $\neq$ \textbf{Infection} (multiplication + troubles). Agents : Bactéries, Virus, Champignons. Portes d'entrée : Peau lésée, Respiratoire, Digestif, Uro-génital.
    \item \textbf{Asepsie} = Empêcher l'arrivée (Stérilisation matériel, Barrières : gants, champs, \textbf{préservatifs}, Appertisation). Règle : « Stérile ne touche que stérile ».
    \item \textbf{Antisepsie} = Détruire sur \textbf{vivant} (Antiseptiques : Bétadine, Chlorhexidine, SHA, Dakin, Eau oxygénée, Hexamidine). Règles : Nettoyer avant, concentration/temps respectés, pas de mélange, adapter au site.
    \item \textbf{Désinfection} = Inerte (sols, surfaces, eau). \textbf{Stérilisation} = Inerte + Spores (Autoclave 121°C/20min).
    \item \textbf{Préservatif} = Asepsie barrière (Latex/Nitrile). Bon usage : Date/Emballage, Réservoir pincé, Avant contact, Lubrifiant eau/silicone, Retrait précoce, Poubelle. Erreurs : Gras, Air, Tardif, Conservation chaleur.
    \item \textbf{Prophylaxie globale :} Hygiène mains, Eau/Aliments sûrs, Vaccins, Préservatifs, Moustiquaires, Assainissement.
\end{itemize}
\end{aretenir}

\begin{exoresolu}[Cas intégrateur : Journée type au CSI de Ngaoundéré]
\textbf{Énoncé.} Au CSI, l'infirmier-chef supervise : (1) Stérilisation des sets de suture à l'autoclave (121°C, 20 min). (2) Préparation d'une césarienne : lavage chirurgical mains (SHA 5 min), gants stériles, champs stériles, antisepsie cutanée patiente (Bétadine dermique large). (3) Soins d'une plaie infectée : nettoyage sérum physiologique, Dakin 1/10, pansement stérile. (4) Conseil planning familial : remise préservatifs masculins + notice illustrée. (5) Nettoyage sol salle d'accouchement (détergent-désinfectant). (6) Traitement eau boisson (pastille chlore Aquatabs® 20 min).

\textbf{Questions.}
\begin{enumerate}
    \item Classer chaque action (1 à 6) dans : Asepsie / Antisepsie / Désinfection / Stérilisation.
    \item Pour l'action (2), expliquer pourquoi on fait \emph{à la fois} asepsie (gants, champs) \emph{et} antisepsie (Bétadine sur patiente).
    \item L'action (6) assure-t-elle une stérilisation de l'eau ? Justifier.
    \item Un couple utilise un préservatif masculin avec de la vaseline comme lubrifiant. Quel risque ? Expliquer le mécanisme.
\end{enumerate}

\tcblower

\textbf{Solution détaillée.}
\begin{enumerate}
    \item (1) \textbf{Stérilisation} (autoclave = destruction spores sur inerte). (2) \textbf{Asepsie} (gants stériles, champs stériles, lavage chirurgical = barrières stériles) \textbf{ET} \textbf{Antisepsie} (Bétadine sur peau vivante patiente). (3) \textbf{Antisepsie} (Dakin sur plaie vivante) + \textbf{Asepsie} (pansement stérile = barrière). (4) \textbf{Asepsie} (Préservatif = barrière physique latex). (5) \textbf{Désinfection} (sol inerte, produit chimique, pas stérilité). (6) \textbf{Désinfection} (eau = milieu inerte/aliment, chlore tue formes végétatives, pas spores).
    \item L'\textbf{asepsie} (gants, champs) protège la patiente des microbes \textbf{extérieurs} (mains équipe, air, instruments). L'\textbf{antisepsie} (Bétadine) détruit la flore cutanée \textbf{endogène} de la patiente (Staphylocoques commensaux) qui pourrait contaminer le site opératoire \emph{de l'intérieur} lors de l'incision. Les deux sont complémentaires : « Zéro germe apporté + Zéro germe résident ».
    \item \textbf{Non.} Aquatabs® (hypochlorite/dichloroisocyanurate) assure une \textbf{désinfection} (potabilisation) : destruction des bactéries, virus, protozoaires (formes végétatives). Les \textbf{spores bactériennes} (ex. \textit{Clostridium}) et certains kystes de parasites résistent au chlore aux doses/temps usuels. L'eau n'est pas stérile (stérilité = 0 germe viable, y compris spores). Pour stériliser l'eau : autoclave ou filtration 0,22 µm (impraticable au quotidien).
    \item \textbf{Risque majeur : Rupture du préservatif.} \textbf{Mécanisme :} La vaseline (corps gras, hydrocarbures) \textbf{solvant le latex} (polymère hydrocarboné). En \textbf{moins de 5 minutes}, le latex perd son élasticité, devient friable, se déchire au moindre étirement. $\rightarrow$ Fuite spermatozoïdes + passage VIH/IST. \textbf{Règle :} Lubrifiant \textbf{uniquement à base d'eau ou silicone} (compatibles latex).
\end{enumerate}
\end{exoresolu}

\begin{exercice}[Entraînement BEPC : Situations variées]
\textbf{Consigne.} Pour chaque situation, identifier la pratique (Asepsie, Antisepsie, Désinfection, Stérilisation) et justifier en une phrase précise (cible + mécanisme).
\begin{enumerate}
    \item Une mère stérilise les biberons de son nourrisson de 2 mois en les faisant bouillir 20 minutes dans une casserole couverte.
    \item Un agent de santé pulvérise une solution de chlore à 0,5\% sur les sols et murs d'une chambre d'isolement Ebola après sortie du patient.
    \item Un chirurgien-dentiste passe ses fraises et pinces à l'autoclave entre chaque patient.
    \item Une jeune fille applique du gel hydroalcoolique sur ses mains avant de manger un sandwich dans un taxi.
    \item Un infirmier nettoie une plaie sale au sérum physiologique, puis applique de la Bétadine® dermique, et pose un pansement stérile.
    \item Un couple utilise un préservatif féminin (FC2) en nitrile pour un rapport sexuel.
\end{enumerate}
\end{exercice}

\begin{corrige}[Exercice : Entraînement BEPC]
\begin{enumerate}
    \item \textbf{Désinfection (thermique).} L'ébullition (100 °C) détruit les formes végétatives mais \textbf{pas les spores} (ex. \textit{Clostridium botulinum} risque toxine). Pour un nourrisson (immunité immature), la stérilisation à l'autoclave (121 °C) ou stérilisateur vapeur/biberons (vapeur 100 °C sous pression légère) est recommandée ; l'ébullition simple est une désinfection améliorée, tolérée si pas d'autoclave.
    \item \textbf{Désinfection.} Support \textbf{inerte} (sols, murs). Produit chimique (chlore) pour réduire charge microbienne environnementale. Pas de stérilité requise ni possible sur surfaces poreuses/grandes surfaces.
    \item \textbf{Stérilisation.} Matériel \textbf{inerte} invasif (contact sang/stérile). Autoclave (chaleur humide 121 °C/20 min) = destruction \textbf{toutes formes} dont spores. Exigence légale/déontologique.
    \item \textbf{Antisepsie.} Support \textbf{vivant} (peau mains saines). SHA (alcool) = destruction flore transitoire. Action rapide, séchage sans rinçage. Pratique hygiène des mains « moment 2 » (avant acte asepsique/alimentaire).
    \item \textbf{Antisepsie} (Bétadine sur plaie vivante) + \textbf{Asepsie} (Pansement stérile = barrière stérile protégeant la plaie de l'extérieur). Nettoyage sérum physiologique = détersion (étape préalable obligatoire).
    \item \textbf{Asepsie (barrière).} Préservatif féminin en \textbf{nitrile} (polymère synthétique) = barrière physique imperméable aux microbes (VIH, spermatozoïdes, IST) et aux fluides, posée sur muqueuses vaginales \textbf{avant} contact. Alternative sans latex (allergies), plus résistante, conduction thermique meilleure.
\end{enumerate}
\end{corrige}

\coursec{L'asepsie : empêcher l'arrivée des microbes}

Dans la section précédente, nous avons distingué la \cle{contamination} (arrivée des microbes sur ou dans l'organisme) de l'\cle{infection} (multiplication et effets pathogènes). La question centrale est maintenant : \textbf{comment empêcher cette contamination ?} C'est tout l'objet de l'\textbf{asepsie}, première ligne de défense avant même que l'infection ne puisse débuter. Observons d'abord ce que l'histoire de la médecine nous enseigne sur la puissance de gestes simples.

\courssub{Définition et principe de l'asepsie}

\begin{definition}[L'asepsie]
L'\textbf{asepsie} est l'ensemble des méthodes et pratiques qui ont pour but d'\textbf{empêcher l'arrivée des microorganismes} sur un organisme vivant (peau, muqueuses, plaie) ou dans un milieu (aliment, eau, instrument chirurgical, milieu de culture). Elle ne tue pas les microbes : elle les \textbf{éloigne} ou les \textbf{empêche de pénétrer}.
\end{definition}

\begin{attention}[Asepsie $\neq$ Antisepsie — Erreur fréquente au BEPC]
Ne confonds jamais les deux termes :
\begin{itemize}
    \item \textbf{Asepsie} : \emph{prévention} de l'arrivée (barrières, propreté, stérilisation du matériel \emph{avant} usage). \textbf{Elle ne détruit pas les microbes.}
    \item \textbf{Antisepsie} : \emph{destruction} ou inhibition des microbes \emph{déjà présents} sur un tissu vivant (peau, muqueuse) à l'aide d'un produit chimique (antiseptique).
\end{itemize}
Au BEPC, une définition qui dit « l'asepsie tue les microbes » est \textbf{fausse} et pénalisée.
\end{attention}

\courssub{Preuve historique : l'expérience de Semmelweis (1847) — Le lavage des mains sauve des vies}

Au milieu du XIX\textsuperscript{e} siècle, à la maternité de l'hôpital général de Vienne, le médecin hongrois \textbf{Ignace Semmelweis} observe une situation dramatique qui va fonder scientifiquement l'asepsie.

\begin{experience}[L'observation de Semmelweis à la maternité de Vienne (1847)]
\textbf{Contexte :} La maternité comporte deux services d'accouchement :
\begin{itemize}
    \item \textbf{Service 1 :} géré par des médecins et des étudiants qui pratiquent des autopsies le matin et accouchent les femmes l'après-midi \textbf{sans se laver les mains} entre les deux activités.
    \item \textbf{Service 2 :} géré par des sages-femmes qui ne font pas d'autopsies.
\end{itemize}
\textbf{Observation (Données chiffrées) :} Le taux de mortalité des accouchées (fièvre puerpérale) est \textbf{environ 18 \%} au Service 1, contre \textbf{environ 2 à 3 \%} au Service 2.

\textbf{Hypothèse de Semmelweis :} Les mains des médecins transportent des « particules cadaveriques » (microbes) des cadavres aux femmes en couches.

\textbf{Expérience (Intervention) :} En mai 1847, Semmelweis impose le \textbf{lavage obligatoire des mains au chlorure de chaux} (solution désinfectante) pour tout le personnel du Service 1 avant chaque examen.

\textbf{Résultat :} Le taux de mortalité au Service 1 chute brutalement pour rejoindre celui du Service 2 (autour de 1-2 \%).
\end{experience}

\begin{popfigure}
\begin{tikzpicture}
\begin{axis}[
    ybar,
    bar width=35pt,
    width=\linewidth,
    height=9cm,
    ymin=0, ymax=20,
    ylabel={Taux de mortalité des accouchées (\%)},
    symbolic x coords={Avant lavage (Service 1), Après lavage (Service 1), Service 2 (Sages-femmes)},
    xtick=data,
    nodes near coords,
    nodes near coords align={vertical},
    ymajorgrids=true,
    grid style=dashed,
    title={Impact du lavage des mains sur la mortalité puerpérale (Semmelweis, 1847)},
    title style={yshift=0.5cm, font=\bfseries},
    enlarge x limits=0.25,
    every axis plot/.append style={fill opacity=0.85}
]
\addplot[popblue, fill=popblue!70] coordinates {(Avant lavage (Service 1),18.3) (Après lavage (Service 1),2.2) (Service 2 (Sages-femmes),2.7)};
\end{axis}
\end{tikzpicture}
\legende{Graphique en barres comparant le taux de mortalité des accouchées avant et après l'instauration du lavage des mains au chlorure de chaux au Service 1, et le taux de référence du Service 2. Données historiques approchées.}
\end{popfigure}

\begin{methode}[Interpréter un résultat expérimental de type « avant / après »]
Pour conclure sur l'efficacité d'une pratique (ici le lavage des mains) :
\begin{enumerate}
    \item \textbf{Identifier les deux situations comparées :} Situation initiale (sans la pratique) vs Situation finale (avec la pratique). Ici : Service 1 avant mai 1847 vs Service 1 après mai 1847.
    \item \textbf{Repérer la variable mesurée :} Le taux de mortalité (\%).
    \item \textbf{Comparer les valeurs :} 18,3 \% $\rightarrow$ 2,2 \% (chute drastique, facteur $\times 8$).
    \item \textbf{Identifier le facteur modifié :} L'introduction du lavage des mains au chlorure de chaux.
    \item \textbf{Conclure :} La modification du facteur (lavage) est la cause probable de la variation observée (baisse mortalité). Le lavage des mains \textbf{empêche l'arrivée des microbes} responsables de la fièvre puerpérale $\rightarrow$ c'est une preuve de l'efficacité de l'asepsie.
\end{enumerate}
\end{methode}

\begin{exemplebox}[Données chiffrées clés à retenir pour le BEPC]
À la maternité de Vienne (1847) :
\begin{itemize}
    \item Mortalité Service 1 (médecins, sans lavage) : \textbf{$\approx$ 18 \%}.
    \item Mortalité Service 1 (après lavage mains chlorure de chaux) : \textbf{$\approx$ 2 \%}.
    \item Mortalité Service 2 (sages-femmes, pas d'autopsies) : \textbf{$\approx$ 2-3 \%}.
\end{itemize}
Le lavage des mains fait passer la mortalité de \textbf{18 \% à moins de 3 \%}, soit une division par \textbf{6 à 9}.
\end{exemplebox}

\begin{savaistu}[Semmelweis, le « sauveur des mères » rejeté par ses pairs]
Malgré des résultats spectaculaires, Semmelweis fut violemment critiqué par l'establishment médical viennois. Ses collègues médecins refusaient d'admettre qu'eux-mêmes étaient la cause de la mort de leurs patientes (« \emph{Les mains d'un gentleman sont propres} »). Il fut écarté de l'hôpital, sombra dans la dépression et mourut en 1865 dans un asile, quelques années avant que les travaux de Pasteur et Lister ne valident définitivement la théorie microbienne. Aujourd'hui, l'OMS considère l'hygiène des mains comme la mesure \textbf{numéro 1} pour prévenir les infections nosocomiales.
\end{savaistu}

\courssub{Les travaux de Lister (1867) : l'asepsie chirurgicale}

Vingt ans après Semmelweis, le chirurgien britannique \textbf{Joseph Lister}, s'appuyant sur les travaux de Pasteur sur les germes atmosphériques, comprend que les infections post-opératoires (gangrène, septicémie) viennent de microbes déposés sur la plaie par l'air, les instruments ou les mains du chirurgien.

Il met en place une \textbf{asepsie rigoureuse au bloc opératoire} :
\begin{itemize}
    \item Pulvérisation d'acide phénique (phénol) dans l'air du bloc (asepsie de l'air).
    \item Lavage des mains et des instruments à l'acide phénique.
    \item Pansements occlusifs imbibés d'antiseptique.
\end{itemize}
\textbf{Résultat :} La mortalité post-opératoire chute de \textbf{45-50 \% à 15 \%} (puis moins avec l'amélioration des techniques). Lister est le père de la \textbf{chirurgie aseptique moderne}. Aujourd'hui, la pulvérisation d'antiseptique dans l'air est abandonnée (toxique), remplacée par la \textbf{stérilisation} (chaleur, autoclave) du matériel et des champs opératoires, le port de \textbf{gants stériles}, \textbf{masque}, \textbf{charle} et \textbf{sur-chaussures}.

\courssub{TP N°6 : Pratique de l'asepsie et de l'antisepsie (Mise en évidence de l'efficacité du lavage des mains)}

Ce TP permet de visualiser la charge microbienne sur les mains et de vérifier l'efficacité du protocole de lavage.

\begin{experience}[TP N°6 : Mise en évidence de l'efficacité du lavage des mains par culture sur gélose]
\textbf{Matériel :} Boîtes de Pétri avec gélose nutritive (PCA ou TS), étiquettes, marqueur, eau, savon liquide, serviette papier stérile, incubateur (ou lieu tiède $\approx$ 30-37 °C).
\textbf{Protocole :}
\begin{enumerate}
    \item \textbf{Repérage :} Diviser le fond de 3 boîtes en 2 secteurs (traits au marqueur sous la boîte). Étiqueter : Boîte A « Avant lavage », Boîte B « Après lavage eau seule », Boîte C « Après lavage savon + eau ».
    \item \textbf{Échantillonnage « Avant » :} Presser le bout des doigts de la main droite (non lavée) sur le secteur correspondant de la Boîte A.
    \item \textbf{Lavage « Eau seule » :} Se laver les mains \textbf{uniquement à l'eau courante} (30 s, sans savon). Sécher sur serviette papier. Presser les doigts sur Boîte B.
    \item \textbf{Lavage « Savon + Eau » :} Se laver les mains \textbf{au savon} en suivant le protocole 6 étapes (30 s minimum). Rincer, sécher sur serviette papier propre. Presser les doigts sur Boîte C.
    \item \textbf{Incubation :} Refermer les boîtes, les retourner (couvercle en bas) et incuber 24 à 48 h à 30-37 °C.
\end{enumerate}
\textbf{Observations attendues :}
\begin{itemize}
    \item Boîte A : Nombreuses colonies (diverses morphologies, couleurs) $\rightarrow$ flore transitoire et résidente abondante.
    \item Boîte B : Colonies encore nombreuses $\rightarrow$ l'eau seule élimine mal les microbes (absence d'action mécanique/détergente).
    \item Boîte C : Très peu de colonies, voire aucune $\rightarrow$ l'action mécanique du savon + rinçage élimine efficacement la flore transitoire et réduit la flore résidente.
\end{itemize}
\textbf{Interprétation :} Le lavage à l'eau seule est inefficace (asepsie incomplète). Le lavage au savon (action tensioactive + mécanique) réalise une asepsie efficace des mains (réduction drastique de la charge microbienne).
\textbf{Conclusion :} L'asepsie des mains nécessite \textbf{savon + frottement 30 s + rinçage + séchage propre}.
\end{experience}

\courssub{Les pratiques d'asepsie au quotidien : du personnel soignant à la collectivité}

L'asepsie ne concerne pas que l'hôpital. Elle s'applique à tous les niveaux : individuel, alimentaire, médical, collectif.

\courssub{Pratiques d'asepsie individuelles}

La barrière la plus simple et la plus efficace est l'\textbf{hygiène des mains}.

\begin{popfigure}
\begin{tikzpicture}[
    node distance=1.1cm and 0.5cm,
    box/.style={rectangle, draw=popdark, thick, fill=popcream, rounded corners, minimum width=3.5cm, minimum height=1.1cm, align=center, font=\small},
    arrow/.style={-Stealth, thick, popdark}
]
\node[box, align=center] (1){1. Mouiller les mains\\à l'eau courante};
\node[box, below=of 1, align=center] (2){2. Savonner\\paumes, dos, doigts};
\node[box, below=of 2, align=center] (3){3. Frotter \textbf{30 s min} :\\paume/paume, paume/dos,\\interlacé, pouces, ongles};
\node[box, below=of 3, align=center] (4){4. Rincer abondamment\\eau courante};
\node[box, below=of 4, align=center] (5){5. Sécher serviette\\propre ou papier\\à usage unique};
\node[box, below=of 5, align=center] (6){6. Fermer robinet\\avec la serviette\\ (ne pas recontaminer)};
\draw[arrow] (1) -- (2);
\draw[arrow] (2) -- (3);
\draw[arrow] (3) -- (4);
\draw[arrow] (4) -- (5);
\draw[arrow] (5) -- (6);
% Side labels
\node[left=0.4cm of 1, text width=2.8cm, align=right, font=\small\bfseries, popblue] {Eau + Savon};
\node[left=0.4cm of 3, text width=2.8cm, align=right, font=\small\bfseries, poporange] {Action mécanique};
\node[left=0.4cm of 6, text width=2.8cm, align=right, font=\small\bfseries, popgreen] {Séchage aseptique};
\end{tikzpicture}
\legende{Les 6 étapes du lavage des mains efficace à l'eau et au savon (recommandation OMS adaptée). L'action mécanique (frottement 30 s) décolle les microbes ; le rinçage les emporte ; le séchage propre évite la récontamination.}
\end{popfigure}

\begin{aretenir}[Moments obligatoires du lavage des mains (Programme 3ème)]
\begin{itemize}
    \item \textbf{Avant :} manger, préparer les repas, soigner une plaie, mettre des lentilles, nourrir un enfant.
    \item \textbf{Après :} être allé aux toilettes, s'être mouché/toussé/éternué, avoir manipulé des déchets, des animaux, de l'argent, des aliments crus, avoir soigné un malade.
\end{itemize}
\end{aretenir}

Autres pratiques individuelles : hygiène corporelle quotidienne (douche), coupe des ongles courts, port de \textbf{gants à usage unique} par le personnel soignant pour tout acte invasif ou contact avec des liquides biologiques (sang, urines, selles).

\courssub{Pratiques d'asepsie alimentaire (Contexte camerounais)}

Au Cameroun, les maladies diarrhéiques (choléra, typhoïde, amibiase) restent une cause majeure de morbidité, souvent liées à une asepsie alimentaire défaillante.

\begin{itemize}
    \item \textbf{Lavage des fruits et légumes :} À l'eau potable, additionnée de quelques gouttes d'eau de Javel (solution à 1 \% de chlore actif) ou vinaigre blanc, puis rincage à l'eau potable. \emph{Exemple :} Laver les tomates, salades, manioc avant consommation crue ou cuisson.
    \item \textbf{Cuisson suffisante :} Température à cœur $\ge$ 70 °C pour détruire la plupart des bactéries pathogènes (Salmonelles, \emph{E. coli}). \emph{Attention :} La cuisson est une \textbf{antisepsie} (destruction), mais elle s'inscrit dans une démarche asepsique globale (empêcher la récontamination après cuisson).
    \item \textbf{Couverture des aliments :} Protège des mouches, poussières, toux. Utiliser des moustiquaires alimentaires ou boîtes hermétiques.
    \item \textbf{Conservation au frais :} Réfrigérateur $\le$ 4 °C (ralentit multiplication) ; congélateur $-18$ °C (arrête multiplication). Ne pas rompre la chaîne du froid.
    \item \textbf{Traitement de l'eau de boisson :} Ébullition 10 min (asepsie thermique), ou chloration (pastilles Aquatabs\textsuperscript{\textregistered}, eau de Javel solution à 1 \% Cl2, 3 gouttes/L, attendre 30 min), ou filtration céramique. \textbf{L'eau de forage ou de source non traitée n'est pas aseptique.}
\end{itemize}

\begin{exoresolu}[Gestion de l'eau dans un quartier de Yaoundé]
\textbf{Énoncé :} Dans un quartier, l'eau du robinet est coupée 3 jours/semaine. Les habitants stockent l'eau dans des bidons de 20 L ouverts. Une épidémie de choléra survient. Proposer 3 mesures d'asepsie pour l'eau de boisson.
\tcblower
\textbf{Solution détaillée :}
\begin{enumerate}
    \item \textbf{Traitement :} Bouillir l'eau 10 min OU chloration (3 gouttes d'eau de Javel à 1 \% Cl2 par litre, attendre 30 min). $\rightarrow$ Détruit \emph{Vibrio cholerae}.
    \item \textbf{Stockage aseptique :} Utiliser des récipients \textbf{propres, fermés (couvercle)}, à robinet (évite d'y tremper les mains/gobelets sales). $\rightarrow$ Empêche récontamination (asepsie).
    \item \textbf{Hygiène du puisage :} Se laver les mains avant de puiser ; utiliser une louche dédiée, propre, accrochée (ne pas poser au sol). $\rightarrow$ Empêche contamination manuportée.
\end{enumerate}
\end{exoresolu}

\courssub{Pratiques d'asepsie médicale (Milieu hospitalier)}

L'objectif : \textbf{zéro microbe} sur le matériel pénétrant dans l'organisme (stérile) ou en contact avec une muqueuse/peau saine.

\begin{center}
\begin{tabularx}{\linewidth}{|l|X|X|}
\hline
\rowcolor{popblueL}
\textbf{Pratique} & \textbf{Principe / Moyen} & \textbf{Exemple d'application} \\
\hline
\textbf{Stérilisation par chaleur humide (Autoclave)} & Vapeur d'eau saturée sous pression : $121$ °C, $1,1$ bar, $20$ min (cycle standard). Détruit \textbf{toutes} formes microbiennes (bactéries, virus, spores). & Instruments chirurgicaux, compresses, champs opératoires, matériel de suture. \\
\hline
\textbf{Stérilisation par chaleur sèche (Étuve / Poupinel)} & Air chaud $160-180$ °C, $1-2$ h. Pour matériel métallique ne rouillant pas, verrerie, poudres. & Pinces, ciseaux, lames de bistouri, seringues en verre (ancien). \\
\hline
\textbf{Matériel à usage unique (Single-use)} & Stérilisé en usine (irradiation $\gamma$, oxyde d'éthylène), emballé stérile. \textbf{Interdiction formelle de réutiliser.} & Seringues, aiguilles, cathéters, gants d'examen, lames de bistouri jetables. \\
\hline
\textbf{Tenue du personnel (Bloc / Soins invasifs)} & Barrières physiques : \textbf{charle stérile}, \textbf{gants stériles}, \textbf{masque chirurgical}, \textbf{calot}, \textbf{sur-chaussures}. & Chirurgie, accouchement, pose de sonde vésicale, perfusion centrale. \\
\hline
\end{tabularx}
\end{center}

\begin{attention}[Stérilisation $\neq$ Désinfection]
\begin{itemize}
    \item \textbf{Stérilisation :} Élimination de \textbf{tous} les microorganismes (y compris spores). Requis pour matériel pénétrant dans le stérile (sang, tissus profonds).
    \item \textbf{Désinfection :} Élimination de la \textbf{plupart} des formes végétatives (bactéries, virus, champignons), \textbf{pas les spores}. Requis pour matériel en contact avec peau saine ou muqueuses non stériles (endoscope, thermomètre, sol).
\end{itemize}
Au BEPC : « Stériliser, c'est détruire tous les microbes, y compris les spores. »
\end{attention}

\courssub{Pratiques d'asepsie collective (Assainissement du milieu)}

L'asepsie ne s'arrête pas au seuil de la maison ou de l'hôpital. L'environnement est un réservoir majeur de microbes.

\begin{itemize}
    \item \textbf{Évacuation des ordures ménagères :} Bacs fermés, collecte régulière, décharge contrôlée (loin des habitations et nappes phréatiques). Évite la prolifération de mouches, rats, cafards (vecteurs mécaniques).
    \item \textbf{Latrines / Toilettes hygiéniques :} Fosse septique étanche ou latrines à fosse améliorée (VIP : \emph{Ventilated Improved Pit}) avec tuyau d'aération moustiquaré. \textbf{Distance minimale 15-30 m des points d'eau (puits, forages).} Empêche la contamination fécale de l'eau de boisson (choléra, typhoïde, poliomyélite, hépatite E).
    \item \textbf{Lutte contre les mouches :} Moustiquaires aux fenêtres, poubelles fermées, élimination des sites de ponte (fèces, déchets organiques en décomposition). La mouche domestique (\emph{Musca domestica}) transporte jusqu'à $10^6$ bactéries sur ses pattes et son proboscis.
    \item \textbf{Assainissement des eaux usées :} Réseaux d'égouts entretenus, stations d'épuration fonctionnelles. Évite les eaux stagnantes (gîtes larvaires moustiques \emph{Anopheles} $\rightarrow$ paludisme ; \emph{Aedes} $\rightarrow$ dengue, chikungunya).
\end{itemize}

\begin{savaistu}[Le choléra au Cameroun : l'asepsie collective en action]
Le Cameroun fait face à des épidémies de choléra récurrentes (Nord, Extrême-Nord, Littoral, Sud-Ouest). La riposte repose sur la \textbf{stratégie « Eau, Hygiène, Assainissement » (EHA)} :
\begin{itemize}
    \item Distribution de pastilles de chlore (Aquatabs\textsuperscript{\textregistered}) pour l'asepsie de l'eau au domicile.
    \item Construction/réhabilitation de latrines VIP dans les écoles et centres de santé.
    \item Sensibilisation au lavage des mains au savon (ou cendres) aux moments critiques.
    \item Enterrement digne et sécurisé des décès (asepsie des corps).
\end{itemize}
C'est l'application concrète, à l'échelle de la population, des principes d'asepsie vus en classe.
\end{savaistu}

\courssub{Synthèse : L'asepsie, première barrière de la défense}

\begin{aretenir}[Bilan : L'asepsie en 3 points clés]
\begin{enumerate}
    \item \textbf{But :} Empêcher l'\textbf{arrivée} des microbes (contamination). \textbf{Pas de destruction.}
    \item \textbf{Preuve historique :} Semmelweis (1847) $\rightarrow$ Lavage mains : mortalité $18\% \rightarrow 2\%$. Lister (1867) $\rightarrow$ Asepsie chirurgicale $\rightarrow$ mortalité divisée par 3.
    \item \textbf{4 niveaux d'action :}
    \begin{itemize}
        \item \textbf{Individuel :} Lavage mains (30 s, savon, moments clés), hygiène corporelle, gants soignant.
        \item \textbf{Alimentaire :} Lavage fruits/légumes, cuisson $\ge 70$ °C, couverture, chaîne du froid, eau traitée (ébullition/chlore).
        \item \textbf{Médical :} Stérilisation (autoclave $121$ °C / étuve $180$ °C), matériel usage unique, tenue stérile (bloc).
        \item \textbf{Collectif :} Ordures gérées, latrines étanches loin de l'eau, lutte anti-mouches/vecteurs, assainissement eaux usées.
    \end{itemize}
\end{enumerate}
\end{aretenir}

\begin{exercice}[Application : Cas pratique d'asepsie à la maison]
Madame M. prépare le repas du soir : poulet grillé, salade de tomates crues, eau de boisson puisée au robinet (coupure annoncée demain). Elle a un enfant de 2 ans qui joue dans la cour.
\begin{enumerate}
    \item Cite 3 gestes d'asepsie qu'elle doit faire \textbf{avant} de cuisiner.
    \item Pour la salade de tomates : quelle pratique d'asepsie alimentaire appliquer ?
    \item Pour l'eau de boisson stockée pour demain : quelles 2 mesures d'asepsie prendre ?
    \item L'enfant tombe, se blesse au genou (plaie superficielle). Quel geste d'asepsie individuel avant de nettoyer la plaie ?
\end{enumerate}
\end{exercice}

\begin{corrige}[Exercice n°1]
\begin{enumerate}
    \item Se laver les mains au savon (30 s) ; attacher ses cheveux / porter un calot ; porter un tablier propre (ou enlever bijoux, montres).
    \item Laver les tomates à l'eau potable (+ quelques gouttes eau de Javel/vinaigre), rincer à l'eau potable, couper sur planche propre, couvrir jusqu'au service.
    \item (1) Traiter l'eau : ébullition 10 min OU chloration (3 gouttes eau de Javel 1\%/L, 30 min attente). (2) Stocker dans récipient \textbf{propre, fermé, à robinet}, hors de portée des enfants/animaux.
    \item Se laver les mains au savon (asepsie individuelle) avant de toucher la plaie ou le matériel de pansement.
\end{enumerate}
\end{corrige}

\textbf{Transition :} Nous venons de voir comment \textbf{empêcher les microbes d'arriver} (asepsie). Mais que faire si la barrière a cédé, si les microbes sont \textbf{déjà là} sur une plaie, une muqueuse, ou nos mains sales ? C'est le rôle de l'\textbf{antisepsie} : \textbf{détruire} ou inhiber ces microbes sur le vivant. C'est l'objet de la section suivante.

\coursec{L'antisepsie : détruire les microbes}

Dans la section précédente, nous avons vu que l'\cle{asepsie} consiste à \emph{empêcher l'arrivée} des microbes : on lave ses mains, on stérilise le matériel, on porte des gants, afin que les microorganismes n'atteignent jamais le corps ou la plaie. Mais que faire lorsque le microbe est \emph{déjà là} ? Un enfant tombe de vélo sur la route de quartier à Yaoundé, son genou est écorché et souillé de poussière : les microbes sont déjà déposés sur la plaie. Il est trop tard pour les empêcher d'arriver ; il faut alors les \emph{détruire}. C'est précisément le rôle de l'\cle{antisepsie}, que nous allons étudier maintenant.

\courssub{Définition de l'antisepsie}

\textbf{Une observation pour commencer.} Dans un centre de santé, l'infirmier qui soigne une écorchure ne se contente pas de la couvrir : il verse d'abord un liquide brun-rouge (la bétadine) ou incolore (l'alcool) sur la plaie. Pourquoi ? Parce que la peau blessée porte déjà des microbes venus de l'air, de la poussière et de la peau elle-même. S'ils restent en place, ils se multiplient et provoquent une \emph{infection} : rougeur, gonflement, douleur, parfois du pus. Le liquide versé a pour but de tuer ces microbes avant qu'ils ne se développent.

\begin{definition}[L'antisepsie]
L'\cle{antisepsie} est l'ensemble des méthodes qui \textbf{détruisent} les microbes ou \textbf{empêchent leur développement} lorsqu'ils sont déjà présents sur la peau, sur les plaies ou sur les objets.
\end{definition}

\begin{definition}[L'antiseptique]
Un \cle{antiseptique} est une substance qui détruit les microbes ou empêche leur développement sur les \textbf{tissus vivants} (peau, plaies, muqueuses). Exemples : l'alcool à 70°, la bétadine, l'eau oxygénée.
\end{definition}

Retiens bien la nuance avec l'asepsie : l'asepsie agit \emph{avant} (elle empêche l'arrivée des microbes), l'antisepsie agit \emph{après} (elle détruit les microbes déjà présents). Les deux pratiques sont complémentaires et se succèdent dans tout soin bien conduit.

\courssub{Les antiseptiques d'usage courant}

Voici les antiseptiques que l'on trouve dans les pharmacies, les centres de santé et les trousses de secours au Cameroun :

\begin{center}
\begin{tabularx}{\linewidth}{|l|X|X|}
\hline
\rowcolor{popblueL} \textbf{Antiseptique} & \textbf{Aspect et usage} & \textbf{Précautions} \\
\hline
Alcool à 70° & Liquide incolore ; désinfecte la peau saine (avant une piqûre) et le pourtour des plaies & Pique sur plaie ouverte ; inflammable ; ne pas utiliser sur les muqueuses \\
\hline
Bétadine (iodée) & Liquide brun-rouge ; s'applique directement sur les plaies & Ne pas mélanger avec les antiseptiques au mercure (Mercurochrome) \\
\hline
Eau oxygénée & Liquide incolore qui mousse au contact de la plaie ; aide à dégager les saletés & Ne pas boucher le flacon hermétiquement après usage ; usage externe uniquement \\
\hline
Solution de Dakin & Eau de Javel diluée ; nettoie les plaies souillées & Usage externe ; ne pas avaler ; se conserve à l'abri de la lumière \\
\hline
Savon antiseptique & Savon contenant une substance antiseptique ; lavage des mains et des plaies & Rincer abondamment à l'eau propre \\
\hline
\end{tabularx}
\end{center}

\begin{attention}[Erreurs fréquentes avec les antiseptiques]
\begin{itemize}
\item \textbf{Ne jamais appliquer un antiseptique sur une plaie sale sans l'avoir lavée.} L'antiseptique ne traverse pas la couche de saleté : les microbes restent vivants en dessous. Le nettoyage à l'eau et au savon \emph{précède toujours} l'application de l'antiseptique.
\item \textbf{Ne pas mélanger certains antiseptiques} entre eux (par exemple bétadine et Mercurochrome) : le mélange peut perdre son efficacité ou devenir irritant. On utilise \emph{un seul} antiseptique à la fois, sauf avis d'un agent de santé.
\item Un antiseptique ne se boit pas et ne s'applique pas dans les yeux : c'est un usage \emph{externe} uniquement.
\end{itemize}
\end{attention}

\begin{exemplebox}[Pourquoi l'alcool à 70° et pas à 90° ?]
On pourrait croire que plus l'alcool est concentré, plus il est efficace. C'est faux ! L'alcool est le plus efficace à \textbf{70°} (c'est-à-dire environ 70 mL d'alcool pour 100 mL de solution). En effet, la présence d'eau dans la solution favorise la \emph{pénétration} de l'alcool à l'intérieur du microbe, qu'il détruit alors de l'intérieur. À 90°, l'alcool trop concentré coagule brutalement les protéines de la surface du microbe et forme comme une « coque » protectrice : le cœur du microbe reste vivant. C'est pourquoi les pharmacies vendent de l'alcool à 70° pour l'antisepsie.
\end{exemplebox}

\courssub{Conduite à tenir devant une plaie simple}

Soigner correctement une petite plaie (écorchure, coupure superficielle) est un savoir-faire indispensable. La démarche se fait en \textbf{trois étapes}, dans un ordre qui ne change jamais :

\begin{methode}[Soigner une plaie simple]
\begin{enumerate}
\item \textbf{Nettoyer} : laver la plaie à l'eau propre et au savon (de préférence antiseptique), en éliminant les saletés visibles (terre, sable, petits graviers). Rincer abondamment. On se sera d'abord lavé les mains !
\item \textbf{Appliquer un antiseptique} : verser ou tamponner doucement l'antiseptique (bétadine, alcool à 70° sur le pourtour, Dakin) sur toute la surface de la plaie.
\item \textbf{Couvrir} : protéger la plaie avec un \textbf{pansement propre} (compresse stérile maintenue par un sparadrap ou un bandage), renouvelé chaque jour.
\end{enumerate}
\end{methode}

\begin{popfigure}
\begin{center}
\begin{tikzpicture}[font=\small]
% Étape 1
\draw[fill=popblueL, rounded corners=4pt] (0,0) rectangle (3.6,2.2);
\node[align=center] at (1.8,1.1) {\textbf{1. Nettoyer}\\ eau propre + savon\\ (enlever les saletés)};
% Étape 2
\draw[fill=poporangeL, rounded corners=4pt] (4.6,0) rectangle (8.2,2.2);
\node[align=center] at (6.4,1.1) {\textbf{2. Antiseptique}\\ bétadine, alcool 70°,\\ Dakin};
% Étape 3
\draw[fill=popgreenL, rounded corners=4pt] (9.2,0) rectangle (12.8,2.2);
\node[align=center] at (11,1.1) {\textbf{3. Couvrir}\\ pansement propre\\ (compresse stérile)};
% Flèches
\draw[-{Stealth[length=3mm]}, very thick, popdark] (3.6,1.1) -- (4.6,1.1);
\draw[-{Stealth[length=3mm]}, very thick, popdark] (8.2,1.1) -- (9.2,1.1);
\end{tikzpicture}
\end{center}
\legende{Les trois étapes du soin d'une plaie simple : le nettoyage à l'eau et au savon précède toujours l'application de l'antiseptique, puis la plaie est protégée par un pansement propre.}
\end{popfigure}

\begin{attention}[Quand faut-il consulter ?]
L'antisepsie a une \textbf{limite} : elle ne supprime pas tous les microbes, surtout ceux situés \emph{en profondeur} dans les tissus. Il faut consulter rapidement un agent de santé si :
\begin{itemize}
\item la plaie est \textbf{profonde} (coupure large, morsure, clou rouillé enfoncé dans le pied) ;
\item la plaie est \textbf{souillée} par de la terre, de la rouille ou des matières fécales ;
\item la plaie devient rouge, gonflée, douloureuse ou coule du pus (signes d'infection) ;
\item la personne a de la fièvre après la blessure.
\end{itemize}
Dans ces cas, l'agent de santé vérifiera notamment la \textbf{vaccination antitétanique} : le tétanos est une maladie grave transmise par un microbe présent dans la terre et la poussière, qui pénètre par les plaies profondes. Le vaccin antitétanique, rappelé régulièrement, en protège efficacement.
\end{attention}

\courssub{Antisepsie et asepsie : deux pratiques complémentaires}

Récapitulons la différence entre ces deux notions, qui se complètent sans jamais se remplacer :

\begin{popfigure}
\begin{center}
\begin{tikzpicture}[font=\small]
\draw[fill=popblueL, rounded corners=4pt] (0,0) rectangle (6,1.4);
\node[align=center, text width=5.6cm] at (3,0.7) {\textbf{ASEPSIE}\\ \emph{empêche l'arrivée} des microbes\\ (action avant le contact)};
\draw[fill=poporangeL, rounded corners=4pt] (7,0) rectangle (13,1.4);
\node[align=center, text width=5.6cm] at (10,0.7) {\textbf{ANTISEPSIE}\\ \emph{détruit} les microbes déjà présents\\ (action après le contact)};
% lignes d'exemples
\node[align=left, text width=5.6cm, anchor=north] at (3,-0.2) {\textbf{But :} aucun microbe n'atteint la plaie ou le patient\\[2pt] \textbf{Moment :} avant et pendant le soin\\[2pt] \textbf{Exemples :} lavage des mains, gants, matériel stérilisé, blouse propre};
\node[align=left, text width=5.6cm, anchor=north] at (10,-0.2) {\textbf{But :} tuer les microbes sur la peau ou la plaie\\[2pt] \textbf{Moment :} au moment du soin, après nettoyage\\[2pt] \textbf{Exemples :} alcool 70°, bétadine, eau oxygénée, Dakin, savon antiseptique};
% accolade complémentarité
\draw[-{Stealth[length=3mm]}, very thick, popgreen] (3,-4.4) -- (10,-4.4);
\node[align=center, text=popgreen] at (6.5,-4.9) {\textbf{Les deux pratiques sont complémentaires :}\\ on empêche d'abord l'arrivée des microbes, puis on détruit ceux déjà présents.};
\end{tikzpicture}
\end{center}
\legende{Comparaison asepsie / antisepsie : but, moment d'action et exemples de pratiques. L'une empêche l'arrivée des microbes, l'autre détruit ceux déjà présents.}
\end{popfigure}

\begin{aretenir}[L'essentiel]
\begin{itemize}
\item L'\cle{antisepsie} détruit les microbes déjà présents sur la peau, les plaies ou les objets ; l'\cle{asepsie} empêche leur arrivée.
\item Un \cle{antiseptique} (alcool à 70°, bétadine, eau oxygénée, Dakin, savon antiseptique) s'applique sur les tissus vivants, toujours \emph{après} le nettoyage à l'eau et au savon.
\item Devant une plaie simple : \textbf{nettoyer} $\rightarrow$ \textbf{antiseptique} $\rightarrow$ \textbf{pansement propre}.
\item Plaie profonde ou souillée : consulter un agent de santé (risque de tétanos).
\end{itemize}
\end{aretenir}

\courssub{Complément : antisepsie et désinfection}

\begin{savaistu}[Antisepsie ou désinfection ?]
On entend souvent le mot « désinfection ». La \cle{désinfection} concerne les \textbf{objets et les surfaces} (sol, tables, toilettes, instruments non stérilisables) : on utilise par exemple l'\emph{eau de Javel} diluée. L'\cle{antisepsie}, elle, concerne les \textbf{tissus vivants} (peau, plaies) : on n'appliquerait jamais d'eau de Javel concentrée sur une plaie ! Petite histoire : la solution de Dakin, encore utilisée aujourd'hui dans les hôpitaux camerounais, n'est autre que de l'eau de Javel très diluée. Elle fut mise au point pendant la Première Guerre mondiale (1914--1918) par le chimiste Henry Dakin pour soigner les plaies infectées des soldats blessés au front. Son efficacité a sauvé d'innombrables vies, et elle reste un antiseptique bon marché, précieux dans les pays où les ressources sont limitées.
\end{savaistu}

\courssub{Exercice résolu}

\begin{exoresolu}[La bonne conduite à tenir]
\textbf{Énoncé.} Pendant un match de football au quartier, Junior glisse et s'écorche le genou sur un terrain poussiéreux. Son ami lui propose de verser directement de la bétadine sur la plaie pleine de terre, puis de la couvrir d'un vieux tissu. Cette conduite est-elle correcte ? Justifie ta réponse et propose la bonne démarche.
\tcblower
\textbf{Solution détaillée.}

Cette conduite est \textbf{incorrecte}, pour deux raisons :
\begin{enumerate}
\item La bétadine est versée sur une plaie \emph{sale} : l'antiseptique ne traverse pas la couche de terre et de poussière, donc les microbes restent vivants sous les saletés et peuvent provoquer une infection.
\item Le vieux tissu n'est pas propre : il peut déposer de nouveaux microbes sur la plaie (c'est le contraire de l'asepsie).
\end{enumerate}

La bonne démarche, en trois étapes :
\begin{enumerate}
\item \textbf{Nettoyer} la plaie à l'eau propre et au savon pour éliminer la terre et la poussière, après s'être lavé les mains ;
\item \textbf{Appliquer un antiseptique} (bétadine ou alcool à 70° sur le pourtour) sur toute la surface de la plaie ;
\item \textbf{Couvrir} avec un pansement propre (compresse stérile), à renouveler chaque jour.
\end{enumerate}

Enfin, comme la plaie est souillée par de la terre, il est prudent de vérifier auprès d'un agent de santé que la \emph{vaccination antitétanique} de Junior est à jour.
\end{exoresolu}

Ainsi, l'antisepsie complète l'asepsie : quand les microbes sont déjà présents, on sait désormais comment les détruire efficacement. Le TP n°6 te permettra de mettre en pratique ces gestes, avant d'aborder une autre pratique essentielle de prévention : l'utilisation des préservatifs.

\coursec{TP N°6 : Pratique de l'asepsie et de l'antisepsie}

Dans la section précédente, nous avons vu que l'\cle{antisepsie} consiste à détruire les microbes présents sur les \textbf{tissus vivants} (peau, muqueuses, plaies) grâce à des substances appelées antiseptiques (alcool à $70^\circ$, bétadine\dots). Nous avons aussi distingué l'antisepsie de l'\cle{asepsie}, qui vise à empêcher l'arrivée des microbes sur les tissus vivants ou les milieux stériles (lavage des mains, port de gants, stérilisation du matériel, désinfection des surfaces inertes à l'eau de Javel diluée). Mais ces affirmations reposent-elles sur des faits ? Peut-on \emph{voir} concrètement que le lavage des mains réduit le nombre de microbes, et qu'un antiseptique les détruit ? C'est précisément l'objet de ce travail pratique : passer du discours à la preuve expérimentale. Comme toujours en sciences, nous allons observer, comparer, interpréter, puis conclure.

\courssub{Objectifs du TP}

\begin{definition}[Objectif du TP N°6]
Mettre en œuvre les gestes d'\cle{asepsie} (lavage des mains au savon) et d'\cle{antisepsie} (action d'un antiseptique sur un milieu ensemencé), puis \textbf{justifier leur efficacité} par l'observation de cultures microbiennes comparées à des témoins.
\end{definition}

Ce TP poursuit deux buts complémentaires :
\begin{itemize}
\item un \textbf{savoir-faire pratique} : réaliser correctement un lavage des mains, manipuler des boîtes de Pétri aseptiquement, déposer un disque imbibé d'antiseptique ;
\item un \textbf{savoir-faire méthodologique} : construire une expérience comparative (avec témoin), recueillir des résultats dans un tableau, et rédiger un compte rendu scientifique complet.
\end{itemize}

\begin{attention}[Le principe de toute expérience comparative]
On ne peut jamais conclure à l'efficacité d'un traitement (lavage, antiseptique) sans le comparer à un \cle{témoin}, c'est-à-dire une situation identique \emph{sans} le traitement. Le témoin montre ce qui se passe « naturellement » ; la différence entre le témoin et l'expérience mesure l'effet du traitement.
\end{attention}

\courssub{Matériel nécessaire}

\begin{center}
\begin{tabularx}{\linewidth}{|l|X|}\hline
\rowcolor{popblueL} \textbf{Matériel} & \textbf{Rôle dans le TP} \\ \hline
Eau et savon & Lavage des mains (geste d'asepsie) \\ \hline
Alcool à $70^\circ$ & Antiseptique de la peau saine \\ \hline
Bétadine (povidone iodée) & Antiseptique des plaies et muqueuses \\ \hline
Eau de Javel diluée & Désinfectant pour surfaces/inertes (asepsie du matériel) \\ \hline
Compresses, pansements & Nettoyage et protection des plaies \\ \hline
Gants & Protection de l'expérimentateur (asepsie) \\ \hline
Boîtes de Pétri avec milieu de culture (gélose) & Support nutritif sur lequel les microbes se développent en colonies visibles \\ \hline
Disques de papier buvard & Support à imbiber d'antiseptique \\ \hline
Étuve (ou endroit chaud, environ \SI{37}{\celsius}) & Incubation : favorise le développement des microbes \\ \hline
Écouvillons stériles & Prélèvement et ensemencement uniforme \\ \hline
\end{tabularx}
\end{center}

\begin{attention}[Sécurité -- Règle d'or]
Les boîtes de Pétri ensemencées ne s'ouvrent \textbf{jamais} après incubation : les colonies développées peuvent contenir des microbes pathogènes. On observe à travers la boîte fermée, et le matériel est ensuite désinfecté (eau de Javel, autoclave) ou détruit par le professeur.
\end{attention}

\courssub{Manipulation 1 : L'asepsie — effet du lavage des mains}

\textbf{Question posée.} Le lavage des mains au savon réduit-il réellement le nombre de microbes présents sur la peau ?

\textbf{Hypothèse.} Une main lavée au savon porte moins de microbes qu'une main non lavée.

\begin{experience}[Protocole de la manipulation 1]
\begin{enumerate}
\item Étiqueter deux boîtes de Pétri contenant de la gélose : boîte A « main non lavée » (témoin), boîte B « main lavée ».
\item Un élève applique délicatement la pulpe des doigts de sa main \textbf{non lavée} sur la gélose de la boîte A (empreinte), sans appuyer fort ni bouger, puis referme immédiatement la boîte.
\item Le même élève se lave soigneusement les mains à l'eau et au savon pendant \textbf{\SI{30}{\second}} (paumes, dos, espaces interdigitaux, ongles, poignets), se rince, se sèche avec une serviette propre ou papier à usage unique.
\item Il applique ensuite les mêmes doigts sur la gélose de la boîte B, referme la boîte.
\item Placer les deux boîtes fermées à l'étuve à \SI{37}{\celsius} pendant \SI{24}{\hour} à \SI{48}{\hour}.
\item Observer à travers les boîtes fermées et compter (ou estimer) le nombre de colonies dans chaque boîte.
\end{enumerate}
\end{experience}

\textbf{Observations.} Après incubation, la boîte A (main non lavée, témoin) présente de \textbf{nombreuses colonies} de microbes : petites taches rondes, blanches, jaunâtres ou brillantes, de tailles variées, parfois confluentess. La boîte B (main lavée) n'en présente que \textbf{quelques-unes, voire aucune}.

\begin{popfigure}
\begin{center}
\begin{tikzpicture}[scale=0.9]
% Boîte A : main non lavée
\draw[thick,popdark,fill=popcream] (0,0) circle (2.2);
\draw[thick,popdark] (0,0) circle (2.35);
% colonies nombreuses
\foreach \x/\y/\r in {0.3/0.5/0.16, -0.8/0.9/0.12, 1.1/-0.4/0.14, -0.4/-0.9/0.18, 0.9/1.1/0.11, -1.3/-0.2/0.13, 0.1/-1.4/0.12, 1.5/0.5/0.15, -0.2/1.5/0.10, 0.7/-1.0/0.13, -1.0/-1.1/0.11, 1.7/-0.9/0.10, -1.6/0.6/0.12, 0.5/1.7/0.09, -0.6/0.1/0.14, 1.2/0.1/0.10, -0.9/0.4/0.09, 0.0/-0.5/0.12}
  \fill[poporange] (\x,\y) circle (\r);
\node[below,popdark] at (0,-2.6) {\textbf{Boîte A : main non lavée (témoin)}};
\node[below,popdark] at (0,-3.1) {Nombreuses colonies};

% Boîte B : main lavée
\begin{scope}[xshift=6.5cm]
\draw[thick,popdark,fill=popcream] (0,0) circle (2.2);
\draw[thick,popdark] (0,0) circle (2.35);
\fill[poporange] (0.6,0.4) circle (0.13);
\fill[poporange] (-0.8,-0.6) circle (0.11);
\fill[poporange] (0.2,-1.2) circle (0.10);
\node[below,popdark] at (0,-2.6) {\textbf{Boîte B : main lavée au savon \SI{30}{\second}}};
\node[below,popdark] at (0,-3.1) {Rares colonies};
\end{scope}
\end{tikzpicture}
\end{center}
\legende{Comparaison des cultures après empreinte d'une main non lavée (boîte A, témoin) et d'une main lavée au savon pendant \SI{30}{\second} (boîte B). Chaque tache orange représente une colonie microbienne issue d'un microbe déposé sur la gélose. La forte réduction du nombre de colonies dans la boîte B établit l'efficacité de l'asepsie par lavage des mains.}
\end{popfigure}

\textbf{Interprétation.} Chaque colonie visible provient du développement d'un microbe (ou d'un petit groupe de microbes) déposé par les doigts sur la gélose, qui s'est multiplié par divisions successives pendant l'incubation. La comparaison avec le témoin (boîte A) montre que le lavage au savon a \textbf{éliminé ou entraîné mécaniquement la grande majorité des microbes} présents sur la peau : le nombre de colonies est fortement réduit dans la boîte B. Le savon, par son action tensioactive, décolle les microbes piégés dans le film hydrolipidique et les graisses cutanées, et le rinçage les entraîne.

\textbf{Conclusion.} Le lavage des mains à l'eau et au savon pendant \SI{30}{\second} est un geste d'\cle{asepsie} efficace : il réduit drastiquement le nombre de microbes sur les mains et limite ainsi leur transmission aux aliments, aux plaies ou aux autres personnes.

\begin{exemplebox}[Un résultat chiffré typique]
Dans une classe, on a compté après incubation : \num{85} colonies dans la boîte A (main non lavée) et \num{7} colonies dans la boîte B (main lavée \SI{30}{\second} au savon). Le nombre de colonies a donc été divisé par environ $85/7 \approx 12$, soit plus de 10. En pratique, un lavage soigneux des mains au savon pendant \SI{30}{\second} peut diviser le nombre de colonies observées \textbf{par 10 ou plus}. C'est pourquoi le personnel soignant des hôpitaux et centres de santé camerounais se lave les mains avant et après chaque soin (règle des « 5 moments » de l'OMS).
\end{exemplebox}

\courssub{Manipulation 2 : L'antisepsie — action d'un antiseptique}

\textbf{Question posée.} Un antiseptique (alcool à $70^\circ$, bétadine) est-il capable de détruire les microbes présents sur un milieu de culture ?

\textbf{Hypothèse.} Autour d'un disque imbibé d'antiseptique déposé sur un milieu ensemencé, les microbes ne se développent pas.

\begin{experience}[Protocole de la manipulation 2]
\begin{enumerate}
\item Étiqueter deux boîtes de Pétri : boîte C et boîte D.
\item Ensemencer \textbf{uniformément} la gélose des deux boîtes avec la même souche microbienne (par étalage à l'écouvillon stérile sur toute la surface).
\item Déposer au centre de la boîte C un disque de papier buvard imbibé d'\textbf{antiseptique} (alcool à $70^\circ$ ou bétadine) à l'aide d'une pince stérile.
\item Déposer au centre de la boîte D un disque imbibé d'\textbf{eau stérile} : c'est le \textbf{témoin} (contrôle du support disque + eau).
\item Incuber les deux boîtes fermées à \SI{37}{\celsius} pendant \SI{24}{\hour} à \SI{48}{\hour}.
\item Observer : mesurer le diamètre de la zone sans colonies autour du disque (auréole d'inhibition).
\end{enumerate}
\end{experience}

\textbf{Observations.} Dans la boîte D (témoin, disque à l'eau stérile), les microbes se développent partout sur la gélose, y compris au contact du disque. Dans la boîte C (antiseptique), on observe autour du disque une \textbf{zone circulaire nette, sans colonies}, appelée \cle{auréole d'inhibition} ; plus loin du disque, hors de cette zone, les microbes se développent normalement.

\begin{popfigure}
\begin{center}
\begin{tikzpicture}[scale=0.9]
% Boîte C : antiseptique
\draw[thick,popdark,fill=popcream] (0,0) circle (2.2);
\draw[thick,popdark] (0,0) circle (2.35);
% fond ensemencé (nombreux petits points) sauf dans l'auréole (rayon 1.2)
\foreach \a in {0,15,...,345}{
  \pgfmathsetmacro{\rr}{1.35 + 0.75*rnd}
  \fill[popgreen!70] ({\rr*cos(\a)},{\rr*sin(\a)}) circle (0.05);
}
\foreach \a in {7,37,...,337}{
  \pgfmathsetmacro{\rr}{1.5 + 0.55*rnd}
  \fill[popgreen!70] ({\rr*cos(\a)},{\rr*sin(\a)}) circle (0.04);
}
% auréole d'inhibition
\draw[thick,popblue,dashed] (0,0) circle (1.2);
% disque central
\draw[thick,popdark,fill=popblueL] (0,0) circle (0.35);
\node[popdark] at (0,0) {\footnotesize disque};
% flèche auréole
\draw[->,thick,popblue] (0.35,0) -- (1.2,0) node[midway,above,popblue]{\footnotesize auréole};
\node[below,popdark] at (0,-2.6) {\textbf{Boîte C : disque + antiseptique}};
\node[below,popdark] at (0,-3.1) {Auréole d'inhibition (rayon \SI{1.2}{\centi\meter})};

% Boîte D : témoin eau stérile
\begin{scope}[xshift=6.5cm]
\draw[thick,popdark,fill=popcream] (0,0) circle (2.2);
\draw[thick,popdark] (0,0) circle (2.35);
\foreach \a in {0,12,...,348}{
  \pgfmathsetmacro{\rr}{0.55 + 1.55*rnd}
  \fill[popgreen!70] ({\rr*cos(\a)},{\rr*sin(\a)}) circle (0.045);
}
\draw[thick,popdark,fill=white] (0,0) circle (0.35);
\node[popdark] at (0,0) {\footnotesize disque};
\node[below,popdark] at (0,-2.6) {\textbf{Boîte D : disque + eau stérile (témoin)}};
\node[below,popdark] at (0,-3.1) {Colonies partout, même au contact du disque};
\end{scope}
\end{tikzpicture}
\end{center}
\legende{Action d'un antiseptique sur un milieu ensemencé uniforme. Boîte C : l'antiseptique diffuse par gradient de concentration depuis le disque et détruit les microbes alentour, créant une auréole d'inhibition (zone sans colonies, cercle en pointillés bleus). Boîte D (témoin) : l'eau stérile n'a aucun effet antimicrobien, les colonies (points verts) se développent jusqu'au contact du disque.}
\end{popfigure}

\textbf{Interprétation.} L'antiseptique \textbf{diffuse} dans la gélose autour du disque selon un gradient de concentration. Dans la zone où sa concentration est supérieure à la concentration minimale inhibitrice (CMI), il \textbf{détruit les microbes} (action bactéricide) ou \textbf{empêche leur multiplication} (action bactériostatique) : c'est l'auréole d'inhibition. La comparaison avec le témoin (boîte D) prouve que c'est bien la molécule antiseptique — et non le disque de papier ni l'eau — qui est responsable de l'absence de colonies.

\textbf{Conclusion.} L'antiseptique détruit ou inhibe les microbes au contact : l'\cle{antisepsie} est efficace pour désinfecter la peau saine avant un acte invasif, nettoyer une plaie ou désinfecter du matériel thermosensible. La taille de l'auréole dépend de la nature de l'antiseptique, de sa concentration, de la sensibilité du microbe et de sa pouvoir de diffusion dans l'agar.

\begin{savaistu}[De l'auréole d'inhibition à l'antibiogramme]
Le principe de l'auréole d'inhibition est exactement celui utilisé dans les laboratoires d'analyses médicales (hôpitaux de référence au Cameroun, Centre Pasteur du Cameroun) pour réaliser un \cle{antibiogramme} : on dépose plusieurs disques imbibés d'antibiotiques différents (normalisés) sur une culture standardisée du microbe responsable d'une infection, et on mesure les diamètres des auréoles après incubation. L'antibiotique qui produit la plus grande auréole (sensible \textbf{S}) est le plus efficace contre ce microbe : c'est celui que le médecin prescrira. Cette méthode permet de lutter contre l'antibiorésistance en choisissant le traitement ciblé.
\end{savaistu}

\courssub{Tableau de résultats et compte rendu}

Les résultats de l'ensemble du TP peuvent être rassemblés dans un tableau de synthèse :

\begin{center}
\begin{tabularx}{\linewidth}{|l|X|X|}\hline
\rowcolor{popblueL} \textbf{Boîte} & \textbf{Traitement} & \textbf{Résultat après incubation} \\ \hline
A & Empreinte main non lavée (témoin) & Nombreuses colonies (ex. \num{85}) \\ \hline
B & Empreinte main lavée au savon \SI{30}{\second} & Rares colonies (ex. \num{7}) \\ \hline
C & Milieu ensemencé + disque antiseptique & Auréole d'inhibition nette (ex. \SI{2.4}{\centi\meter} de diamètre) \\ \hline
D & Milieu ensemencé + disque eau stérile (témoin) & Colonies partout, aucune auréole \\ \hline
\end{tabularx}
\end{center}

\begin{methode}[Rédiger un compte rendu de TP -- Standard BEPC]
Un compte rendu de TP suit toujours le même plan rigoureux, dans l'ordre :
\begin{enumerate}
\item \textbf{Objectif} : la question précise à laquelle le TP doit répondre (une phrase).
\item \textbf{Matériel} : la liste du matériel et des produits utilisés.
\item \textbf{Protocole} : les étapes de la manipulation, numérotées, rédigées au passé composé ou à l'infinitif, en \textbf{précisant explicitement quel est le témoin} et comment il a été réalisé.
\item \textbf{Résultats} : observations brutes (ce qu'on voit), présentées dans un tableau et/ou des schémas légendés, \textbf{sans aucune interprétation} à ce stade.
\item \textbf{Interprétation} : on explique les résultats en les \textbf{comparant systématiquement au témoin} et en les reliant aux notions du cours (mécanismes).
\item \textbf{Conclusion} : on répond à la question de départ (objectif), de façon justifiée par les observations et l'interprétation.
\end{enumerate}
\end{methode}

\begin{attention}[Erreur fréquente au BEPC : conclure sans comparer au témoin]
Écrire « la boîte B contient peu de colonies, donc le savon est efficace » est une conclusion \textbf{incomplète et non scientifique} : peut-être que la main était déjà propre avant le lavage, ou que l'élève a peu touché la gélose ! Toute conclusion expérimentale repose sur la \textbf{comparaison explicite avec le témoin} : c'est la différence entre la boîte A (témoin) et la boîte B qui prouve l'effet du lavage. De même pour l'antiseptique : c'est la comparaison C vs D qui prouve l'effet de la molécule. \textbf{Au BEPC, une conclusion sans référence au témoin perd des points.}
\end{attention}

\begin{exoresolu}[Rédiger l'interprétation et la conclusion du TP]
\textbf{Énoncé.} À l'issue du TP N°6, un élève a obtenu les résultats suivants : boîte A (main non lavée) : \num{92} colonies ; boîte B (main lavée \SI{30}{\second} au savon) : \num{8} colonies ; boîte C (disque + bétadine) : auréole d'inhibition de \SI{2.4}{\centi\meter} de diamètre ; boîte D (disque + eau stérile) : aucune auréole. Rédiger l'interprétation et la conclusion de ce TP.
\tcblower
\textbf{Solution détaillée.}

\emph{Interprétation.} La comparaison des boîtes A et B montre que le lavage des mains au savon a réduit le nombre de colonies de \num{92} à \num{8}, soit une division par plus de 11 ($92/8 = 11,5$) : le savon a donc éliminé la grande majorité des microbes présents sur la peau par action mécanique (décollage et entraînement). La comparaison des boîtes C et D montre que seule la bétadine (et non le disque, puisque le témoin D ne présente aucune auréole) a empêché le développement des microbes sur un rayon de \SI{1.2}{\centi\meter} : la bétadine a donc détruit ou inhibé les microbes autour du disque par diffusion dans l'agar.

\emph{Conclusion.} Le lavage des mains à l'eau et au savon pendant \SI{30}{\second} est un geste d'asepsie efficace : il réduit fortement le nombre de microbes sur les mains (facteur $>10$) et limite ainsi leur transmission. La bétadine, comme les autres antiseptiques, détruit les microbes au contact par diffusion : l'antisepsie est efficace pour désinfecter la peau et les plaies. Asepsie et antisepsie sont donc deux pratiques complémentaires et indispensables pour éviter la contamination et l'infection.
\end{exoresolu}

\begin{exercice}[À toi de jouer -- Entraînement BEPC]
Lors d'un TP, un groupe obtient les résultats suivants :
\begin{itemize}
\item Boîte témoin (main non lavée) : \num{60} colonies ;
\item Boîte « main rincée à l'eau seule (sans savon) » : \num{45} colonies ;
\item Boîte « main lavée au savon \SI{30}{\second} » : \num{5} colonies.
\end{itemize}
\begin{enumerate}
\item Calculer le facteur de réduction du nombre de colonies entre le rinçage à l'eau seule et le lavage au savon.
\item Que montre la comparaison entre l'eau seule et le savon sur l'efficacité respective de ces deux gestes ?
\item Pourquoi la boîte témoin (main non lavée) est-elle indispensable pour valider l'expérience ?
\end{enumerate}
\end{exercice}

\begin{corrige}[Exercice 1]
\begin{enumerate}
\item Le lavage au savon divise le nombre de colonies par $45/5 = 9$ par rapport au rinçage à l'eau seule (et par $60/5 = 12$ par rapport à la main non lavée).
\item L'eau seule n'élimine qu'une partie des microbes (passage de \num{60} à \num{45} colonies, soit une réduction de 25\,\%) ; le savon est beaucoup plus efficace (\num{5} colonies, réduction de plus de 90\,\%) car son action tensioactive décolle les microbes piégés dans les graisses et saletés cutanées, que le rinçage entraîne ensuite.
\item Le témoin montre le nombre de colonies de base sans aucun traitement : c'est la référence absolue qui permet de quantifier l'effet de \emph{chaque} traitement testé (eau seule, savon). Sans lui, impossible d'affirmer que l'eau ou le savon ont eu un effet réel, ni de comparer leurs efficacités relatives.
\end{enumerate}
\end{corrige}

\begin{aretenir}[L'essentiel du TP N°6]
\begin{itemize}
\item L'\cle{asepsie} empêche l'arrivée des microbes : le lavage des mains au savon pendant \SI{30}{\second} peut diviser le nombre de colonies par 10 ou plus (action mécanique tensioactive).
\item L'\cle{antisepsie} détruit ou inhibe les microbes sur les tissus vivants : un disque imbibé d'antiseptique crée une \cle{auréole d'inhibition} sur un milieu ensemencé (preuve par diffusion).
\item Toute conclusion expérimentale rigoureuse repose sur la \textbf{comparaison explicite avec un témoin}.
\item Un compte rendu de TP suit le plan standard : Objectif, Matériel, Protocole (avec témoin), Résultats (tableau/schémas), Interprétation (comparaison + mécanismes), Conclusion (réponse justifiée).
\end{itemize}
\end{aretenir}

Ce TP nous a donc permis de prouver expérimentalement l'efficacité de l'asepsie (lavage des mains) et de l'antisepsie (antiseptiques) sur les microbes de la peau et des surfaces. Mais certaines contaminations ne passent ni par les mains ni par les plaies : elles se transmettent lors des rapports sexuels par contact direct avec des liquides biologiques infectés. Une pratique de protection spécifique est alors nécessaire : c'est l'utilisation des préservatifs, que nous étudierons dans la section suivante.

\coursec{L'utilisation des préservatifs}

Après avoir découvert au laboratoire comment l'asepsie et l'antisepsie permettent de maîtriser la contamination microbienne sur les objets et la peau (TP N°6), nous allons maintenant étudier une application vitale de ces principes dans la vie quotidienne : la protection lors des rapports sexuels. Les rapports sexuels non protégés constituent en effet une voie majeure de contamination par les agents des Infections Sexuellement Transmissibles (IST) — VIH/sida, blennorragie, syphilis, hépatite B — et sont la cause principale des grossesses précoces non désirées chez les adolescentes. Comment une simple gaine peut-elle constituer une barrière efficace ? C'est ce que nous allons analyser avec rigueur.

\courssub{Le préservatif masculin : description et mode d'emploi correct}

Le préservatif masculin (aussi appelé « condom ») est le moyen de protection le plus répandu. C'est une gaine fine, généralement en latex (caoutchouc naturel) ou en polyuréthane pour les personnes allergiques au latex, qui se déroule sur le pénis en érection avant tout contact sexuel.

\begin{definition}[Préservatif masculin]
Gaine fine et élastique, à usage unique, qui recouvre le pénis en érection lors d'un rapport sexuel. Elle constitue une barrière physique étanche empêchant le contact direct entre les muqueuses et les sécrétions génitales (sperme, sécrétions vaginales, sang), bloquant ainsi le passage des spermatozoïdes et des agents pathogènes (bactéries, virus, champignons).
\end{definition}

Pour que cette barrière soit efficace, son utilisation doit respecter un protocole strict, analogue aux règles d'asepsie que nous avons manipulées au TP N°6 : toute erreur de manipulation peut créer une brèche (rupture, glissement) par où passent les microbes.

\begin{methode}[Mise en place correcte du préservatif masculin (procédure standard OMS)]
\textbf{Étape 1 : Vérification avant ouverture.} Contrôler la date de péremption imprimée sur l'emballage (un préservatif périmé est fragilisé). Vérifier la présence du marquage de conformité (ex. norme ISO 4074, marque CE). S'assurer que l'emballage est intact (pas de trou, pas d'écrasement).
\textbf{Étape 2 : Ouverture soignée.} Ouvrir l'emballage par le bord prédécoupé, avec les doigts. \textbf{Interdiction formelle :} ne pas utiliser les dents, des ciseaux, un couteau ou tout objet tranchant qui risquerait de déchirer le préservatif.
\textbf{Étape 3 : Sens de déroulement.} Avant de le mettre, déterminer le sens de déroulement : le rebord (bourrelet) doit être à l'extérieur, le préservatif doit se dérouler facilement vers le bas. Ne pas le dérouler avant de le mettre.
\textbf{Étape 4 : Mise en place \emph{avant tout contact génital}.} Pincer le réservoir (bout du préservatif) entre le pouce et l'index pour chasser l'air (évite l'éclatement par surpression lors de l'éjaculation). Le poser sur le gland du pénis en érection. De l'autre main, le dérouler délicatement jusqu'à la base du pénis.
\textbf{Étape 5 : Pendant le rapport.} S'assurer qu'il reste en place. En cas de rupture ou de glissement, s'arrêter immédiatement, se retirer, jeter le préservatif et en mettre un neuf.
\textbf{Étape 6 : Retrait après l'éjaculation.} Se retirer \textbf{avant que le pénis ne redevienne mou}, en tenant le préservatif à la base pour éviter tout débordement de sperme.
\textbf{Étape 7 : Élimination.} Faire un nœud au préservatif usagé, le jeter dans une poubelle (pas dans les toilettes : risque de bouchage et de pollution). \textbf{Jamais le réutiliser.}
\end{methode}

\begin{attention}[Erreurs fréquentes et idées reçues dangereuses]
\begin{itemize}
    \item \textbf{Réutilisation :} Un préservatif est à \textbf{usage unique}. Le laver, le retourner ou le réutiliser détruit l'intégrité du latex et n'élimine pas les microbes.
    \item \textbf{Double préservatif :} Mettre deux préservatifs l'un sur l'autre (ou un masculin + un féminin) \textbf{augmente} le risque de rupture par frottement entre les deux couches de latex. Un seul, bien mis, suffit.
    \item \textbf{Lubrifiants incompatibles :} N'utiliser \textbf{que des lubrifiants à base d'eau ou de silicone} (vendus en pharmacie/supermarché). Les corps gras (huile, vaseline, beurre de karité, crème hydratante, salive en grande quantité) dégradent le latex en quelques minutes et provoquent la rupture.
    \item \textbf{Conservation inadaptée :} Ne pas garder un préservatif dans une poche de pantalon (chaleur + frottement), dans un portefeuille (écrasement), au soleil ou dans la boîte à gants d'une voiture (chaleur excessive). Le conserver à température ambiante, à l'abri de la lumière et de l'humidité.
    \item \textbf{Oubli du réservoir :} Ne pas pincer le bout pour chasser l'air crée une poche d'air comprimé qui éclate au moment de l'éjaculation.
\end{itemize}
\end{attention}

\courssub{Le préservatif féminin : une alternative sous contrôle de la femme}

Moins connu mais tout aussi efficace, le préservatif féminin (ex. FC2) offre une protection dont la mise en place est gérée par la femme.

\begin{definition}[Préservatif féminin]
Gaine souple en nitrile (polymère synthétique non allergisant), plus large et plus longue que le préservatif masculin, munie de deux anneaux flexibles : un anneau interne (au fond de la gaine) qui sert à l'insérer dans le vagin, et un anneau externe (à l'ouverture) qui reste à l'extérieur et recouvre la vulve. Il peut être mis en place plusieurs heures avant le rapport.
\end{definition}

\begin{popfigure}
\begin{tikzpicture}[scale=0.9, every node/.style={font=\small}]
    % Definition of virus pic
    \tikzset{virus/.pic={\draw[poppink, line width=0.5pt] (0,0) circle (2pt); \foreach \a in {0,45,...,315} \draw[poppink, line width=0.5pt] (0,0) -- (\a:4pt);}}
    
    % Anatomie simplifiée (coupe sagittale)
    % Vagin / Col utérus
    \draw[popdark, line width=1.5pt] (-2, -1.5) .. controls (-2, 0) and (2, 0) .. (2, -1.5);
    \draw[popdark, line width=1.5pt] (-2, 1.5) .. controls (-2, 0) and (2, 0) .. (2, 1.5);
    % Col utérus
    \draw[popdark, fill=popcream] (-0.5, -1.5) rectangle (0.5, -2.2);
    \node[popdark] at (0, -2.5) {Col de l'utérus};
    
    % Préservatif féminin (nitrile)
    \draw[popblue, line width=2pt, fill=popblue!10] (-2.2, -1.3) .. controls (-2.2, 1.2) and (2.5, 1.2) .. (2.5, -1.3);
    \draw[popblue, line width=2pt, fill=popblue!10] (-2.2, 1.3) .. controls (-2.2, 1.8) and (2.5, 1.8) .. (2.5, 1.3);
    
    % Anneau interne (au fond)
    \draw[poporange, line width=3pt] (-1.8, -1.0) ellipse (0.8 and 0.15);
    \node[poporange, anchor=north] at (-1.8, -1.2) {Anneau interne};
    
    % Anneau externe (à l'entrée)
    \draw[poporange, line width=3pt] (2.2, 1.3) ellipse (0.8 and 0.15);
    \node[poporange, anchor=south] at (2.2, 1.6) {Anneau externe (recouvre vulve)};
    
    % Pénis (schématique) traversant le préservatif
    \draw[popmuted, line width=4pt] (0, 2.5) -- (0, -1.0);
    \node[popmuted, anchor=south] at (0, 2.6) {Pénis};
    
    % Microbes bloqués (représentation symbolique)
    \foreach \x/\y in {0.5/0.5, -0.8/0.2, 1.0/-0.5, -1.2/-0.8} {
        \pic at (\x, \y) {virus};
        \draw[poppink, ->, line width=1pt] (\x, \y) -- ++(0.3, 0.3);
        \draw[popblue, line width=2pt, -|] (\x+0.3, \y+0.3) -- ++(0.1, 0.1); % Barrière
    }
    \node[poppink, anchor=west] at (3, 1.0) {Agents pathogènes};
    \node[popblue, anchor=west] at (3, 0.5) {Bloqués par la paroi};
    
    % Flèches "Pas de contact"
    \draw[popgreen, <->, thick] (-2.5, 0) -- (-3.5, 0) node[left, popgreen] {Pas de contact muqueuses/muqueuses};
    \draw[popgreen, <->, thick] (-2.5, -1.0) -- (-3.5, -1.0) node[left, popgreen] {Pas de contact sperme/vagin};
\end{tikzpicture}
\legende{Principe de barrière du préservatif féminin (coupe sagittale schématique). La gaine en nitrile (bleu) tapisse le vagin et recouvre la vulve (anneau externe orange). L'anneau interne (orange) maintient le fond contre le col. Les microbes (rose) sont bloqués par la paroi intacte. Il n'y a aucun contact direct entre les muqueuses des partenaires ni entre le sperme et le vagin.}
\end{popfigure}

\begin{savaistu}[Le préservatif féminin au Cameroun]
Au Cameroun, le préservatif féminin (souvent distribué sous la marque FC2 ou « Female Condom ») est disponible dans les centres de santé, les pharmacies communautaires et via les ONG de lutte contre le sida (comme le CNLS, CAMNAFAW, ou les programmes du Fonds Mondial). Il est particulièrement apprécié car il peut être mis en place \emph{avant} le rapport (jusqu'à 8h avant), sans interrompre les préliminaires, et il ne nécessite pas d'érection pour être maintenu en place. Il résiste aux lubrifiants à base d'huile (contrairement au latex). Son coût est plus élevé, mais des programmes de subvention le rendent accessible (souvent 100-200 FCFA l'unité en programme public).
\end{savaistu}

\courssub{Mécanisme de protection : une barrière physique assimilée à l'asepsie}

Le principe fondamental est identique à l'asepsie chirurgicale vue au TP N°6 : \textbf{empêcher la rencontre entre le contaminant (microbes, spermatozoïdes) et le milieu récepteur (muqueuses, sang, ovule)}.

\begin{propriete}[Mécanisme barrière du préservatif (Savoir admis)]
Le préservatif (masculin ou féminin) agit comme une \textbf{barrière physique continue et étanche}.
\begin{itemize}
    \item \textbf{Taille des pores :} Le latex et le nitrile sont des polymères dont la structure moléculaire ne présente pas de pores traversants. Les virus les plus petits (VIH : \SI{100}{\nano\meter}, VHB : \SI{42}{\nano\meter}) sont des milliers de fois plus gros que les espaces intermoléculaires du matériau. Ils ne peuvent pas traverser la paroi intacte.
    \item \textbf{Imperméabilité aux fluides :} Le sperme, les sécrétions vaginales, le sang (vecteurs des IST) sont bloqués.
    \item \textbf{Séparation des muqueuses :} Il évite le contact direct muqueuse-muqueuse (pénis/vagin, pénis/anus, bouche/génital), supprimant la porte d'entrée des germes (lésions microscopiques, inflammation).
\end{itemize}
C'est une \textbf{asepsie mécanique} appliquée à la sexualité : on isole la source de contamination.
\end{propriete}

\begin{popfigure}
\begin{tikzpicture}[scale=1.1, every node/.style={font=\small}]
    % Definition of virus pic
    \tikzset{virus/.pic={\draw[poppink, line width=0.5pt] (0,0) circle (2pt); \foreach \a in {0,45,...,315} \draw[poppink, line width=0.5pt] (0,0) -- (\a:4pt);}}
    
    % Schéma comparatif : Sans préservatif vs Avec préservatif
    % --- SANS ---
    \begin{scope}[xshift=-5cm]
        \node[popdark, align=center, font=\bfseries] at (0, 2.5) {SANS PRÉSERVATIF};
        % Partenaire 1 (Pénis)
        \draw[popmuted, fill=popmuted!20, line width=1pt] (-1.5, 0) rectangle (-0.5, 1.5);
        \node[popdark] at (-1, 1.7) {Pénis};
        % Microbes sur le pénis
        \foreach \y in {0.2, 0.6, 1.0, 1.4} {
            \pic at (-1, \y) {virus};
        }
        % Flèches de passage
        \draw[poppink, ->, thick] (-0.5, 0.7) -- (0.5, 0.7) node[midway, above] {Contact direct};
        \draw[poppink, ->, thick] (-0.5, 1.0) -- (0.5, 1.0);
        % Partenaire 2 (Vagin)
        \draw[popmuted, fill=popmuted!20, line width=1pt] (0.5, 0) rectangle (1.5, 1.5);
        \node[popdark] at (1, 1.7) {Vagin};
        % Microbes entrés
        \foreach \y in {0.3, 0.7, 1.1} {
            \pic at (1, \y) {virus};
        }
        \node[poppink, align=center, font=\bfseries] at (0, -0.8) {Contamination :\\IST + Grossesse possible};
    \end{scope}
    
    % --- AVEC ---
    \begin{scope}[xshift=5cm]
        \node[popdark, align=center, font=\bfseries] at (0, 2.5) {AVEC PRÉSERVATIF (Bien mis)};
        % Partenaire 1
        \draw[popmuted, fill=popmuted!20, line width=1pt] (-1.5, 0) rectangle (-0.5, 1.5);
        \node[popdark] at (-1, 1.7) {Pénis};
        % Microbes sur le pénis (côté externe)
        \foreach \y in {0.2, 0.6, 1.0, 1.4} {
            \pic at (-1, \y) {virus};
        }
        % Préservatif (Latex)
        \draw[popblue, fill=popblue!15, line width=2pt] (-0.5, -0.2) rectangle (0.5, 1.7);
        \node[popblue, rotate=90, anchor=south] at (-0.7, 0.75) {Latex / Nitrile intact};
        % Flèches bloquées
        \draw[poppink, ->, thick] (-0.5, 0.7) -- (-0.1, 0.7);
        \draw[popblue, line width=3pt] (-0.1, 0.7) -- (0.1, 0.7); % Barrière visuelle
        \node[popblue, font=\bfseries] at (0, 0.9) {BLOQUÉ};
        % Partenaire 2 (Protégé)
        \draw[popmuted, fill=popgreen!10, line width=1pt] (0.5, 0) rectangle (1.5, 1.5);
        \node[popdark] at (1, 1.7) {Vagin (sain)};
        \node[popgreen, align=center, font=\bfseries] at (0, -0.8) {Protection :\\Barrière étanche};
    \end{scope}
\end{tikzpicture}
\legende{Comparaison schématique du mécanisme barrière. À gauche : contact direct entre muqueuses et fluides contaminés $\rightarrow$ passage des agents pathogènes (VIH, \textit{Neisseria gonorrhoeae}, \textit{Treponema pallidum}, VHB) et des spermatozoïdes. À droite : le préservatif (bleu) interpose une paroi polymère continue sans pore. Les microbes (rose) restent confinés du côté externe ; la muqueuse réceptrice (vert) reste stérile.}
\end{popfigure}

\courssub{Double efficacité : prévention des IST et contraception — Limites réelles}

Le préservatif est la \textbf{seule méthode} qui cumule deux protections simultanées.

\begin{propriete}[Double efficacité du préservatif (Savoir admis)]
\begin{enumerate}
    \item \textbf{Prévention des IST (dont VIH/sida) :} Utilisé correctement et systématiquement (à chaque rapport, du début à la fin), le préservatif masculin en latex réduit le risque de transmission du VIH d'environ \textbf{80\% à 95\%} (méta-analyses OMS/ONUSIDA). Pour les IST bactériennes (blennorragie, syphilis, chlamydiae) transmises par fluides, l'efficacité est proche de \textbf{90-98\%}. Pour les IST transmises par contact cutané-muqueux direct (herpès, HPV/condylomes), la protection est partielle (\textasciitilde 70\%) car les lésions peuvent siéger en dehors de la zone couverte (base du pénis, scrotum, cuisses, vulve).
    \item \textbf{Contraception :} Efficacité théorique (usage parfait) : \textbf{98\%} (indice de Pearl \textasciitilde 2). Efficacité réelle (usage courant, avec erreurs) : \textbf{85-87\%} (indice de Pearl \textasciitilde 13-15). Cela signifie qu'avec une utilisation typique, environ 13 à 15 femmes sur 100 utilisant le préservatif comme unique contraception tomberont enceintes sur un an.
\end{enumerate}
\end{propriete}

\begin{popfigure}
\begin{tikzpicture}
\begin{axis}[
    width=\linewidth, height=7cm,
    ybar, bar width=12pt,
    ymin=0, ymax=100,
    ylabel={Risque relatif de transmission du VIH (\% du risque sans protection)},
    symbolic x coords={Sans préservatif, Avec préservatif (usage parfait), Avec préservatif (usage réel)},
    xtick=data,
    nodes near coords={\pgfmathprintnumber\pgfplotspointmeta\,\%},
    nodes near coords style={font=\small, popdark},
    ymajorgrids=true,
    grid style={popline},
    axis lines*=left,
    enlarge x limits=0.3,
    tick label style={font=\small, align=center},
    title style={font=\bfseries\small, popdark},
    title={Efficacité comparative du préservatif masculin sur la transmission du VIH (Données indicatives OMS/ONUSIDA)},
]
\addplot[poppink, fill=poppink!60] coordinates {(Sans préservatif, 100)};
\addplot[popblue, fill=popblue!60] coordinates {(Avec préservatif (usage parfait), 5)}; % ~95% reduction
\addplot[poporange, fill=poporange!60] coordinates {(Avec préservatif (usage réel), 15)}; % ~85% reduction
\legend{Risque résiduel}
\end{axis}
\end{tikzpicture}
\legende{Graphique comparatif du risque relatif de transmission du VIH. Sans protection : risque de référence 100\%. Avec usage parfait (correct, systématique, sans rupture) : risque résiduel \textasciitilde 5\% (réduction 95\%). Avec usage réel (erreurs de mise en place, rupture, glissement, usage non systématique) : risque résiduel \textasciitilde 15-20\% (réduction 80-85\%). L'écart entre usage parfait et usage réel souligne l'importance de l'éducation à la bonne utilisation.}
\end{popfigure}

\begin{attention}[Limites et causes d'échec — Pourquoi « usage réel » $\neq$ « usage parfait » ?]
\begin{itemize}
    \item \textbf{Rupture :} Due à la chaleur, péremption, lubrifiant gras, bijoux (piercing, bague), ongles longs, absence de réservoir pincé, rapport trop long/vigoureux sans lubrifiant additionnel.
    \item \textbf{Glissement :} Préservatif trop grand, perte d'érection avant retrait, retrait tardif (pénis mou).
    \item \textbf{Mauvaise utilisation :} Mis après le début de la pénétration (liquide pré-éjaculatoire déjà présent, contient spermatozoïdes et VIH), enlevé avant la fin, réutilisé, mis à l'envers puis retourné (contamination face externe).
    \item \textbf{Non utilisation systématique :} « Une fois sans, ce n'est pas grave » $\rightarrow$ une seule exposition suffit pour contracter le VIH ou une IST, ou tomber enceinte.
\end{itemize}
\end{attention}

\courssub{Comportements complémentaires : la stratégie combinée (Abstinence - Fidélité - Préservatif - Dépistage)}

Le préservatif est un outil technique puissant, mais son efficacité populationnelle repose sur des comportements individuels et collectifs. La stratégie globale de lutte contre le VIH/sida et les IST au Cameroun (et internationalement) repose sur l'approche combinée.

\begin{aretenir}[Stratégie de prévention combinée (Adaptée OMS/CNLS Cameroun)]
Pour une protection optimale de sa santé sexuelle et reproductive, on associe :
\begin{enumerate}
    \item \textbf{Abstinence (ou retardement des premiers rapports) :} Risque zéro absolu. Choix personnel, surtout pertinent pour les jeunes adolescents (classe de 3ème : 14-15 ans). C'est la seule méthode 100\% efficace contre IST et grossesse.
    \item \textbf{Fidélité réciproque au sein d'un couple stable \emph{et} dépisté :} Si les deux partenaires sont séro-négatifs (VIH, hépatites, syphilis) et restent exclusifs, le risque IST est nul. Nécessite la confiance, le dialogue et \textbf{le dépistage préalable commun}.
    \item \textbf{Utilisation correcte et systématique du préservatif (masculin ou féminin) :} Pour tout rapport hors couple stable dépisté, ou en plus dans le couple si désir de contraception non hormonale ou si un partenaire a une IST en traitement. C'est le « filet de sécurité ».
    \item \textbf{Dépistage régulier (VIH, Hépatite B, Syphilis) :} Connaître son statut sérologique permet d'accéder au traitement (ARV pour VIH = charge virale indétectable = intransmissible \textbf{I=I}), de protéger l'autre et d'éviter les complications (stérilité, cancer du col utérin lié HPV). Au Cameroun, le dépistage VIH est gratuit dans tous les centres de santé publics et agréés.
    \item \textbf{Information sans tabou et consultation précoce :} Parler sexualité, consentement, protection avec un personnel de santé (infirmier, médecin, sage-femme, conseiller pair-éducateur), un enseignant de SVT, ou une structure jeunesse (ex. CMPJ - Centres Multifonctionnels de Promotion des Jeunes). Ne pas attendre les symptômes : beaucoup d'IST sont asymptomatiques (surtout chez la fille/femme).
    \item \textbf{Vaccination :} Vaccin anti-VHB (intégré au PEV Cameroun depuis 2005, 3 doses), Vaccin anti-HPV (filles 9-14 ans, campagne nationale) pour prévenir le cancer du col de l'utérus.
\end{enumerate}
\end{aretenir}

\begin{exemplebox}[Situation concrète : « Premier rapport protégé »]
\textbf{Contexte :} Paul (16 ans) et Grâce (15 ans), élèves en 3ème, décident d'avoir leur premier rapport sexuel. Ils sont allés ensemble au Centre de Santé du district pour se faire dépister (VIH, VHB, Syphilis) : tous deux négatifs. Ils achètent des préservatifs masculins (marque certifiée ISO) en pharmacie.
\textbf{Analyse :}
\begin{itemize}
    \item \textbf{Points forts :} Dépistage préalable (élimine risque IST antérieur), achat en pharmacie (conservation correcte, non périmés), décision commune, préservatifs certifiés.
    \item \textbf{Points de vigilance pour l'acte :} Vérifier date + emballage. Ouvrir avec les doigts. Mettre \textbf{AVANT} tout contact génital (liquide pré-éjaculatoire). Pincer réservoir. Dérouler jusqu'à la base. Lubrifiant à base d'eau si besoin. Retrait immédiat après éjaculation en tenant la base. Nœud + poubelle.
    \item \textbf{Résultat :} Protection quasi-totale contre VIH/IST (barrière physique) et contraception efficace (blocage spermatozoïdes). Ils ont appliqué l'asepsie sexuelle.
\end{itemize}
\end{exemplebox}

\begin{exoresolu}[Exercice : Analyse d'erreurs de manipulation]
\textbf{Énoncé :} Lors d'une séance d'éducation par les pairs, un élève affirme : « Moi, je mets deux préservatifs l'un sur l'autre, comme ça je suis doublement protégé, et je les garde dans mon portefeuille pour les avoir tout le temps. » En t'appuyant sur les connaissances du cours, rédige une argumentation scientifique pour corriger ces deux idées reçues.

\tcblower
\textbf{Solution détaillée :}
\begin{enumerate}
    \item \textbf{Double préservatif (superposition) :} \textbf{FAUX et DANGEREUX.} Le frottement entre les deux couches de latex (ou latex/nitrile) pendant les mouvements du rapport génère de la chaleur et une usure mécanique accélérée. Cela \textbf{augmente très significativement le risque de rupture simultanée des deux préservatifs}. Un seul préservatif, correctement mis, offre une résistance mécanique suffisante (norme ISO 4074 : test de gonflage $>$ \SI{18}{\liter} d'air, test d'éclatement $>$ \SI{1}{\kilo\pascal}). La superposition annule cette garantie.
    \item \textbf{Conservation dans le portefeuille :} \textbf{FAUX et DANGEREUX.} Le portefeuille subit : (a) une température élevée (proche du corps, \textasciitilde \SI{37}{\celsius}), (b) une pression mécanique constante (assise, pliage), (c) des frottements contre les cartes/pièces. Ces trois facteurs dégradent les chaînes polymères du latex (oxydation thermique, micro-fissures par fatigue mécanique). Un préservatif ainsi stocké pendant plusieurs semaines a une probabilité de rupture à l'usage multipliée par 3 à 5. \textbf{Conservation correcte :} Boîte rigide (étui à préservatifs) dans un tiroir ou sac à dos (température ambiante, pas de compression).
\end{enumerate}
\textbf{Conclusion :} Ces pratiques diminuent la protection au lieu de l'augmenter. L'efficacité repose sur \textbf{un seul préservatif, neuf, bien conservé, bien mis}.
\end{exoresolu}

\begin{savaistu}[Histoire et contexte : Du boyau au latex]
Les premiers « préservatifs » documentés (Égypte antique, Italie XVIe siècle - Fallope) étaient en boyau d'animal (mouton), lin trempé dans des solutions chimiques, ou écaille de tortue. Ils servaient surtout contre la syphilis. Le caoutchouc vulcanisé (Goodyear, 1839) permet les premiers modèles réutilisables (épais, couture). Le \textbf{latex} (hévéa naturel, traité par trempage) révolutionne le domaine dans les années 1930 : fin, élastique, jetable, peu cher. Au Cameroun, la promotion du préservatif s'intensifie dès les années 1987-1990 avec l'épidémie de sida. Aujourd'hui, la distribution gratuite ou subventionnée (programmes CNLS, Fonds Mondial, UNFPA) en fait un outil de santé publique majeur. Le préservatif féminin (polyuréthane puis nitrile) arrive dans les années 1990 (FC1 puis FC2), offrant une alternative non-latex et gérée par la femme.
\end{savaistu}

\courssub{Synthèse : Le préservatif, outil d'asepsie sexuelle et de responsabilité}

\begin{aretenir}[Bilan : Ce qu'il faut retenir sur le préservatif]
\begin{itemize}
    \item \textbf{Nature :} Barrière physique (latex, nitrile, polyuréthane) à usage unique.
    \item \textbf{Principe :} \textbf{Asepsie mécanique} : empêche le contact fluides/muqueuses et le passage des microbes (VIH, bactéries IST) et des spermatozoïdes.
    \item \textbf{Double protection :} Seule méthode contraceptive protégeant aussi contre les IST (dont sida).
    \item \textbf{Condition d'efficacité :} \textbf{Usage correct} (procédure stricte) \textbf{ET systématique} (à chaque rapport, du début à la fin).
    \item \textbf{Erreurs fatales :} Réutilisation, double usage, lubrifiant gras, conservation chaleur/pression, oubli du réservoir, mise en place tardive/retrait précoce.
    \item \textbf{Complémentarité :} S'inscrit dans une stratégie globale : Abstinence / Fidélité dépistée / Préservatif / Dépistage / Vaccination / Dialogue sans tabou.
\end{itemize}
\end{aretenir}

\begin{exercice}[Application : Cas pratiques]
\textbf{1.} Marie (17 ans) utilise la pilule contraceptive (œstro-progestative) prise régulièrement. Elle a un nouveau petit ami. Doivent-ils utiliser le préservatif ? Justifie scientifiquement (deux arguments distincts).
\textbf{2.} Un préservatif est resté 6 mois dans la boîte à gants d'une voiture à Yaoundé (température moyenne \SI{35}{\celsius}). Est-il utilisable ? Explique au niveau moléculaire (polymère).
\textbf{3.} Lors d'un rapport, le préservatif masculin glisse et reste dans le vagin après le retrait du pénis. Quelle est la conduite à tenir immédiate pour la prévention : (a) Grossesse ? (b) IST/VIH ? Cite les services de santé compétents au Cameroun (nom générique).
\textbf{4.} Pourquoi le préservatif protège-t-il moins bien contre le HPV (papillomavirus) et l'herpès génital que contre le VIH ou la blennorragie ? (Indice : mode de transmission).
\end{exercice}

\begin{corrige}[Exercice : Cas pratiques]
\textbf{1.} \textbf{OUI, indispensable.}
    \begin{itemize}
        \item Argument 1 (IST) : La pilule ne protège \textbf{pas du tout} contre les IST (VIH, hépatites, blennorragie, syphilis, HPV, herpès). Seul le préservatif assure cette barrière.
        \item Argument 2 (Sécurité contraceptive) : L'efficacité réelle de la pilule (usage courant) est \textasciitilde 91\% (oubli de prise, vomissements, interactions médicamenteuses). L'association « Pilule + Préservatif » (double protection) ramène le risque de grossesse quasi à zéro et couvre le risque IST.
    \end{itemize}
\textbf{2.} \textbf{NON, inutilisable (dangereux).} Le latex (polyisoprène \textit{cis}-1,4) est un polymère sensible à l'oxydation thermique. À \SI{35}{\celsius} constant pendant 6 mois, les chaînes polymères se cassent (cisaillage thermique) et s'oxydent (radicaux libres $\rightarrow$ rupture liaisons C-C). Le matériau devient cassant, perd son élasticité (module de Young augmenté) et présente des micro-fissures invisibles. Il éclatera sous la pression éjaculatoire (\SIrange{5}{10}{\kilo\pascal}). \textbf{Action :} Jeter, racheter.
\textbf{3.} Conduite à tenir immédiate (dans les 72h max, idéalement < 24h) :
    \begin{itemize}
        \item (a) \textbf{Grossesse :} Contraception d'urgence (Lévonorgestrel \SI{1,5}{mg} dose unique ou Ulipristal \SI{30}{mg}). Disponible sans ordonnance en pharmacie, gratuite dans les centres de santé publics pour les mineurs. Efficacité maximale si prise tout de suite.
        \item (b) \textbf{IST/VIH :} Se rendre \textbf{immédiatement} (idéalement < 4h, max 48-72h) dans un \textbf{Centre de Prise en Charge (CPEC)} ou un \textbf{Centre de Dépistage Volontaire (CDV)} ou service des urgences d'un hôpital de district/régional pour évaluation \textbf{PEP (Prophylaxie Post-Exposition)} : traitement antirétroviral d'urgence sur 28 jours (si statut source inconnu ou positif). Dépistage de base (VIH, VHB, Syphilis) et vaccination VHB si non immunisé. Traitement présomptif IST bactériennes (Céfixime + Azithromycine ou Doxycycline) selon protocole national.
    \end{itemize}
\textbf{4.} Le VIH et la blennorragie/syphilis (stade primaire/secondaire) se transmettent essentiellement par \textbf{fluides} (sperme, sang, sécrétions vaginales) $\rightarrow$ bloqués par la barrière continue du préservatif. Le HPV et l'Herpès (HSV-2) se transmettent par \textbf{contact direct peau-muqueuse / muqueuse-muqueuse} au niveau des zones \textbf{non couvertes} par le préservatif masculin (base du pénis, bourses, pubis, cuisses, périnée, vulve, anus). Le préservatif féminin (couvrant la vulve) offre une meilleure protection contre ces IST cutanées, mais pas totale.
\end{corrige}

\coursec{Synthèse et application à la vie courante}

Après avoir étudié l'utilisation des préservatifs comme barrière mécanique contre les contaminations sexuelles, nous disposons maintenant des trois piliers de la prévention primaire : l'\textbf{asepsie}, l'\textbf{antisepsie} et la \textbf{protection par préservatif}. Cette section finale a pour but de synthétiser ces notions, de t'apprendre à les mobiliser face à des situations réelles et de mesurer leur impact concret sur la santé publique, ici au Cameroun comme ailleurs.

\courssub{Bilan récapitulatif des trois pratiques préventives}

Pour choisir la bonne pratique, il faut d'abord bien distinguer leur objectif, leur cible et leur domaine d'application. Le tableau ci-dessous résume les points essentiels à retenir pour le BEPC.

\begin{center}
\begin{tabularx}{\linewidth}{|l|X|X|X|}
\hline
\rowcolor{popblueL}
\textbf{Pratique} & \textbf{Objectif principal} & \textbf{Action sur le microbe} & \textbf{Principaux domaines d'application} \\
\hline
\textbf{Asepsie} & Empêcher l'arrivée des microbes (prévention de la contamination) & Aucune action directe : crée une barrière physique ou chimique stérile & Chirurgie, soins invasifs (perfusion, cathéter), préparation des aliments (hygiène cuisine), manipulations en laboratoire \\
\hline
\textbf{Antisepsie} & Détruire ou inhiber les microbes déjà présents sur un vivant (peau, muqueuse) & Action chimique (désinfectant) ou physique (chaleur) sur la flore microbienne & Désinfection cutanée avant injection/incision, lavage des mains au savon ou solution hydroalcoolique, soin d'une plaie, hygiène bucco-dentaire \\
\hline
\textbf{Préservatif (masculin/féminin)} & Bloquer le passage des agents pathogènes (VIH, IST, spermatozoïdes) lors des rapports sexuels & Barrière mécanique étanche (latex, polyuréthane) & Relations sexuelles vaginales, anales, orales ; prévention des IST/VIH et contraception \\
\hline
\end{tabularx}
\end{center}

\begin{aretenir}[Distinction fondamentale Asepsie / Antisepsie]
L'\textbf{asepsie} s'applique sur des \textbf{objets, des milieux ou des surfaces inertes} (instruments, champs opératoires, mains avant stérilisation) pour les maintenir \textbf{stériles} (absence de tout germe vivant). L'\textbf{antisepsie} s'applique \textbf{uniquement sur les tissus vivants} (peau, muqueuses) pour \textbf{réduire la charge microbienne} ; elle ne stérilise pas (la flore résidente persiste).
\end{aretenir}

\courssub{Démarche de résolution d'une situation-problème de prévention}

Face à une situation concrète (examen, vie courante), ne devine pas : suis cette démarche rigoureuse en trois étapes. C'est la clé pour justifier tes réponses et obtenir les points au BEPC.

\begin{methode}[Analyser et justifier une mesure de prévention]
\textbf{Étape 1 : Identifier le risque et la porte d'entrée.}
Lis attentivement l'énoncé. Quel est le microbe suspecté (bactérie, virus, champignon) ? Par où peut-il pénétrer dans l'organisme ? (Porte d'entrée : peau lésée, muqueuse respiratoire, digestive, génitale, voie sanguine directe).

\textbf{Étape 2 : Choisir la pratique adaptée.}
\begin{itemize}
    \item Si le risque vient d'un \textbf{objet/instrument} qui va pénétrer le corps $\rightarrow$ \textbf{Asepsie} (stérilisation, champ stérile, matériel à usage unique).
    \item Si le risque vient de \textbf{la peau ou d'une muqueuse saine ou souillée} (mains, site d'injection, plaie) $\rightarrow$ \textbf{Antisepsie} (lavage, désinfectant).
    \item Si le risque est une \textbf{contamination sexuelle} (rapport sexuel) $\rightarrow$ \textbf{Préservatif} (barrière mécanique).
\end{itemize}

\textbf{Étape 3 : Justifier par les définitions du cours.}
Rédige une phrase complète : « On utilise l'antisepsie ici car il s'agit de détruire les microbes présents sur la peau (tissu vivant) avant une piqûre, ce qui correspond à la définition de l'antisepsie. » N'oublie pas le vocabulaire précis : \emph{stériliser}, \emph{désinfecter}, \emph{barrière mécanique}.
\end{methode}

\begin{attention}[Piège classique : Confondre prévention et traitement]
\textbf{Erreur fréquente :} Proposer la prise d'antibiotiques \emph{avant} ou \emph{juste après} un rapport à risque ou une blessure pour « prévenir l'infection ».
\textbf{Correction :} Les \textbf{antibiotiques} sont des \textbf{traitements curatifs} (ils tuent les bactéries \emph{après} l'installation de l'infection). Ils ne sont \textbf{pas} des moyens de prévention primaire (sauf prophylaxie médicale très encadrée, ex : prophylaxie post-exposition au VIH - PPE - qui utilise des antirétroviraux, pas des antibiotiques classiques). L'asepsie, l'antisepsie et le préservatif agissent \textbf{avant} l'installation de l'infection. Ne confonds jamais \emph{prévenir} (éviter l'arrivée) et \emph{guérir} (éliminer l'installé).
\end{attention}

\courssub{Analyse de situations concrètes de la vie courante}

Appliquons la méthode à quatre situations typiques du quotidien camerounais. Pour chacune, identifie la porte d'entrée, la pratique choisie et la justification.

\begin{exoresolu}[Situation 1 : La vendeuse de « beignets haricots » au carrefour]
\textbf{Énoncé :} Marie vend des beignets chauds au carrefour Nlongkak à Yaoundé. Elle manipule la pâte, l'huile bouillante, l'argent et sert les clients sans se laver les mains entre deux opérations. Quel risque ? Quelle pratique préventive appliquer ?

\textbf{Solution détaillée :}
\begin{enumerate}
    \item \textbf{Porte d'entrée :} Voie digestive des clients (ingestion des beignets contaminés).
    \item \textbf{Microbes sources :} Mains de Marie (flore transitoire : \emph{E. coli}, \emph{Salmonella}, \emph{Staphylococcus aureus} provenant de l'argent, du nez, des selles si hygiène défaillante).
    \item \textbf{Pratique adaptée :} \textbf{Asepsie} (hygiène alimentaire / manipulation aseptique) \textbf{ET Antisepsie} (lavage des mains).
    \item \textbf{Justification :}
    \begin{itemize}
        \item L'argent et l'environnement ne sont pas stériles. Pour \textbf{empêcher l'arrivée} des microbes sur l'aliment (pâte cuite), Marie doit respecter l'\textbf{asepsie} : ustensiles propres, protection des aliments (vitrine), séparation argent/aliment (utiliser une pince ou un sachet).
        \item Surtout, le \textbf{lavage des mains au savon} (antisepsie mécanique + chimique) \textbf{détruit/élimine} la flore transitoire avant de toucher la pâte cuite.
    \end{itemize}
    \item \textbf{Bilan :} L'association Asepsie (protection de l'aliment) + Antisepsie (mains propres) brise la chaîne de contamination fécale-orale.
\end{enumerate}
\tcblower
\textbf{Point méthode :} Ici, deux pratiques s'imbriquent. L'asepsie protège le \emph{milieu récepteur} (le beignet), l'antisepsie traite le \emph{vecteur} (la main). Au BEPC, cite les deux si la situation le justifie.
\end{exoresolu}

\begin{exoresolu}[Situation 2 : Blessure à la machette au champ de cacao]
\textbf{Énoncé :} Paul se coupe le mollet avec une machette rouillée en entretenant son champ de cacao à Ngomedzap. La plaie est sale (terre, débris végétaux). Il est seul, loin du centre de santé. Que faire immédiatement pour éviter l'infection (tétanos, gangrène gazeuse) ?

\textbf{Solution détaillée :}
\begin{enumerate}
    \item \textbf{Porte d'entrée :} Peau lésée (plaie ouverte) $\rightarrow$ voie directe vers le sang et les tissus profonds.
    \item \textbf{Microbes sources :} Terre (\emph{Clostridium tetani}, \emph{Clostridium perfringens}), rouille, flore cutanée.
    \item \textbf{Pratique adaptée :} \textbf{Antisepsie} (nettoyage + désinfection de la plaie).
    \item \textbf{Justification :} La plaie est un \textbf{tissu vivant} contaminé. On ne peut pas la stériliser (asepsie impossible sur tissu vivant). Il faut \textbf{détruire les microbes présents} dans la plaie sans nécroser les tissus sains.
    \item \textbf{Protocole de terrain (antisepsie) :}
    \begin{enumerate}
        \item Nettoyer à l'eau potable (ou eau de source bouillie/filtrée) + savon pour enlever la terre (action mécanique majeure).
        \item Rincer abondamment.
        \item Appliquer un antiseptique \textbf{non irritant, non coloré} (ex : Dakin dilué, Bétadine dermique 10\% iodée, Chlorhexidine 0,05\%). \textbf{Éviter l'alcool 70° ou l'eau oxygénée pur sur plaie ouverte} (douleur, nécrose, bulles d'O$_2$ favorisant anaérobies).
        \item Couvrir d'un pansement stérile (asepsie du pansement) ou compresses propres.
    \end{enumerate}
    \item \textbf{Urgence :} Évacuation au centre de santé pour rappel antitétanique (sérovaccination) si vaccination non à jour. L'antisepsie ne remplace pas la prophylaxie vaccinale !
\end{enumerate}
\tcblower
\textbf{Attention :} « Mettre de l'alcool sur la plaie » est une erreur fréquente (douleur intense, destruction tissulaire, inefficacité sur spores de \emph{Clostridium} dans la profondeur). L'eau + savon est la première antisepsie vitale.
\end{exoresolu}

\begin{exoresolu}[Situation 3 : Accouchement en maternité (salle de travail)]
\textbf{Énoncé :} Une sage-femme au Centre de Santé Intégré de Bafia prépare l'accouchement d'une primipare. Elle met en place le champ opératoire, sort les instruments du stérilisateur, se lave les mains à la brosse povidone-iodée, met des gants stériles. Classe chaque geste : Asepsie ou Antisepsie ?

\textbf{Solution détaillée :}
\begin{center}
\begin{tabularx}{\linewidth}{|l|X|l|}
\hline
\rowcolor{popgreenL}
\textbf{Geste} & \textbf{Classification} & \textbf{Justification (Cible)} \\
\hline
Stérilisation des instruments (autoclave 121°C, 20 min) & \textbf{Asepsie} & Objets inertes $\rightarrow$ stérilité totale \\
Champ opératoire stérile (draps) & \textbf{Asepsie} & Surface inerte $\rightarrow$ barrière stérile \\
Gants stériles (usage unique) & \textbf{Asepsie} & Objet inerte $\rightarrow$ barrière stérile main/patiente \\
Lavage mains sage-femme (brosse + povidone) & \textbf{Antisepsie} & Peau vivante (mains) $\rightarrow$ réduction flore transitoire/résidente \\
Désinfection vulve/périnée patiente (povidone) & \textbf{Antisepsie} & Muqueuse/peau vivante $\rightarrow$ réduction flore locale \\
\hline
\end{tabularx}
\end{center}
\textbf{Synthèse :} L'\textbf{asepsie} concerne tout ce qui est \textbf{apporté de l'extérieur} (instruments, champs, gants) et doit rester stérile. L'\textbf{antisepsie} concerne les \textbf{surfaces vivantes} (mains soignante, peau patiente) qui ne peuvent être stérilisées. Les deux sont \textbf{indissociables} pour prévenir l'infection du site opératoire (endométrite, septicémie néonatale).
\end{exoresolu}

\begin{exoresolu}[Situation 4 : Rapport sexuel occasionnel non protégé]
\textbf{Énoncé :} Kevin, 24 ans, étudiant à Douala, a un rapport sexuel vaginal non protégé avec une partenaire dont il ne connaît pas le statut sérologique. Il craint le VIH/Sida et les IST. Quelle pratique de prévention a été oubliée ? Quelle conduite tenir maintenant ?

\textbf{Solution détaillée :}
\begin{enumerate}
    \item \textbf{Pratique oubliée (Prévention primaire) :} Utilisation du \textbf{préservatif masculin (ou féminin)}.
    \item \textbf{Justification :} La porte d'entrée est la \textbf{muqueuse génitale} (pénis, vagin, col utérin). Le préservatif est la \textbf{seule} pratique qui crée une \textbf{barrière mécanique étanche} empêchant le contact direct entre les liquides biologiques (sperme, sécrétions vaginales, sang) porteurs du VIH, HBV, HCV, \emph{Treponema pallidum}, \emph{Neisseria gonorrhoeae}, \emph{Chlamydia trachomatis}, HPV, etc.
    \item \textbf{Conduite immédiate (Post-exposition - Prévention secondaire) :} Se rendre \textbf{au plus vite (idéalement < 4h, max 48-72h)} dans un Centre de Prise en Charge (CPEC) ou Centre de Dépistage Volontaire (CDV) pour une \textbf{Prophylaxie Post-Exposition (PPE)} au VIH (trithérapie antirétrovirale 28 jours) + dépistage IST + contraception d'urgence si féminin.
    \item \textbf{Distinction cruciale :} Le préservatif = \textbf{Prévention PRIMAIRE} (empêche l'arrivée). PPE = \textbf{Prévention SECONDAIRE} (traitement d'urgence après risque pour empêcher l'installation). L'antisepsie génitale (toilette, uriner) \textbf{ne protège PAS} du VIH/IST (virus déjà dans les cellules/muqueuses).
\end{enumerate}
\tcblower
\textbf{Message clé :} « Mieux vaut prévenir que guérir. » Le préservatif est gratuit dans les structures publiques camerounaises. Son usage correct et systématique reste la référence.
\end{exoresolu}

\courssub{Arbre de décision : de la situation à la pratique adaptée}

Pour t'aider à visualiser le raisonnement, voici un arbre de décision à mémoriser. Il résume la logique « Porte d'entrée $\rightarrow$ Pratique ».

\begin{popfigure}
\begin{tikzpicture}[
    node distance=1.2cm and 1.5cm,
    decision/.style={diamond, draw=popblue, thick, fill=popblueL, text width=3.5cm, align=center, font=\small},
    action/.style={rectangle, rounded corners, draw=popgreen, thick, fill=popgreenL, text width=3cm, align=center, font=\small\bfseries},
    start/.style={ellipse, draw=poporange, thick, fill=poporangeL, text width=2.5cm, align=center, font=\small\bfseries},
    arrow/.style={-Stealth, thick, popdark},
    label/.style={font=\footnotesize, popmuted, sloped, pos=0.5, above}
]
    % Nœuds
    \node[start, align=center] (start){Situation à risque\\contamination};
    \node[decision, below of=start, align=center] (q1){Le risque vient-il d'un\\objet/instrument/milieu\\inerte qui va toucher\\l'intérieur du corps ?};
    \node[action, below left=1.5cm and -1cm of q1, align=center] (a1){ASEPSIE\\Stérilisation, champ stérile,\\matériel usage unique};
    \node[decision, below right=1.5cm and -1cm of q1, align=center] (q2){Le risque vient-il d'une\\peau/muqueuse vivante\\(mains, site injection,\\plaie, périnée) ?};
    \node[action, below left=1.5cm and -1cm of q2, align=center] (a2){ANTISEPSIE\\Lavage savon, solution\\hydroalcoolique, désinfectant\\cutané/muqueux};
    \node[decision, below=2.5cm of q1, align=center] (q3){Le risque est-il une\\relation sexuelle\\vaginale/anale/orale ?};
    \node[action, below left=1.5cm and -1cm of q3, align=center] (a3){PRÉSERVATIF\\Barrière mécanique\\latex/polyuréthane};
    \node[action, below right=1.5cm and -1cm of q3, align=center] (a4){AUTRE PRÉVENTION\\Vaccination, hygiène de vie,\\dépistage, PPE (post-expo)};

    % Flèches
    \draw[arrow] (start) -- (q1);
    \draw[arrow] (q1) -- node[label, near start, anchor=east] {OUI} (a1);
    \draw[arrow] (q1) -- node[label, near start, anchor=west] {NON} (q2);
    \draw[arrow] (q2) -- node[label, near start, anchor=east] {OUI} (a2);
    \draw[arrow] (q2) -- node[label, near start, anchor=west] {NON} (q3);
    \draw[arrow] (q3) -- node[label, near start, anchor=east] {OUI} (a3);
    \draw[arrow] (q3) -- node[label, near start, anchor=west] {NON} (a4);
\end{tikzpicture}
\legende{Arbre de décision pour choisir la pratique préventive adaptée (Asepsie, Antisepsie, Préservatif) en fonction de la nature du risque et de la porte d'entrée du microbe.}
\end{popfigure}

\courssub{Bilan de la leçon et ouverture}

\begin{aretenir}[Bilan : Trois gestes simples, accessibles, efficaces]
Éviter la contamination et l'infection ne relève pas de la haute technologie, mais de \textbf{gestes quotidiens rigoureux} :
\begin{enumerate}
    \item \textbf{L'asepsie} (stériliser, protéger les champs, hygiène alimentaire) : \emph{« Zéro germe sur l'objet »}.
    \item \textbf{L'antisepsie} (laver, désinfecter le vivant) : \emph{« Moins de germes sur la peau/muqueuse »}.
    \item \textbf{Le préservatif} (barrière sexuelle) : \emph{« Pas de contact des liquides »}.
\end{enumerate}
Leur efficacité n'est pas théorique : elle est \textbf{démontrée expérimentalement} (TP N°6 : boîte de Pétri stérile vs contaminée ; études épidémiologiques : lavage mains $-40\%$ diarrhées ; préservatif $>95\%$ efficacité VIH si usage correct).
\end{aretenir}

\begin{exemplebox}[Donnée chiffrée : L'impact du lavage des mains]
Selon l'OMS et de multiples études de terrain (dont au Cameroun dans les zones rurales et urbaines), la promotion du \textbf{lavage des mains au savon aux moments critiques} (après défécation, avant manger, avant cuisiner) permet de réduire :
\begin{itemize}
    \item d'environ \textbf{40 \%} l'incidence des \textbf{maladies diarrhéiques} (2ème cause de mortalité < 5 ans au Cameroun).
    \item d'environ \textbf{20-25 \%} les \textbf{infections respiratoires aiguës} (pneumonies).
\end{itemize}
C'est l'intervention de santé publique au \textbf{rapport coût-efficacité le plus élevé} qui existe. Un savon coûte $\sim$ 150-200 FCFA ; une hospitalisation pour choléra ou dysenterie en coûte des milliers.
\end{exemplebox}

\begin{savaistu}[L'OMS et le « geste qui sauve »]
L'Organisation Mondiale de la Santé (OMS) considère le \textbf{lavage des mains} comme \textbf{« la mesure unique la plus importante, la plus simple et la moins coûteuse pour prévenir la transmission des infections »}. La Journée Mondiale du Lavage des Mains (15 octobre) est célébrée dans les écoles camerounaises. En 2020-2021, face au COVID-19, les « lave-mains » (seaux à robinet + savon/eau de Javel diluée) sont devenus la norme à l'entrée des marchés, écoles et administrations à Yaoundé, Douala, Bafoussam... Preuve que l'antisepsie simple (eau + savon) est une arme de masse contre les épidémies.
\end{savaistu}

\courssub{Ouverture : Vers les défenses de l'organisme (Leçons suivantes)}

Ces trois pratiques (asepsie, antisepsie, préservatif) constituent la \textbf{prévention primaire} : elles agissent \textbf{à l'extérieur} de l'organisme pour \textbf{bloquer l'entrée} du microbe. Elles sont notre \textbf{première ligne de défense}.

Mais si, malgré tout, un microbe franchit ces barrières (blessure non soignée, rapport non protégé, aliment contaminé, aérosol inhalé), l'organisme n'est pas désarmé. Il possède ses \textbf{propres systèmes de défense} :
\begin{itemize}
    \item Les \textbf{barrières naturelles} (peau intacte, cils vibratiles, mucus, pH acide, flore commensale) — \emph{leçon suivante}.
    \item La \textbf{réponse inflammatoire} (non spécifique, immédiate).
    \item L'\textbf{immunité spécifique} (lymphocytes B/T, anticorps, mémoire immunitaire) — \emph{leçons sur l'immunité}.
    \item La \textbf{vaccination} (éducation du système immunitaire \emph{avant} la rencontre).
\end{itemize}

Comprendre comment \textbf{nous aidons notre corps à se défendre} (hygiène, vaccins) et comment \textbf{notre corps se défend tout seul} (immunité) est le fil conducteur de toute cette unité « Microorganismes et immunité ». La leçon prochaine entrera dans le vif du sujet : \textbf{les barrières naturelles et la réponse inflammatoire}. Prépare-toi à observer la guerre microscopique qui se joue en toi à chaque instant !

\begin{exercice}[Entraînement BEPC : Situations mixtes]
Pour chacune des situations ci-dessous, indique la pratique préventive principale (Asepsie, Antisepsie, Préservatif) et justifie en une phrase en utilisant le vocabulaire du cours (porte d'entrée, cible, action).
\begin{enumerate}
    \item Un infirmier prépare une perfusion : il désinfecte le site de ponction du bras du patient à la chlorhexidine.
    \item Au laboratoire, l'enseignant passe les boucles de culture à la flamme du bec Bunsen avant d'ensemencer une gélose.
    \item Une mère de famille lave soigneusement les feuilles de « ndolè » (vernonia) à l'eau additionnée de quelques gouttes d'eau de Javel avant de les cuisiner.
    \item Un couple utilise un préservatif masculin lors de chaque rapport sexuel.
    \item Un chirurgien met des gants stériles après s'être lavé les mains à la brosse iodée.
\end{enumerate}
\end{exercice}

\begin{corrige}[Exercice n°1 - Entraînement BEPC]
\begin{enumerate}
    \item \textbf{Antisepsie.} La peau du bras (tissu vivant, porte d'entrée potentielle veineuse) est désinfectée pour détruire la flore microbienne locale avant l'effraction cutanée.
    \item \textbf{Asepsie.} La boucle de culture (objet inerte, matériel) est stérilisée à la flamme pour empêcher l'arrivée de contaminants sur le milieu de culture (milieu stérile).
    \item \textbf{Asepsie (hygiène alimentaire).} L'eau de Javel (désinfectant) traite l'eau de lavage (milieu inerte) et la surface des feuilles (considérées comme surface inerte alimentaire) pour empêcher l'arrivée de microbes (kystes amibes, bactéries) dans l'aliment consommé cru ou peu cuit. \textit{Note : On parle souvent d'asepsie alimentaire pour le traitement de l'eau/légumes, bien que l'eau de Javel soit un antiseptique/désinfectant ; la cible est l'aliment/milieu, pas un tissu vivant du patient.}
    \item \textbf{Préservatif.} Barrière mécanique sur le pénis (muqueuse vivante) empêchant le contact des liquides biologiques (porte d'entrée muqueuse génitale) lors du rapport sexuel.
    \item \textbf{Double pratique :} Lavage mains brossé iodé = \textbf{Antisepsie} (peau vivante main soignant). Gants stériles = \textbf{Asepsie} (objet inerte barrière stérile entre main soignant et site opératoire stérile).
\end{enumerate}
\end{corrige}

\newpage

\coursec{Activité d'intégration}

\coursec{Leçon 5 : Pratiques pour éviter la contamination et l'infection}

\courssub{Objectifs de l'activité}
\begin{itemize}
    \item Exploiter des données historiques (Semmelweis) et expérimentales (empreintes de mains) pour mettre en évidence l'efficacité de l'asepsie (lavage des mains) et de l'antisepsie (désinfection).
    \item Distinguer clairement les concepts d'\textbf{asepsie} (empêcher l'arrivée des microbes), d'\textbf{antisepsie} (détruire les microbes présents) et de \textbf{protection par préservatif} (barrière mécanique).
    \item Appliquer ces connaissances à des situations concrètes de la vie courante au Cameroun (hygiène alimentaire, prise en charge d'une plaie, prévention des IST/VIH).
    \item Rédiger des fiches de conseils de prévention justifiées scientifiquement, en mobilisant le vocabulaire précis du programme.
\end{itemize}

\courssub{Situation : Enquête à la maternité — Pourquoi le lavage des mains sauve des vies ?}

\textbf{Contexte historique : Vienne, 1847.} Le Dr Ignaz Semmelweis dirige deux services de maternité à l'Hôpital général de Vienne. Le \textbf{Service 1} est géré par des médecins et des étudiants qui passent directement de l'autopsie de cadavres (salle de dissection) aux accouchements, sans se laver les mains. Le \textbf{Service 2} est géré par des sages-femmes qui ne pratiquent pas d'autopsies. Semmelweis observe une différence frappante de mortalité par \emph{fièvre puerpérale} (infection grave après l'accouchement). Il impose alors le lavage des mains avec une solution chlorée (eau de Javel) avant tout examen.

\textbf{Document 1 : Taux de mortalité des accouchées à la maternité de Vienne (1841--1848)}

\begin{center}
\begin{tabularx}{\linewidth}{|l|X|X|X|}\hline
\rowcolor{popblueL}
\textbf{Année} & \textbf{Service 1 (Médecins/Étudiants)} & \textbf{Service 2 (Sages-femmes)} & \textbf{Observation} \\ \hline
1841 & 7,8 \% & 3,5 \% & Avant intervention \\ \hline
1842 & 15,8 \% & 6,9 \% &  \\ \hline
1843 & 12,7 \% & 7,8 \% &  \\ \hline
1844 & 11,4 \% & 8,2 \% &  \\ \hline
1845 & 10,3 \% & 7,1 \% &  \\ \hline
1846 & 11,4 \% & 7,2 \% &  \\ \hline
\rowcolor{poporangeL}
1847 (Mai--Déc) & \textbf{2,7 \%} & 6,3 \% & \textbf{Instauration lavage mains (Service 1)} \\ \hline
1848 & 1,3 \% & 5,1 \% & Maintien de la mesure \\ \hline
\end{tabularx}
\end{center}
\begin{popfigure}
\legende{Document 1 : Évolution de la mortalité par fièvre puerpérale. Source : Semmelweis, \emph{Die Ätiologie, der Begriff und die Prophylaxis des Kindbettfiebers}, 1861.}
\end{popfigure}

\textbf{Document 2 : Résultats du TP N°6 -- Comparaison d'empreintes de mains sur gélose nutritive (Incubation 24h à 37°C)}

\begin{center}
\begin{tabularx}{\linewidth}{|l|X|X|}\hline
\rowcolor{popblueL}
\textbf{Condition} & \textbf{Aspect de la boîte de Pétri} & \textbf{Nombre moyen de colonies (UFC/boîte)} \\ \hline
Mains non lavées (témoin) & Nombreuses colonies variées (blanches, jaunes, filamenteuses), recouvrement > 80\% & 145 $\pm$ 12 \\ \hline
Mains lavées à l'eau savonneuse (30 s, frottement) & Colonies rares, isolées, recouvrement < 5\% & 8 $\pm$ 3 \\ \hline
Mains lavées + friction hydroalcoolique (30 s) & Aucune colonie visible (stérile) ou 0--1 colonie & 0,2 $\pm$ 0,4 \\ \hline
\end{tabularx}
\end{center}
\begin{popfigure}
\legende{Document 2 : Résultats moyens d'une classe de 3ème (n = 30 élèves). UFC = Unité Formant Colonie.}
\end{popfigure}

\textbf{Document 3 : Trois situations camerounaises à analyser}

\begin{enumerate}
    \item \textbf{Situation A -- Hygiène alimentaire (Marché Mokolo, Yaoundé) :} Madame Ngo est vendeuse de beignets (puff-puff) et de haricots bouillis. Elle manipule l'argent, sert les clients, pétrit la pâte et trempe les beignets dans l'huile bouillante avec les mêmes mains. Il n'y a pas de point d'eau sur son étal. Elle essuie parfois ses mains sur son pagne.
    \item \textbf{Situation B -- Prise en charge d'une plaie (Collège bilingue, Douala) :} Lors d'un match de football, Kevin (14 ans) se blesse au genou sur un gravier. La plaie saigne modérément et est souillée par de la terre. L'infirmier de l'école dispose d'eau, de savon, de compresses stériles, de Bétadine\footnote{Bétadine\textregistered\ (povidone-iodée) : antiseptique à base d'iode.} et de pansements.
    \item \textbf{Situation C -- Prévention des IST/VIH (Quartier résidentiel, Bafoussam) :} Paul et Grâce, tous deux âgés de 17 ans, sont en couple depuis six mois. Ils envisagent d'avoir des rapports sexuels. Ils ont entendu parler du VIH/Sida et des grossesses précoces au club « Jeunes Leaders » de leur église. Ils ne savent pas concrètement comment se protéger efficacement.
\end{enumerate}

\courssub{Consignes}

\textbf{Partie 1 : Analyse des documents scientifiques (Travail de groupe -- 25 min)}
\begin{enumerate}
    \item À partir du Document 1, calculer la \textbf{diminution moyenne du taux de mortalité} dans le Service 1 entre la période 1841--1846 (moyenne des 6 années) et l'année 1848. Exprimer cette diminution en points de pourcentage et en pourcentage relatif.
    \item Comparer l'évolution du Service 1 et du Service 2 en 1847--1848. Quelle conclusion Semmelweis a-t-il pu tirer sur l'origine de la contamination ?
    \item À partir du Document 2, construire un \textbf{diagramme en barres} comparant le nombre moyen de colonies pour les trois conditions. Utiliser une échelle logarithmique si nécessaire (justifier).
    \item Classer les trois conditions du Document 2 du plus « sale » au plus « propre ». Associer chaque condition à une pratique : \emph{absence de mesure}, \emph{asepsie simple}, \emph{asepsie + antisepsie}.
\end{enumerate}

\textbf{Partie 2 : Application aux situations camerounaises (Travail de groupe -- 20 min)}
\begin{enumerate}\setcounter{enumi}{4}
    \item Pour chaque situation (A, B, C), identifier le(s) risque(s) de \textbf{contamination} (entrée de microbes) et d'\textbf{infection} (développement pathogène).
    \item Proposer des mesures concrètes de prévention pour chaque situation en utilisant \textbf{obligatoirement} les termes : \cle{asepsie}, \cle{antisepsie}, \cle{préservatif} (le cas échéant). Justifier chaque mesure par le mécanisme d'action (barrière physique, destruction chimique, élimination mécanique).
    \item Rédiger une \textbf{fiche conseil} unique (format A5) destinée à être affichée dans l'infirmerie du collège, résumant les 3 grands principes : « Protéger les aliments », « Soigner une plaie », « Se protéger au cours des rapports ». La fiche doit être illustrée de pictogrammes simples.
\end{enumerate}

\textbf{Partie 3 : Restitution et synthèse (Classe entière -- 15 min)}
\begin{enumerate}\setcounter{enumi}{7}
    \item Un rapporteur par groupe présente sa fiche conseil. La classe valide ou corrige le vocabulaire scientifique.
    \item L'enseignant synthétise : définir formellement l'asepsie, l'antisepsie et le rôle du préservatif. Compléter le tableau de synthèse ci-dessous.
\end{enumerate}

\begin{center}
\begin{tabularx}{\linewidth}{|l|X|X|X|}\hline
\rowcolor{popblueL}
\textbf{Pratique} & \textbf{Définition} & \textbf{Mécanisme d'action} & \textbf{Exemple concret (Cameroun)} \\ \hline
Asepsie & & & \\ \hline
Antisepsie & & & \\ \hline
Préservatif & & & \\ \hline
\end{tabularx}
\end{center}

\courssub{Ressources à mobiliser}

\begin{aretenir}[Rappels de cours -- Unité 5]
\begin{enumerate}
\item \textbf{Contamination :} Pénétration de microorganismes (bactéries, virus, champignons) dans l'organisme ou sur un milieu de culture.
\item \textbf{Infection :} Développement et multiplication de ces microorganismes entraînant des symptômes (fièvre, pus, diarrhée, etc.).
\item \textbf{Asepsie :} Ensemble des techniques visant à \emph{empêcher l'arrivée} des microbes sur un milieu stérile ou dans l'organisme (lavage des mains, port de gants stériles, stérilisation du matériel, champ opératoire). \emph{Action : élimination mécanique / barrière physique.}
\item \textbf{Antisepsie :} Ensemble des techniques visant à \emph{détruire ou inhiber} les microbes \emph{déjà présents} sur une peau vivante ou une muqueuse (désinfectants : Bétadine\textregistered, Dakin\textregistered, solution hydroalcoolique, alcool 70°). \emph{Action : chimique (oxydation, dénaturation protéique).} \textbf{Ne jamais utiliser d'antiseptique sur du matériel inerte (stérilisation requise).}
\item \textbf{Préservatif (masculin ou féminin) :} Dispositif médical à usage unique constituant une \textbf{barrière mécanique étanche} empêchant le contact direct entre les muqueuses et les liquides biologiques (sperme, sécrétions vaginales, sang). Seule méthode contraceptive protégeant efficacement contre le VIH/Sida et la plupart des IST.
\end{enumerate}
\end{aretenir}

\begin{methode}[Calcul de pourcentage relatif]
Pour calculer l'évolution relative d'une valeur :
\[
\text{\% d'évolution} = \frac{\text{Valeur finale} - \text{Valeur initiale}}{\text{Valeur initiale}} \times 100
\]
Une diminution donne un nombre négatif.
\end{methode}

\begin{methode}[Diagramme en barres avec échelle logarithmique]
Si les valeurs varient sur plusieurs ordres de grandeur (ex : 145 vs 0,2), une échelle linéaire écrase les petites valeurs. Une échelle logarithmique (base 10) place 1, 10, 100, 1000 à distances égales. En Python/matplotlib : \texttt{plt.yscale('log')}. Sur papier : graduations 0,1 ; 1 ; 10 ; 100.
\end{methode}

\courssub{Démarche et corrigé commenté}

\begin{exoresolu}[Activité d'intégration : Pratiques de prévention]
\textbf{Énoncé -- Consignes de l'activité}

\textbf{Partie 1 : Analyse des documents scientifiques}
\begin{enumerate}
    \item Calculer la diminution moyenne du taux de mortalité (Service 1) entre 1841--1846 et 1848 (en points et en \% relatif).
    \item Comparer l'évolution des Services 1 et 2 en 1847--1848. Conclure sur l'origine de la contamination.
    \item Construire un diagramme en barres (Document 2) avec échelle logarithmique justifiée.
    \item Classer les conditions (plus sale $\to$ plus propre) et associer : absence de mesure / asepsie simple / asepsie + antisepsie.
\end{enumerate}

\textbf{Partie 2 : Application aux situations camerounaises}
\begin{enumerate}\setcounter{enumi}{4}
    \item Identifier risques de contamination et d'infection pour A, B, C.
    \item Proposer mesures de prévention avec termes imposés (\cle{asepsie}, \cle{antisepsie}, \cle{préservatif}) et justification mécanistique.
    \item Rédiger une fiche conseil A5 illustrée (3 panneaux : Aliments, Plaie, Sexualité).
\end{enumerate}

\textbf{Partie 3 : Synthèse}
\begin{enumerate}\setcounter{enumi}{7}
    \item Restitution collective et validation du vocabulaire.
    \item Compléter le tableau de synthèse (Définition, Mécanisme, Exemple Cameroun).
\end{enumerate}

\tcblower

\textbf{Corrigé détaillé}

\textbf{Partie 1 : Analyse des documents scientifiques}

\textbf{Question 1 -- Calculs sur le Document 1.}
\begin{itemize}
    \item \textbf{Moyenne 1841--1846 (Service 1) :}
    \[
    \frac{7,8 + 15,8 + 12,7 + 11,4 + 10,3 + 11,4}{6} = \frac{69,4}{6} \approx 11,57 \%
    \]
    \item \textbf{Taux 1848 (Service 1) :} 1,3 \%.
    \item \textbf{Diminution en points :} $11,57 - 1,3 = 10,27$ points de pourcentage.
    \item \textbf{Diminution relative :}
    \[
    \frac{1,3 - 11,57}{11,57} \times 100 \approx \frac{-10,27}{11,57} \times 100 \approx -88,8 \%
    \]
    \textbf{Interprétation :} Le lavage des mains a réduit la mortalité de près de \textbf{89 \%}.
\end{itemize}

\textbf{Question 2 -- Comparaison Services 1 et 2.}
Avant 1847, le Service 1 (médecins) a une mortalité \textbf{2 à 3 fois plus élevée} que le Service 2 (sages-femmes). En 1847, après instauration du lavage des mains \textbf{uniquement au Service 1}, son taux chute brutalement (11,4 \% $\to$ 2,7 \%) pour devenir \textbf{inférieur} à celui du Service 2 (6,3 \%). En 1848, l'écart se creuse (1,3 \% vs 5,1 \%).
\textbf{Conclusion :} Les médecins/étudiants transportaient des « particules cadaveriques » (microbes) de la salle de dissection aux accouchées via leurs mains souillées. Le lavage des mains (asepsie) supprime ce vecteur de contamination. C'est la preuve expérimentale du rôle des mains dans la transmission nosocomiale.

\textbf{Question 3 -- Diagramme en barres (Document 2).}
Les valeurs sont : 145 ; 8 ; 0,2. L'amplitude est de $145 / 0,2 = 725$ (presque 3 ordres de grandeur).
\textbf{Choix :} \textbf{Échelle logarithmique (base 10) obligatoire.} Sur une échelle linéaire, la barre « 0,2 » serait invisible (0,1 mm pour 145 mm).
\begin{popfigure}
\begin{tikzpicture}
\begin{axis}[
    ymode=log,
    ymin=0.1, ymax=200,
    width=\linewidth, height=7cm,
    ylabel={Nombre moyen de colonies (UFC/boîte) -- Échelle log},
    symbolic x coords={Non lavées, Eau + Savon, Eau + Savon + Gel HA},
    xtick=data,
    x tick label style={rotate=30, anchor=east, font=\small},
    ymajorgrids=true,
    grid style=dashed,
    bar width=25pt,
    nodes near coords,
    nodes near coords align={vertical},
    every node near coord/.append style={font=\tiny, rotate=90, anchor=west, yshift=2pt}
]
\addplot[popblue, fill=popblue!60] coordinates {
    (Non lavées, 145)
    (Eau + Savon, 8)
    (Eau + Savon + Gel HA, 0.2)
};
\end{axis}
\end{tikzpicture}
\legende{Figure 1 : Efficacité comparée du lavage simple et de l'antisepsie sur la flore transitoire des mains. Échelle logarithmique.}
\end{popfigure}

\textbf{Question 4 -- Classement et association.}
\begin{enumerate}
    \item \textbf{Mains non lavées} (145 UFC) : \emph{Absence de mesure} -- Flore transitoire et résidente abondante.
    \item \textbf{Mains lavées eau+savon} (8 UFC) : \emph{Asepsie simple (mécanique)} -- Élimination physique de la flore transitoire (salissures, microbes récents) par frottement + tension de surface du savon. La flore résidente persiste.
    \item \textbf{Mains lavées + gel hydroalcoolique} (0,2 UFC) : \emph{Asepsie + Antisepsie (chimique)} -- Le gel (alcool 60-70\%) dénature les protéines des microbes résiduels (flore résidente + transitoire survivante). Stérilisation quasi-totale de la surface cutanée.
\end{enumerate}

\textbf{Partie 2 : Application aux situations camerounaises}

\textbf{Question 5 -- Risques identifiés.}

\begin{center}
\begin{tabularx}{\linewidth}{|l|X|X|}\hline
\rowcolor{popblueL}
\textbf{Situation} & \textbf{Risque de contamination (Entrée microbes)} & \textbf{Risque d'infection (Développement)} \\ \hline
A -- Vendeuse beignets & Mains souillées (argent, pagne, poussière) $\to$ Aliments (pâte, haricots). Absence de point d'eau = impossibilité asepsie. & Toxi-infections alimentaires collectives (TIAC) : \emph{Salmonella, Staphylococcus aureus, E. coli}. Diarrhées, fièvres chez clients. \\ \hline
B -- Plaie Kevin & Terre, graviers, bactéries peau (Staphylocoques) $\to$ Tissus profonds (derme, muscle). Plaie ouverte = porte d'entrée majeure. & Infection locale (abcès, phlegmon), risque septicémie, tétanos (si non vacciné), gangrène gazeuse (\emph{Clostridium}). \\ \hline
C -- Couple Paul/Grace & Contact muqueuses / liquides biologiques (sperme, sécrétions, sang) sans barrière. & VIH/Sida (incurable), IST (syphilis, hépatite B, gonorrhée, chlamydias, HPV), Grossesse non désirée. \\ \hline
\end{tabularx}
\end{center}

\textbf{Question 6 -- Mesures de prévention justifiées.}

\begin{itemize}
    \item \textbf{Situation A (Alimentaire) :}
    \begin{itemize}
        \item \textbf{Asepsie :} Installation d'un \textbf{lave-mains portable} (bidon à robinet + savon + serviette à usage unique) sur l'étal. Lavage \textbf{obligatoire} après manipulation argent, avant pétrissage, après essuyage mains. Port de \textbf{gants jetables} (changés régulièrement) ou utilisation de pinces/ustensiles pour servir (barrière physique).
        \item \textbf{Antisepsie :} Gel hydroalcoolique en complément (si eau rare), mais \textbf{ne remplace pas le lavage} si mains visiblement sales (graisse, pâte).
        \item \textbf{Cuisson :} L'huile bouillante (180°C) tue les microbes \emph{sur le beignet}, mais ne décontamine pas les mains ni les haricots servis tièdes. \emph{La prévention doit être en amont (asepsie).}
    \end{itemize}
    \item \textbf{Situation B (Plaie) :}
    \begin{itemize}
        \item \textbf{Asepsie :} Infirmier : lavage mains eau+savon + friction gel HA \textbf{avant} et \textbf{après} soin. Port de \textbf{gants stériles} à usage unique. Utilisation de \textbf{compresses stériles} (emballage intact, date non dépassée) pour nettoyer.
        \item \textbf{Antisepsie :} Nettoyage plaie \textbf{de l'intérieur vers l'extérieur} (centrifuge) avec du \textbf{sérum physiologique} ou de l'eau savonneuse pour enlever terre/gravats (mécanique). Puis application \textbf{Bétadine\textregistered\ (povidone-iodée) 10\%} ou \textbf{Dakin\textregistered} (hypochlorite) sur pourtour et fond de plaie. \textbf{Ne pas mettre d'alcool 70° ni d'eau oxygénée} dans la plaie (nécrose tissulaire, douleur intense, retard cicatrisation).
        \item \textbf{Pansement :} Compresse stérile + bande (barrière physique aseptique). Vérification vaccination antitétanique.
    \end{itemize}
    \item \textbf{Situation C (Sexualité) :}
    \begin{itemize}
        \item \textbf{Préservatif masculin (latex/polyisoprène) :} \textbf{Seule méthode} barrière efficace contre VIH/IST + contraception. Utilisation \textbf{à chaque rapport}, de la mise en place au retrait. Vérification date péremption, marque CE/ISO, stockage au frais/sec (pas portefeuille/porte-monnaie). Ouverture soigneuse (pas dents/ciseaux). Pinçage réservoir $\to$ déroulement sur pénis en érection. Retrait immédiat après éjaculation (tenir base).
        \item \textbf{Preservatif féminin (nitrile) :} Alternative, inséré dans vagin avant rapport. Moins disponible au Cameroun.
        \item \textbf{Autres :} Dépistage VIH/IST volontaire et gratuit (Centre de Dépistage Volontaire -- CDV) avant arrêt préservatif si fidélité mutuelle prouvée. Vaccination HPV (filles 9-14 ans, programme EPI Cameroun).
        \item \textbf{Attention :} Pilule, implant, stérilet, retrait, calendrier \textbf{ne protègent PAS contre le VIH/IST}.
    \end{itemize}
\end{itemize}

\textbf{Question 7 -- Fiche Conseil (Modèle pour affichage infirmerie).}

\begin{exemplebox}[Modèle de fiche conseil à afficher en infirmerie]
\textbf{1. Protéger les aliments (asepsie).} Laver les mains (eau et savon, 30 secondes) \textbf{avant} de cuisiner ou de servir, \textbf{après} avoir touché l'argent ou les toilettes. Utensiles propres, gants jetables, eau potable pour laver les légumes. \emph{Barrière physique : les microbes restent dehors.}

\textbf{2. Soigner une plaie.} \emph{Asepsie} : laver les mains, utiliser des gants et des compresses stériles. \emph{Antisepsie} : nettoyer la plaie au sérum physiologique, puis appliquer un antiseptique (Bétadine\textregistered{} ou Dakin\textregistered{}, jamais d'alcool pur). \emph{Détruire les microbes dedans et dessus.}

\textbf{3. Se protéger (sexualité).} Préservatif masculin ou féminin (latex ou nitrile), à \textbf{chaque} rapport, \textbf{du début à la fin} ; vérifier la date et la marque CE ; dépistage VIH/IST gratuit en CDV. \emph{Seul moyen de se protéger du VIH/Sida et des IST.}
\end{exemplebox}

\textbf{Partie 3 : Synthèse -- Tableau récapitulatif complété}

\begin{center}
\begin{tabularx}{\linewidth}{|l|X|X|X|}\hline
\rowcolor{popblueL}
\textbf{Pratique} & \textbf{Définition} & \textbf{Mécanisme d'action} & \textbf{Exemple concret (Cameroun)} \\ \hline
\textbf{Asepsie} & Ensemble des techniques visant à \emph{empêcher l'arrivée} de microorganismes sur un milieu stérile ou dans l'organisme. & \textbf{Élimination mécanique} (lavage, rinçage, essuyage) et \textbf{Barrière physique} (gants stériles, champs opératoires, ustensiles propres, emballages scellés). & Lavage mains vendeuse beignets (eau+savon) ; Port gants stériles infirmier ; Stérilisation matériel chirurgical (autoclave 121°C, 20 min) ; Champ stérile accouchement. \\ \hline
\textbf{Antisepsie} & Ensemble des techniques visant à \emph{détruire ou inhiber} les microorganismes \emph{vivants} présents sur une peau saine, une muqueuse ou une plaie. & \textbf{Action chimique} : oxydation (eau oxygénée, Dakin\textregistered), halogénation (Bétadine\textregistered\ = iode), dénaturation protéique (alcool 60-70\%, gel hydroalcoolique), cations métalliques (chlorhexidine). & Désinfection plaie Kevin (Bétadine\textregistered) ; Friction mains gel HA soignant ; Préparation cutanée avant injection/prise de sang ; Bain de bouche chlorhexidine. \\ \hline
\textbf{Préservatif} & Dispositif médical à usage unique (latex, polyisoprène, nitrile) constituant une barrière mécanique entre les muqueuses et les liquides biologiques. & \textbf{Barrière mécanique étanche} (pores < taille virus VIH $\approx$ 100 nm). Empêche contact sperme/sécrétions/sang $\leftrightarrow$ muqueuses. & Rapport sexuel protégé Paul/Grace ; Prévention VIH/Sida, IST (syphilis, HPV, hépatite B) + Contraception (efficacité 98\% usage parfait). \\ \hline
\end{tabularx}
\end{center}

\end{exoresolu}

\courssub{Bilan de l'activité}

Cette activité d'intégration a permis de :

\begin{enumerate}
    \item \textbf{Valider historiquement et expérimentalement} l'efficacité de l'asepsie (Semmelweis : -89\% mortalité ; TP : division par 18 du nombre de colonies par lavage simple) et de l'antisepsie (gel HA : quasi-stérilisation).
    \item \textbf{Ancrer les définitions officielles} dans des réalités vécues : l'asepsie = « garder le propre propre » (barrière, lavage) ; l'antisepsie = « rendre propre le sale » (désinfectant sur peau/muqueuse/plaie) ; le préservatif = « barrière étanche » (seule protection VIH/IST).
    \item \textbf{Dépasser les idées reçues fréquentes au Cameroun :}
    \begin{itemize}
        \item \textit{« L'alcool 90° ou l'eau de Javel sur une plaie, c'est mieux. »} \textbf{FAUX} : nécrose tissulaire, douleur, retard cicatrisation. On utilise Bétadine\textregistered\ ou Dakin\textregistered\ dilués.
        \item \textit{« Le gel hydroalcoolique remplace le lavage des mains. »} \textbf{FAUX} sur mains visiblement sales (graisse, terre, sang) : le savon élimine la charge organique qui inactiverait l'alcool. Gel = complément ou si pas d'eau.
        \item \textit{« La pilule protège du Sida. »} \textbf{FAUX} : contraception seulement. Préservatif = double protection.
    \end{itemize}
    \item \textbf{Développer des compétences transverses :} calcul de pourcentages (mathématiques), lecture de tableaux/graphiques à échelle logarithmique, rédaction argumentée (français), éducation à la citoyenneté et à la santé (vie courante, genre, droits reproductifs).
\end{enumerate}

\cle{Mots-clés à retenir pour l'évaluation :} \textbf{Asepsie, Antisepsie, Contamination, Infection, Flore transitoire/résidente, Barrière mécanique, Préservatif, Semmelweis, Fièvre puerpérale, Toxi-infection alimentaire, Tétanos, VIH/Sida, IST.}

\newpage

\coursec{Méthodes et exercices résolus}

\courssub{Méthodes et exercices résolus}

\begin{methode}[Différencier l'asepsie et l'antisepsie]
    \textbf{Objectif :} Ne pas confondre ces deux notions fondamentales de la prophylaxie.
    \begin{enumerate}
        \item \textbf{Identifier le but visé :} Empêcher l'arrivée des microbes (asepsie) \textbf{OU} détruire/inhiber les microbes déjà présents (antisepsie).
        \item \textbf{Repérer le support d'action :} Matériel inerte, air, mains saines (asepsie) \textbf{OU} peau vivante, muqueuses, plaie (antisepsie).
        \item \textbf{Citer les moyens physiques/chimiques :} Stérilisation (chaleur, autoclave), filtration, gants stériles, champ opératoire (asepsie) \textbf{OU} Produits chimiques (bétadine, chlorhexidine, Dakin, alcool modifié), lavage simple au savon (antisepsie).
        \item \textbf{Rédiger la définition attendue au BEPC :} Utiliser les termes « empêcher la pénétration » pour l'asepsie et « détruire ou inhiber » pour l'antisepsie.
    \end{enumerate}
    \cle{Asepsie} : ensemble des techniques visant à \textbf{empêcher l'arrivée} de microorganismes sur un milieu stérile ou une zone saine. \cle{Antisepsie} : ensemble des techniques visant à \textbf{détruire ou inhiber} les microorganismes \textbf{sur un milieu vivant} (peau, muqueuse, plaie).
\end{methode}

\begin{exoresolu}[Classification de pratiques hospitalières et domestiques]
    \textbf{Énoncé :} Parmi les pratiques suivantes, cochez celles qui relèvent de l'asepsie (A) et celles qui relèvent de l'antisepsie (S). Justifiez brièvement chaque choix.
    \begin{enumerate}
        \item Stérilisation d'une pince à disséquer à l'autoclave (\SI{121}{\celsius}, \SI{20}{min}).
        \item Application de Bétadine dermique (povidone iodée) sur la peau saine avant une injection.
        \item Port de gants stériles par le chirurgien pour opérer.
        \item Lavage des mains au savon liquide et à l'eau par un infirmier avant un soin.
        \item Passage d'une lame de scalpel à la flamme du bec Bunsen.
        \item Désinfection d'une plaie superficielle au Dakin (hypochlorite de sodium dilué).
    \end{enumerate}
    \tcblower
    \textbf{Solution détaillée :}
    \begin{center}
    \begin{tabularx}{\linewidth}{|c|X|X|}\hline
    \rowcolor{popblueL} \textbf{N°} & \textbf{Classement} & \textbf{Justification} \\ \hline
    1 & \textbf{Asepsie (A)} & La stérilisation à la chaleur humide (autoclave) détruit \emph{tous} les microbes (y compris spores) sur un \textbf{matériel inerte} pour le rendre stérile. C'est une technique de prévention de l'arrivée des germes. \\ \hline
    2 & \textbf{Antisepsie (S)} & La Bétadine est un antiseptique chimique appliqué sur la \textbf{peau vivante} (intègre) pour réduire la flore cutanée résidente et transitoire avant un geste invasif. \\ \hline
    3 & \textbf{Asepsie (A)} & Les gants stériles constituent une \textbf{barrière physique} empêchant le passage des microbes des mains du chirurgien (porteur sain potentiel) vers le site opératoire (milieu stérile). \\ \hline
    4 & \textbf{Antisepsie (S)} & Le lavage simple au savon (même antiseptique) élimine mécaniquement et chimiquement la flore transitoire sur la \textbf{peau vivante} des mains. Il ne stérilise pas les mains. \\ \hline
    5 & \textbf{Asepsie (A)} & Le flambage détruit les microbes sur un \textbf{matériel métallique inerte} (lame) juste avant usage pour garantir sa stérilité. \\ \hline
    6 & \textbf{Antisepsie (S)} & Le Dakin est un antiseptique actif sur les tissus nécrosés et le pus ; il est appliqué sur une \textbf{plaie (milieu vivant lésé)} pour détruire les germes installés. \\ \hline
    \end{tabularx}
    \end{center}
    \textbf{Conclusion :} La distinction fondamentale repose sur la nature du support : \textbf{inerte} pour l'asepsie (stérilisation), \textbf{vivant} pour l'antisepsie (désinfection).
\end{exoresolu}

\begin{methode}[Réaliser un lavage des mains hygiénique et chirurgical (TP N°6)]
    \textbf{Objectif :} Maîtriser le protocole standardisé (recommandations OMS) pour réduire la flore transitoire (lavage hygiénique) ou la flore résidente + transitoire (lavage chirurgical).
    \begin{enumerate}
        \item \textbf{Préparation :} Retirer bijoux, montres, manches remontées au-dessus du coude. Ongles courts, sans vernis.
        \item \textbf{Mouillage :} Mouiller mains et avant-bras (chirurgical) ou mains/poignets (hygiénique) à l'eau tiède.
        \item \textbf{Savonnage (produit) :} Appliquer dose de savon (hygiénique) ou solution hydroalcoolique / savon antiseptique (chirurgical).
        \item \textbf{Friction (6 étapes OMS, \SI{30}{\second} à \SI{60}{\second}) :}
            \begin{itemize}
                \item Paume contre paume.
                \item Paume sur dos de main (entrelacer doigts), puis inverser.
                \item Paume contre paume, doigts entrelacés.
                \item Dos des doigts dans paume opposée (griffes).
                \item Rotation pouce dans paume opposée.
                \item Pulpe des doigts dans paume opposée (rotation).
            \end{itemize}
        \item \textbf{Rinçage :} Eau courante, mains vers le bas (éviter recontamination par l'eau sale coulant vers le coude).
        \item \textbf{Séchage :} Serviette à usage unique (papier) ou air pulsé. \textbf{Ne pas toucher le robinet avec les mains propres} (utiliser le coude ou le papier).
        \item \textbf{Différence clé :} Lavage chirurgical = avant-bras inclus + durée plus longue (\SI{3}{\minute} à \SI{5}{\minute} savon antiseptique OU \SI{1}{\minute}\SI{30}{\second} solution hydroalcoolique sur mains sèches) + brossage ongles (si première fois de la journée).
    \end{enumerate}
    \cle{Règle d'or} : « Mains propres ne touchent que du propre (stérile ou aseptisé) ».
\end{methode}

\begin{exoresolu}[Analyse d'un protocole de lavage des mains au laboratoire]
    \textbf{Énoncé :} Un élève de 3ème réalise le TP N°6 « Pratique de l'asepsie et de l'antisepsie ». Il doit ensemencer une boîte de Pétri gélosée (milieu nutritif stérile) avec ses doigts après différents traitements. Il note ses observations après incubation \SI{24}{\hour} à \SI{37}{\celsius}.
    \begin{center}
    \begin{tabular}{|c|c|c|}\hline
    \rowcolor{popblueL} \textbf{Boîte} & \textbf{Traitement des mains} & \textbf{Aspect du gélose après 24h} \\ \hline
    Témoin (T) & Aucune (mains sales) & Nombreuses colonies variées, conflentes \\ \hline
    A & Lavage rapide eau seule (\SI{5}{\second}) & Colonies nombreuses, peu différentes du témoin \\ \hline
    B & Lavage hygiénique savon \SI{30}{\second} (6 étapes) & Quelques colonies isolées (\textasciitilde 5-10) \\ \hline
    C & Friction solution hydroalcoolique \SI{30}{\second} & Aucune colonie visible \\ \hline
    D & Gants stériles mis sur mains sales & Aucune colonie visible \\ \hline
    \end{tabular}
    \end{center}
    \textbf{Questions :}
    \begin{enumerate}
        \item Interprétez les résultats des boîtes A et B par rapport au témoin T.
        \item Expliquez pourquoi la boîte C ne montre aucune colonie alors que la boîte B en montre quelques-unes.
        \item La boîte D prouve-t-elle que les gants stériles remplacent le lavage des mains ? Justifiez.
        \item Quelle conclusion tirer sur l'efficacité respective de l'asepsie (gants) et de l'antisepsie (lavage/friction) ?
    \end{enumerate}
    \tcblower
    \textbf{Solution détaillée :}
    \begin{enumerate}
        \item \textbf{Boîte A (Eau seule) :} L'eau seule élimine mécaniquement une partie des salissures visibles et de la flore transitoire de surface, mais \textbf{ne tue pas les microbes} (pas d'action antiseptique) et n'élimine pas la flore résidente. Le nombre de colonies reste élevé, proche du témoin. \textbf{Boîte B (Savon \SI{30}{\second}) :} Le savon (tensioactif) décolle les lipides cutanés piégeant les microbes ; la friction mécanique (6 étapes) + rinçage élimine \textbf{efficacement la flore transitoire}. La flore résidente (profonde, non pathogène habituellement) persiste, d'où les quelques colonies résiduelles. C'est une \textbf{antisepsie} incomplète sur la flore résidente.
        \item \textbf{Boîte C (Solution hydroalcoolique - SHA) :} La SHA (ex: gel hydroalcoolique à \SI{60}{\percent}-\SI{70}{\percent} éthanol) a une action \textbf{bactéricide, virucide, fongicide rapide} par dénaturation des protéines. Elle détruit \textbf{à la fois la flore transitoire ET une grande partie de la flore résidente} superficielle en \SI{30}{\second} sur mains \textbf{propres et sèches}. C'est une antisepsie chimique plus puissante que le lavage simple au savon sur la flore résidente.
        \item \textbf{Non.} La boîte D montre que le gant stérile (asepsie = barrière physique) empêche le contact direct main/gélose. \textbf{Mais} si les mains sont sales/sueurs sous le gant, toute micro-perforation (fréquente, invisible) ou contamination lors du retrait du gant (technique « gant dans gant ») relarguerait des millions de germes. Le lavage (antisepsie) \textbf{préalable} est indispensable pour réduire la charge microbienne \textbf{sous} le gant. L'asepsie (gant) et l'antisepsie (lavage) sont \textbf{complémentaires}, non interchangeables.
        \item \textbf{Conclusion :} L'\textbf{asepsie} (gants stériles, boîte D) est la méthode la plus fiable pour \textbf{empêcher tout contact} entre la main (source de germes) et le milieu stérile. L'\textbf{antisepsie} (lavage savon B, SHA C) \textbf{réduit la charge microbienne} sur la peau vivante. La SHA (C) est plus efficace que le savon (B) sur la flore résidente. En pratique clinique (soins, chirurgie), on associe \textbf{antisepsie des mains (lavage/friction) PUIS asepsie (gants stériles)} pour une sécurité maximale.
    \end{enumerate}
\end{exoresolu}

\begin{methode}[Choisir et appliquer un antiseptique adapté (Antisepsie)]
    \textbf{Objectif :} Sélectionner le bon produit selon la zone (peau saine, muqueuse, plaie) et respecter le protocole d'application (temps de contact).
    \begin{enumerate}
        \item \textbf{Identifier la zone à traiter :}
            \begin{itemize}
                \item \textbf{Peau saine intacte} (pré-opératoire, injection) : Produits à large spectre, rémanence (action prolongée). Ex: \textbf{Povidone iodée (Bétadine dermique)}, \textbf{Chlorhexidine alcoolique}.
                \item \textbf{Muqueuses / Plaies ouvertes / Brûlures} : Produits non irritants, non cytotoxiques pour les cellules de la cicatrisation. Ex: \textbf{Chlorhexidine aqueuse (\SI{0,05}{\percent}-\SI{0,1}{\percent})}, \textbf{Dakin (hypochlorite Na \SI{0,1}{\percent}-\SI{0,5}{\percent})}, \textbf{Eau oxygénée (\ce{H2O2} - action mécanique moussante)}. \textbf{Interdit :} Alcool (brûlure, nécrose), Bétadine dermique (alcoolisée) sur muqueuse/plaie profonde (risque absorption iode thyroïdien).
            \end{itemize}
        \item \textbf{Vérifier la DLUO (Date Limite d'Utilisation Optimale) et l'aspect} (pas de trouble, pas de dépôt).
        \item \textbf{Nettoyer préalablement} la zone (eau + savon doux, sérum physiologique pour plaie) pour éliminer saletés, sang, pus qui inactivent l'antiseptique (charge organique).
        \item \textbf{Appliquer} par tamponnage (compresse stérile) du centre vers la périphérie (éviter récontamination du centre). Ne pas frotter (plaie).
        \item \textbf{Respecter le temps de contact} (indiqué sur notice : ex \SI{1}{\minute} à \SI{2}{\minute} pour Bétadine, \SI{1}{\minute} pour Chlorhexidine) \textbf{AVANT} tout geste invasif ou pansement. Laisser sécher à l'air.
        \item \textbf{Rincer si nécessaire} (Dakin, Eau oxygénée sur plaie) ou laisser agir (Bétadine, Chlorhexidine sur peau saine).
    \end{enumerate}
    \cle{Antisepsie} $\neq$ Stérilisation : L'antiseptique ne tue pas les spores, n'agit pas en profondeur dans les tissus nécrosés (nécessite débridement chirurgical).
\end{methode}

\begin{exoresolu}[Cas cliniques : Choix de l'antiseptique au dispensaire]
    \textbf{Énoncé :} Au centre de santé de district de Dschang, l'infirmier doit préparer trois patients. Pour chaque situation, indiquez l'antiseptique le plus adapté parmi la liste (Bétadine dermique, Bétadine scrub, Chlorhexidine aqueuse \SI{0,05}{\percent}, Dakin \SI{0,1}{\percent}, Alcool \SI{70}{\percent}, Eau oxygénée 10 vol) et justifiez le choix en une phrase.
    \begin{enumerate}
        \item Patient A : Préparation du champ opératoire pour une appendicectomie (peau saine du ventre).
        \item Patient B : Désinfection d'une plaie lacérée sale au genou (terre, graviers) avant sutures.
        \item Patient C : Désinfection du méat urinaire avant pose de sonde vésicale (muqueuse).
        \item Patient D : Désinfection de la peau saine pour une prise de sang veineuse (pli du coude).
    \end{enumerate}
    \tcblower
    \textbf{Solution détaillée :}
    \begin{enumerate}
        \item \textbf{Patient A (Peau saine, chirurgie) :} \textbf{Bétadine dermique (Povidone iodée alcoolique) OU Chlorhexidine alcoolique \SI{0,5}{\percent}-\SI{2}{\percent}.} \textit{Justification :} Large spectre (bactéries, virus, champignons), action rémanente (persistance) grâce à l'iode libre ou la chlorhexidine fixée sur kératine, potentialisée par l'alcool pour une action immédiate. Idéal pour champ opératoire.
        \item \textbf{Patient B (Plaie sale, tissus lésés) :} \textbf{Eau oxygénée 10 vol (premier lavage mécanique) PUIS Dakin \SI{0,1}{\percent} OU Chlorhexidine aqueuse \SI{0,05}{\percent}.} \textit{Justification :} L'eau oxygénée décolle les corps étrangers (effet moussant) ; le Dakin ou la Chlorhexidine aqueuse sont peu irritants, non cytotoxiques pour les fibroblastes (cicatrisation), actifs en présence de matière organique (sang, pus). \textbf{Interdit :} Bétadine dermique (alcool = douleur/nécrose, iode = absorption systémique sur grande surface lésée), Alcool (douleur, nécrose).
        \item \textbf{Patient C (Muqueuse urétrale) :} \textbf{Chlorhexidine aqueuse \SI{0,05}{\percent} (ou Bétadine solution aqueuse \SI{10}{\percent} si pas d'allergie iode).} \textit{Justification :} Seule la forme \textbf{aqueuse (sans alcool)} est tolérée sur muqueuse (non irritante, non cytotoxique). L'alcool et la Bétadine dermique (alcoolisée) sont contre-indiqués (brûlure chimique sévère).
        \item \textbf{Patient D (Peau saine, geste simple) :} \textbf{Alcool \SI{70}{\percent} (ou Chlorhexidine alcoolique, Bétadine dermique).} \textit{Justification :} Action rapide (quelques secondes), large spectre, coût faible, évaporation rapide sans rinçage. Suffisant pour une ponction veineuse (risque infectieux faible vs chirurgie). Temps de contact : laisser sécher (\SI{30}{\second}).
    \end{enumerate}
    \textbf{Rappel BEPC :} Sur peau saine $\rightarrow$ Antiseptique alcoolique (rapide, rémanent). Sur plaie/muqueuse $\rightarrow$ Antiseptique aqueux (non irritant, non cytotoxique).
\end{exoresolu}

\begin{methode}[Utiliser correctement le préservatif masculin (Prophylaxie IST/VIH/Sida)]
    \textbf{Objectif :} Appliquer la procédure exacte pour garantir l'efficacité contraceptive et la prévention du VIH/IST (efficacité \textasciitilde \SI{98}{\percent} en usage parfait).
    \begin{enumerate}
        \item \textbf{Vérifier avant ouverture :} Date de péremption (non dépassée), marque CE/NF/ISO, emballage intact (pas de déchirure, poche d'air présente = test du « coussin »).
        \item \textbf{Ouvrir avec précaution :} Déchirer le long du bord prédécoupé \textbf{avec les doigts} (jamais ciseaux, dents, ongles, bijoux). Risque de micro-déchirure invisible.
        \item \textbf{Sens de déroulement :} Placer le préservatif sur le gland \textbf{en érection complète} (avant tout contact génital). Vérifier le sens : le réservoir (bout) doit être à l'extérieur, le bord doit se dérouler vers le bas (comme un sombrero). Si mis à l'envers $\rightarrow$ jeter et prendre un neuf (contact pré-éjaculatoire = risque).
        \item \textbf{Pincer le réservoir :} \textbf{Indispensable.} Pincer le bout entre pouce et index pour chasser l'air (évite l'éclatement par surpression lors de l'éjaculation).
        \item \textbf{Dérouler jusqu'à la base :} De l'autre main, dérouler sur toute la longueur du pénis en érection. Vérifier qu'il tient bien.
        \item \textbf{Pendant le rapport :} Vérifier occasionnellement qu'il ne glisse pas ou ne se déchire pas. Lubrifiant \textbf{uniquement à base d'eau ou de silicone} (jamais gras : vaseline, huile, beurre, crème $\rightarrow$ rupture latex en \SI{1}{\minute}).
        \item \textbf{Retrait immédiat après éjaculation :} \textbf{Avant la perte d'érection.} Tenir le préservatif à la base du pénis pour éviter qu'il ne glisse et ne reste dans le vagin/anus. Retirer délicatement.
        \item \textbf{Élimination :} Nœud au bout, jeter à la poubelle (pas dans les toilettes = bouchage). \textbf{Jamais de réutilisation.}
    \end{enumerate}
    \cle{Préservatif} = seule méthode \textbf{double protection} (Contraception + Prévention IST/VIH). \cle{Pré-éjaculat} = peut contenir spermatozoïdes et VIH $\rightarrow$ mettre préservatif \textbf{AVANT} tout contact.
\end{methode}

\begin{exoresolu}[Analyse d'erreurs fréquentes d'utilisation du préservatif]
    \textbf{Énoncé :} Lors d'une séance d'éducation pour la santé au lycée de Ngaoundéré, l'infirmière scolaire présente 5 scénarios d'utilisation du préservatif masculin. Pour chaque scénario, dites si l'utilisation est \textbf{correcte (C)} ou \textbf{incorrecte (I)}. Si incorrecte, précisez l'erreur et le risque encouru.
    \begin{enumerate}
        \item Paul vérifie la date (valide), ouvre l'emballage avec ses dents car ses mains sont grasses.
        \item Marie et son partenaire mettent le préservatif après la pénétration, juste avant l'éjaculation.
        \item Jean pince le réservoir, déroule le préservatif jusqu'à la base, utilise de la vaseline comme lubrifiant.
        \item Après l'éjaculation, Ahmed attend que son érection disparaisse complètement avant de retirer le préservatif, sans le tenir à la base.
        \item Fatou (partenaire) vérifie l'emballage, l'ouvre avec les doigts, place le préservatif sur le pénis en érection de son partenaire en pinçant le réservoir, le déroule, utilise un gel lubrifiant à base d'eau. Après l'éjaculation, le partenaire tient la base et retire le préservatif avant la fin de l'érection.
    \end{enumerate}
    \tcblower
    \textbf{Solution détaillée :}
    \begin{enumerate}
        \item \textbf{I (Incorrect).} \textbf{Erreur :} Ouverture avec les dents. \textbf{Risque :} Micro-déchirure du latex invisible à l'œil nu $\rightarrow$ rupture pendant le rapport $\rightarrow$ grossesse non désirée + transmission VIH/IST. \textbf{Correct :} Ouvrir avec les doigts le long de la languette.
        \item \textbf{I (Incorrect).} \textbf{Erreur :} Mise en place tardive (après pénétration). \textbf{Risque :} Contact du pré-éjaculat (liquide séminal) avec la muqueuse vaginale. Le pré-éjaculat peut contenir des spermatozoïdes (fuite antérieure) et \textbf{contient le VIH} (charge virale). $\rightarrow$ Risque grossesse + contamination VIH/IST.
        \item \textbf{I (Incorrect).} \textbf{Erreur :} Utilisation de vaseline (corps gras). \textbf{Risque :} Dégradation chimique du latex (perte d'élasticité) en \textbf{moins d'une minute} $\rightarrow$ rupture quasi certaine. \textbf{Correct :} Lubrifiant à base d'eau ou silicone (ex: K-Y, Durex Play, gel médical).
        \item \textbf{I (Incorrect).} \textbf{Erreur :} Retrait tardif (perte d'érection) sans tenir la base. \textbf{Risque :} Le pénis rétrécit, le préservatif devient lâche $\rightarrow$ glisse et reste dans le vagin/anus (rétention) $\rightarrow$ fuite de sperme $\rightarrow$ grossesse + IST. \textbf{Correct :} Retrait \textbf{immédiat} après éjaculation, \textbf{en tenant fermement la base} du préservatif.
        \item \textbf{C (Correct).} \textbf{Analyse :} Toutes les étapes sont respectées : vérification emballage/date, ouverture doigts, mise avant contact (érection), pincement réservoir, déroulement complet, lubrifiant aqueux, retrait immédiat en tenant la base. C'est l'utilisation \textbf{parfaite (usage parfait)}.
    \end{enumerate}
    \textbf{Conclusion :} L'efficacité réelle (usage courant \textasciitilde \SI{85}{\percent}) est inférieure à l'efficacité théorique (usage parfait \textasciitilde \SI{98}{\percent}) à cause de ces erreurs humaines. L'éducation aux gestes précis (Méthode ci-dessus) est cruciale.
\end{exoresolu}

\begin{methode}[Rédiger une argumentation scientifique pour une mesure de prévention (Savoir-faire : Rédiger et justifier)]
    \textbf{Objectif :} Structurer une réponse écrite (type BEPC) justifiant une mesure de santé publique (campagne dépistage, distribution préservatifs, vaccination, hygiène).
    \begin{enumerate}
        \item \textbf{Analyser la situation / la famille de situations :} Identifier le problème de santé (ex: recrudescence VIH chez les 15-24 ans, épidémie choléra, tétanos néonatal).
        \item \textbf{Mobiliser les savoirs scientifiques (Le « Pourquoi » biologique) :}
            \begin{itemize}
                \item Mode de transmission (sexuel, sanguin, materno-fœtal, féco-orale, respiratoire).
                \item Mécanisme physiopathologique (ex: VIH détruit lymphocytes T CD4+ $\rightarrow$ immunodépression).
                \item Facteurs de risque (multiplicité partenaires, non port préservatif, eau non traitée, non vaccination).
            \end{itemize}
        \item \textbf{Identifier la mesure de prévention proposée (Le « Quoi ») :} Primaire (éviter contamination : préservatif, vaccin, eau potable, éducation), Secondaire (dépistage précoce : test VIH, frottis), Tertiaire (limiter séquelles : ARV, rééducation).
        \item \textbf{Expliquer le mécanisme d'action de la mesure (Le « Comment ») :}
            \begin{itemize}
                \item Barrière physique (préservatif, moustiquaire, gant).
                \item Immunité active (vaccin $\rightarrow$ Ac mémoire).
                \item Rupture chaîne transmission (traitement ARV $\rightarrow$ charge virale indétectable = I=I \textit{Indétectable = Intransmissible}, assainissement).
            \end{itemize}
        \item \textbf{Conclure par l'impact attendu (Résultat) :} Réduction incidence/prévalence, mortalité, coûts sanitaires. Utiliser des connecteurs logiques : \textbf{Donc, par conséquent, ce qui permet, en effet, or, cependant.}
    \end{enumerate}
    \cle{Structure type BEPC} : \textbf{Constat (Données) $\rightarrow$ Explication biologique (Savoirs) $\rightarrow$ Mesure $\rightarrow$ Mécanisme d'efficacité $\rightarrow$ Conclusion/Impact.}
\end{methode}

\begin{exoresolu}[Rédaction : Justifier la distribution gratuite de préservatifs aux lycéens]
    \textbf{Énoncé :} Le Ministère de la Santé Publique (MINSANTÉ) et le Ministère de l'Éducation Secondaire (MINESEC) lancent une campagne de distribution gratuite de préservatifs masculins et féminins dans les lycées du Cameroun, accompagnée de séances d'éducation à la vie sexuelle. En vous appuyant sur vos connaissances sur le VIH/Sida et les IST, rédigez un argumentaire scientifique de 15 à 20 lignes pour justifier cette mesure auprès des parents d'élèves réticents.
    \tcblower
    \textbf{Solution détaillée (Modèle de rédaction attendue) :}
    \textbf{Argumentaire scientifique pour la distribution de préservatifs en milieu scolaire}

    \textbf{1. Constat épidémiologique alarmant :} Au Cameroun, la prévalence du VIH reste élevée chez les jeunes de 15-24 ans (environ \SI{1,5}{\percent} à \SI{2}{\percent} selon les enquêtes EDS-MICS récentes), avec une incidence plus marquée chez les filles. Les IST (chlamydiose, gonococcie, syphilis) sont très fréquentes et souvent asymptomatiques, favorisant la transmission du VIH (lésions génitales = portes d'entrée).

    \textbf{2. Rappel physiopathologique :} Le VIH se transmet par les liquides biologiques contaminés (sperme, sécrétions vaginales, sang, lait maternel) lors de rapports sexuels non protégés. Le virus cible les lymphocytes T CD4+, piliers de l'immunité adaptative. Leur destruction progressive conduit au Sida (immunodépression sévère), fatal sans traitement ARV à vie. \textbf{Il n'existe ni vaccin préventif, ni traitement curatif éradiquant le virus.}

    \textbf{3. Efficacité prouvée du préservatif :} Le préservatif (masculin en latex/polyisoprène, féminin en nitrile) est la \textbf{seule méthode offrant une double protection} : contraceptive (barrière aux spermatozoïdes) et prophylaxique (barrière étanche au VIH, virus hépatites B/C, bactéries IST). Son efficacité en usage parfait atteint \SI{98}{\percent} pour le VIH. La distribution gratuite lève l'obstacle financier (pouvoir d'achat élèves) et la disponibilité immédiate évite le « rapport à risque par défaut de préservatif ».

    \textbf{4. Nécessité de l'éducation associée (Savoir-être/Savoir-faire) :} La distribution seule est insuffisante. Les séances d'éducation permettent d'enseigner la \textbf{méthode correcte d'utilisation} (vérification date, pincement réservoir, lubrifiant aqueux, retrait immédiat), de déconstruire les idées fausses (« ça diminue le plaisir », « c'est pour les prostituées ») et de promouvoir le \textbf{consentement} et le dépistage volontaire.

    \textbf{5. Conclusion :} Cette mesure de \textbf{prévention primaire} combine une barrière physique efficace (préservatif) et l'acquisition de compétences psychosociales (négociation, usage correct). Elle vise à réduire l'incidence du VIH et des IST chez les adolescents, à prévenir les grossesses précoces (cause majeure de décrochage scolaire des filles) et à former une génération responsable. C'est un investissement en santé publique et en capital humain, conforme aux ODD (Objectifs de Développement Durable) 3 et 4.
\end{exoresolu}

\begin{methode}[Interpréter les résultats du TP N°6 (Boîtes de Pétri : Asepsie vs Antisepsie)]
    \textbf{Objectif :} Savoir lire une boîte de Pétri gélosée après incubation et en déduire l'efficacité d'une technique aseptique ou antiseptique.
    \begin{enumerate}
        \item \textbf{Observer la boîte témoin (Témoin positif) :} Doit montrer une croissance abondante (confluent ou nombreuses colonies variées). \textit{Valide que le milieu est fertile et l'inoculum contaminant.}
        \item \textbf{Observer la boîte stérile (Témoin négatif / Contrôle stérilité) :} Milieu non inoculé, incubé. Doit rester \textbf{stérile (aucune colonie)}. \textit{Valide la stérilisation du milieu et l'asepsie de la manipulation (pas de contamination aérienne/manipulateur).}
        \item \textbf{Compter/Estimer les UFC (Unités Formant Colonies) :} « Nombreuses (>100) », « Quelques-unes (10-50) », « Rares (<10) », « Aucune (0) ». Noter la morphologie (couleur, aspect, hémolyse si gélose sang).
        \item \textbf{Comparer les conditions testées :}
            \begin{itemize}
                \item \textbf{Test Asepsie (ex: Gants stériles, Champ stérile, Matériel stérilisé) :} Résultat attendu = \textbf{0 colonie} (stérilité maintenue). Toute colonie = rupture d'asepsie (gant percé, asepsie défaillante).
                \item \textbf{Test Antisepsie (ex: Lavage mains, SHA, Bétadine, Alcool) :} Résultat = \textbf{Réduction significative} du nombre de colonies par rapport au témoin « mains sales ». \textbf{Jamais 0 colonie} sur peau vivante (flore résidente persistante, sauf SHA prolongée). Comparer « Lavage eau » vs « Lavage savon » vs « SHA ».
            \end{itemize}
        \item \textbf{Conclure scientifiquement :} « La technique X (asepsie) a maintenu la stérilité (0 UFC) / La technique Y (antisepsie) a réduit la charge microbienne de Z\% mais n'a pas stérilisé la peau. »
    \end{enumerate}
    \cle{UFC} : Unité Formant Colonie = une colonie visible = origine d'un seul germe viable (ou amas) déposé sur le gélose.
\end{methode}

\begin{exoresolu}[Interprétation de photographies de boîtes de Pétri (TP N°6)]
    \textbf{Énoncé :} Lors du TP N°6, quatre boîtes de Pétri (gélose Chapman, sélective \textit{Staphylococcus}) ont été inoculées par empreinte digitale (doigt pulpe) après différents traitements. Incubation \SI{24}{\hour} à \SI{37}{\celsius}. Les photos ci-dessous (schématisées) montrent les résultats.
    \begin{popfigure}
    \begin{tikzpicture}[scale=0.7, every node/.style={transform shape}]
        % Style boite petri
        \def\petri#1#2#3#4{ % x, y, label, colonies
            \begin{scope}[shift={(#1,#2)}]
                \draw[popdark, thick] (0,0) circle (2.5);
                \draw[popdark, thick] (0,0) circle (2.2); % bord interne
                \node[popdark, font=\bfseries\small] at (0, -3) {#3};
                % Colonies
                \foreach \c in {#4} {
                    \fill[poporange!80!popdark] \c circle (0.12);
                }
            \end{scope}
        }
        % Boîte 1 : Témoin Mains Sales (Confluent)
        \petri{0}{0}{Boîte 1 : Témoin (Mains sales)}{
            (0.2,0.5), (-0.5,-0.3), (1.2,0.8), (-1.5,0.2), (0.8,-1.2), (-0.8,-1.5), (1.8,0.1), (-1.2,1.5), (0,1.8), (1.5,-0.8),
            (-0.2,0.1), (0.5,1.2), (-1.8,-0.5), (0.3,-0.8), (-1.0,0.8), (1.0,1.0), (-0.5,1.5), (1.5,1.2), (-1.5,-1.0), (0.8,1.5)
        }
        % Boîte 2 : Lavage Eau seule
        \petri{6}{0}{Boîte 2 : Lavage Eau 10s}{
            (0.2,0.5), (-0.5,-0.3), (1.2,0.8), (-1.5,0.2), (0.8,-1.2), (-0.8,-1.5), (1.8,0.1), (-1.2,1.5), (0,1.8), (1.5,-0.8),
            (-0.2,0.1), (0.5,1.2), (-1.8,-0.5), (0.3,-0.8), (-1.0,0.8)
        }
        % Boîte 3 : Lavage Savon 30s (6 etapes)
        \petri{12}{0}{Boîte 3 : Lavage Savon 30s}{
            (0.5,0.2), (-1.0,-0.5), (1.5,1.0), (-0.5,1.5)
        }
        % Boîte 4 : Friction SHA 30s
        \petri{18}{0}{Boîte 4 : Friction SHA 30s}{
            % Aucune colonie
        }
    \end{tikzpicture}
    \legende{Résultats schématisés du TP N°6 sur gélose Chapman (sélective Staphylocoques) après incubation 24h à 37°C. Chaque point orange représente une colonie (UFC).}
    \end{popfigure}
    \textbf{Questions :}
    \begin{enumerate}
        \item Décrivez l'aspect des cultures sur les boîtes 1 et 2. Que conclure sur l'efficacité du lavage à l'eau seule ?
        \item Comparez les boîtes 3 et 4. Quelle technique antiseptique est la plus efficace sur la flore cutanée ? Pourquoi ?
        \item La boîte 4 est-elle stérile ? Cela signifie-t-il que la main est stérile ? Justifiez.
        \item Si l'on avait mis un gant stérile sur la main sale (sans lavage) et fait l'empreinte gantée, quel résultat attendriez-vous ? (Relier à la méthode Asepsie).
    \end{enumerate}
    \tcblower
    \textbf{Solution détaillée :}
    \begin{enumerate}
        \item \textbf{Boîte 1 (Témoin) :} Culture \textbf{confluente} (colonies très nombreuses, jointives, recouvrant tout le gélose). Traduction : charge microbienne massive (flore transitoire + résidente) sur mains sales. \textbf{Boîte 2 (Eau seule) :} Culture \textbf{très abondante, quasi identique au témoin} (légèrement moins dense peut-être). \textbf{Conclusion :} Le lavage à l'eau seule (\SI{10}{\second}, sans savon, sans friction mécanique rigoureuse) \textbf{n'est pas une antisepsie efficace}. Il élimine les salissures macroscopiques mais laisse la quasi-totalité de la flore microbienne (transitoire et résidente).
        \item \textbf{Boîte 3 (Savon \SI{30}{\second}) :} \textbf{Quelques colonies isolées} (4-5 UFC). \textbf{Boîte 4 (SHA \SI{30}{\second}) :} \textbf{Aucune colonie visible} (0 UFC). \textbf{Comparaison :} La \textbf{friction hydroalcoolique (SHA) est plus efficace} que le lavage au savon simple. \textbf{Explication :} Le savon (tensioactif) + friction mécanique élimine \textbf{mécaniquement} la flore transitoire (décollement) mais laisse la flore résidente (profonde, adhérente). La SHA (alcool \SI{60}{\percent}-\SI{70}{\percent}) a une action \textbf{chimique bactéricide/fongicide/virucide} qui détruit \textbf{les deux flores} (transitoire et superficielle de la résidente) par dénaturation protéique, en \SI{30}{\second} sur mains propres/sèches.
        \item \textbf{Non, la main n'est pas stérile.} \textbf{0 UFC sur gélose Chapman} signifie : « Aucune bactérie viable capable de croître sur ce milieu sélectif (Staphylocoques) n'a été déposée ». \textbf{Raisons :} (1) La flore résidente profonde (follicules pileux, glandes sébacées) n'a pas été atteinte par la SHA (action de surface). Elle remontera en 1-2h (recolonisation). (2) La gélose Chapman est sélective pour \textit{Staphylococcus} ; d'autres germes (Corynebactéries, champignons) pourraient ne pas pousser. (3) Limite de détection : < 1 UFC/cm². \textbf{Antisepsie $\neq$ Stérilisation.}
        \item \textbf{Résultat attendu : Boîte stérile (0 UFC).} \textbf{Justification :} Le gant stérile est une barrière physique (\textbf{Asepsie}) qui empêche \textbf{tout contact} entre la peau (source de germes) et le milieu de culture. Même si la main est très sale (Boîte 1), le gant stérile intact garantit l'asepsie du geste. \textbf{Condition :} Gant intact (pas de micro-perforation), mise en place aseptique (mains sèches), pas de contamination du gant par l'environnement lors de la pose. Cela illustre la supériorité de l'asepsie (barrière) sur l'antisepsie (réduction) pour garantir la stérilité d'un champ opératoire.
    \end{enumerate}
\end{exoresolu}

\begin{attention}[Les 5 erreurs qui coûtent des points au BEPC (SVTEEHB 3ème)]
    \begin{enumerate}
        \item \textbf{Confondre Asepsie et Antisepsie (Erreur n°1) :} Écrire « L'asepsie tue les microbes » ou « L'antisepsie empêche l'arrivée des microbes ». \textbf{Sanction :} 0 point à la définition. \textbf{Rappel :} Asepsie = \textbf{Empêcher l'arrivée} (support inerte, barrière, stérilisation). Antisepsie = \textbf{Détruire/Inhiber sur vivant} (peau, muqueuse, plaie).
        \item \textbf{Dire que le lavage des mains « stérilise » les mains (Erreur n°2) :} Aucun lavage (même chirurgical) ne stérilise la peau vivante (flore résidente profonde). On parle de \textbf{« décontamination »} ou \textbf{« antisepsie »}. Seule la stérilisation (autoclave, Poupinel, filtration) rend un matériel \textbf{stérile} (absence de tout germe viable, spores comprises).
        \item \textbf{Oublier le pincement du réservoir ou le retrait « base tenue » pour le préservatif (Erreur n°3) :} Dans une description de procédure, l'oubli de \textbf{pincer le réservoir} (risque éclatement) ou de \textbf{tenir la base au retrait} (risque glissement/fuite) fait perdre les points de « savoir-faire technique ». Ces deux gestes sont \textbf{obligatoires} au programme.
        \item \textbf{Proposer de l'alcool ou de la Bétadine dermique sur une plaie ou une muqueuse (Erreur n°4) :} Confondre « peau saine » et « plaie/muqueuse ». L'alcool et l'iode alcoolisé (Bétadine dermique) sont \textbf{contre-indiqués} (douleur, nécrose tissulaire, retard cicatrisation, absorption iode). Sur plaie/muqueuse : \textbf{Chlorhexidine aqueuse, Dakin, Eau oxygénée}.
        \item \textbf{Rédiger une justification sans mécanisme biologique (Erreur n°5) :} Réponse du type « C'est bien car ça protège ». \textbf{Zéro point.} Au BEPC, « Justifier » = \textbf{Mobiliser un savoir} (ex: « Le préservatif est une barrière physique étanche au virus VIH présent dans le sperme, empêchant sa rencontre avec les lymphocytes T CD4+ de la muqueuse »). Toujours faire le lien : \textbf{Mesure $\rightarrow$ Mécanisme biologique $\rightarrow$ Efficacité.}
    \end{enumerate}
\end{attention}

\newpage

\coursec{Exercices}

\courssub{Je vérifie mes connaissances}

\begin{exercice}[QCM : asepsie ou antisepsie ?]
Pour chaque affirmation, choisir la (ou les) bonne(s) réponse(s) en justifiant brièvement.
\begin{enumerate}
\item L'\cle{asepsie} consiste à :
\begin{enumerate}
\item détruire les microbes déjà présents sur la peau ;
\item empêcher l'arrivée des microbes dans un milieu ;
\item tuer les microbes à l'intérieur de l'organisme.
\end{enumerate}
\item La \cle{antisepsie} consiste à :
\begin{enumerate}
\item empêcher la pénétration des microbes ;
\item détruire les microbes présents sur la peau ou sur une plaie ;
\item stériliser le matériel chirurgical.
\end{enumerate}
\item Un antiseptique s'applique :
\begin{enumerate}
\item sur les tissus vivants (peau, plaie) ;
\item uniquement sur les objets ;
\item uniquement par voie buccale.
\end{enumerate}
\item Le préservatif est une pratique :
\begin{enumerate}
\item d'antisepsie ;
\item d'asepsie ;
\item de vaccination.
\end{enumerate}
\end{enumerate}
\end{exercice}

\begin{exercice}[Vrai ou faux ?]
Répondre par VRAI ou FAUX, puis justifier chaque réponse par une phrase.
\begin{enumerate}
\item L'asepsie détruit les microbes déjà présents sur une plaie.
\item Se laver les mains à l'eau et au savon avant de soigner une plaie est une pratique d'asepsie.
\item La bétadine est un antiseptique.
\item Le préservatif, correctement utilisé, protège à la fois des infections sexuellement transmissibles et des grossesses non désirées.
\item Un objet stérilisé ne contient aucun microbe vivant.
\end{enumerate}
\end{exercice}

\begin{exercice}[Définitions]
Définir en une ou deux phrases chacun des termes suivants, en donnant pour chacun un exemple de la vie courante : \cle{contamination}, \cle{infection}, \cle{asepsie}, \cle{antisepsie}, \cle{antiseptique}.
\end{exercice}

\begin{exercice}[Classer des pratiques]
Classer les pratiques suivantes en deux colonnes « asepsie » ou « antisepsie », en justifiant chaque classement :
\begin{enumerate}
\item faire bouillir l'eau de boisson ;
\item appliquer de la bétadine sur une plaie ;
\item porter des gants propres pour faire un pansement ;
\item utiliser un préservatif ;
\item bien cuire la viande avant de la manger ;
\item nettoyer la peau à l'alcool avant une piqûre.
\end{enumerate}
\end{exercice}

\courssub{Je m'entraîne}

\begin{exercice}[Le lavage des mains]
Au dispensaire du quartier, l'infirmier se lave soigneusement les mains à l'eau et au savon avant chaque soin, puis enfile des gants propres.
\begin{enumerate}
\item À quelle pratique (asepsie ou antisepsie) correspondent ces gestes ? Justifier.
\item Expliquer pourquoi ces gestes réduisent le risque de contamination des malades.
\item Citer deux autres gestes d'asepsie observés dans un centre de santé.
\end{enumerate}
\end{exercice}

\begin{exercice}[Compte rendu de TP : mains lavées et mains non lavées]
Au cours du TP N°6, un élève réalise l'expérience suivante : il appuie les doigts de sa main \emph{non lavée} sur la gélose d'une boîte de Pétri (boîte A), puis, après lavage soigneux des mains au savon, il appuie les doigts de sa main \emph{lavée} sur la gélose d'une seconde boîte (boîte B). Les deux boîtes sont placées à l'étuve à \SI{37}{\celsius} pendant 48 heures. Résultats : nombreuses colonies de microbes dans la boîte A ; très peu de colonies dans la boîte B.
\begin{enumerate}
\item Reproduire et légender le schéma des deux boîtes de Pétri après 48 heures.
\item Formuler l'hypothèse que cette expérience permet de tester.
\item Rédiger l'interprétation des résultats obtenus.
\item Rédiger une conclusion indiquant la pratique (asepsie ou antisepsie) mise en évidence.
\end{enumerate}
\end{exercice}

\begin{exercice}[Une blessure au genou]
En jouant au football, Junior tombe et son genou est éraflé par un clou rouillé. La plaie saigne légèrement et est souillée de terre.
\begin{enumerate}
\item Décrire, étape par étape, la conduite à tenir pour soigner cette plaie.
\item Pour chaque étape, préciser s'il s'agit d'une pratique d'asepsie ou d'antisepsie, en justifiant.
\item Expliquer pourquoi il est dangereux de laisser la plaie souillée sans soin.
\end{enumerate}
\end{exercice}

\begin{exercice}[L'eau de boisson au village]
Dans un village, l'eau du puits contient des microbes responsables de diarrhées. On conseille aux familles de faire bouillir l'eau avant de la boire et de la conserver dans des récipients propres et fermés.
\begin{enumerate}
\item Identifier, parmi ces conseils, la pratique qui relève de l'antisepsie et celle qui relève de l'asepsie. Justifier.
\item Expliquer pourquoi la conservation dans un récipient fermé est indispensable après ébullition.
\end{enumerate}
\end{exercice}

\courssub{Je me prépare au Bac}

\begin{exercice}[Type BEPC : l'expérience de Semmelweis]
En 1847, dans une maternité de Vienne (Autriche), le médecin Ignace Semmelweis constate que de nombreuses femmes meurent après l'accouchement d'une infection appelée « fièvre puerpérale ». Il impose alors au personnel soignant le lavage des mains avec une solution antiseptique avant chaque accouchement. Le tableau ci-dessous donne les résultats obtenus.

\begin{center}
\begin{tabularx}{\linewidth}{|l|X|X|}
\hline
\rowcolor{popblueL} \textbf{Période} & \textbf{Avant le lavage des mains} & \textbf{Après le lavage des mains} \\
\hline
Nombre d'accouchements & \num{4000} & \num{3500} \\
\hline
Nombre de décès par fièvre puerpérale & 400 & 35 \\
\hline
\end{tabularx}
\end{center}

\begin{enumerate}
\item Calculer le taux de mortalité (en \%) pour chaque période. On rappelle : taux de mortalité $= \dfrac{\text{nombre de décès}}{\text{nombre d'accouchements}} \times 100$.
\item Calculer la variation (en points de pourcentage) du taux de mortalité entre les deux périodes.
\item À quelle pratique (asepsie ou antisepsie) correspond le lavage des mains avec une solution antiseptique ? Justifier.
\item Rédiger une conclusion de 3 à 5 lignes sur l'efficacité du lavage des mains, en exploitant les résultats calculés.
\end{enumerate}
\end{exercice}

\begin{exercice}[Type BEPC : situation-problème au centre de santé]
Au centre de santé de Nkolbisson, une infirmière reçoit trois patients le même jour : un enfant blessé au pied par une bouteille cassée, une femme enceinte venue pour une consultation, et un adolescent venu se renseigner sur la prévention des infections sexuellement transmissibles (IST).
\begin{enumerate}
\item Pour l'enfant blessé : décrire les gestes de soins à réaliser, en distinguant clairement les gestes d'asepsie et les gestes d'antisepsie, et en justifiant chacun d'eux.
\item Expliquer pourquoi l'infirmière doit se laver les mains et changer de gants entre deux patients.
\item Pour l'adolescent : expliquer le double rôle du préservatif dans la prévention des IST et des grossesses précoces.
\item Rédiger un court message de sensibilisation (5 à 8 lignes) destiné aux élèves d'un lycée, expliquant comment éviter la contamination par les IST.
\end{enumerate}
\end{exercice}

\begin{exercice}[Type BEPC : étude documentaire sur les IST au Cameroun]
\textbf{Document.} Les infections sexuellement transmissibles (IST), comme la blennorragie, la syphilis ou l'infection à VIH, se transmettent lors de rapports sexuels non protégés avec un partenaire contaminé. Au Cameroun, les campagnes de sensibilisation recommandent l'utilisation systématique du préservatif, qui constitue une barrière empêchant le passage des microbes d'un partenaire à l'autre. Certaines IST non soignées peuvent entraîner de graves complications, comme la stérilité.
\begin{enumerate}
\item Définir les termes : contamination, infection.
\item Expliquer, à partir du document, comment le préservatif empêche la contamination lors d'un rapport sexuel. À quelle pratique (asepsie ou antisepsie) son utilisation correspond-elle ? Justifier.
\item Citer deux IST mentionnées dans le document et une complication possible d'une IST non soignée.
\item Le préservatif protège-t-il aussi des grossesses non désirées ? Justifier.
\item Rédiger une conclusion dégagant l'importance des pratiques de prévention (asepsie, antisepsie, préservatif) pour la santé de chacun.
\end{enumerate}
\end{exercice}

\newpage

\coursec{Corrigés des exercices}

\begin{corrige}[Exercice 1 — QCM : asepsie ou antisepsie ?]
\begin{enumerate}
\item \textbf{Réponse b.} L'\cle{asepsie} est l'ensemble des mesures qui \textbf{empêchent l'arrivée des microbes} dans un milieu ou sur un tissu. Détruire les microbes déjà présents (réponse a) relève de l'antisepsie ; tuer les microbes à l'intérieur de l'organisme (réponse c) relève des médicaments (antibiotiques), ce qui n'est ni l'asepsie ni l'antisepsie.
\item \textbf{Réponse b.} L'\cle{antisepsie} consiste à \textbf{détruire les microbes présents} sur la peau ou sur une plaie, au moyen d'un antiseptique. Empêcher la pénétration des microbes (réponse a) est l'asepsie ; stériliser le matériel (réponse c) est une opération réalisée sur des objets, pas sur les tissus vivants.
\item \textbf{Réponse a.} Un antiseptique s'applique sur les \textbf{tissus vivants} (peau, plaie, muqueuse). Un produit utilisé sur les objets est un désinfectant ; un antiseptique ne s'avale pas.
\item \textbf{Réponse b.} Le préservatif est une \textbf{barrière} qui empêche le passage des microbes d'un partenaire à l'autre : il empêche l'arrivée des microbes, c'est donc une pratique d'\textbf{asepsie} (il ne détruit aucun microbe).
\end{enumerate}
\textbf{Points de vigilance :} ne pas confondre asepsie (empêcher l'arrivée) et antisepsie (détruire les microbes présents) ; ne pas confondre antiseptique (tissus vivants) et désinfectant (objets).
\end{corrige}

\begin{corrige}[Exercice 2 — Vrai ou faux ?]
\begin{enumerate}
\item \textbf{FAUX.} Détruire les microbes déjà présents sur une plaie est le rôle de l'\textbf{antisepsie} ; l'asepsie empêche seulement l'arrivée de nouveaux microbes.
\item \textbf{VRAI.} Se laver les mains avant le soin empêche les microbes présents sur les mains d'atteindre la plaie : c'est une mesure d'\textbf{asepsie}.
\item \textbf{VRAI.} La bétadine (solution iodée) détruit les microbes sur la peau et les plaies : c'est un \textbf{antiseptique} couramment utilisé dans les centres de santé.
\item \textbf{VRAI.} Correctement utilisé, le préservatif empêche le passage des microbes (prévention des IST, dont le VIH) \textbf{et} le passage des spermatozoïdes (prévention des grossesses non désirées).
\item \textbf{VRAI.} La \textbf{stérilisation} (par exemple à l'autoclave ou par la chaleur forte) détruit tous les microbes, y compris leurs formes résistantes : un objet stérilisé ne contient aucun microbe vivant.
\end{enumerate}
\textbf{Point de vigilance :} la justification doit toujours faire appel à la définition exacte de la pratique concernée.
\end{corrige}

\begin{corrige}[Exercice 3 — Définitions]
\begin{itemize}
\item \textbf{Contamination :} pénétration de microbes (agents pathogènes) dans l'organisme ou dans un milieu. \emph{Exemple :} des microbes de la terre pénètrent dans une plaie ouverte.
\item \textbf{Infection :} multiplication des microbes dans l'organisme après la contamination, entraînant des troubles (rougeur, fièvre, pus...). \emph{Exemple :} une plaie mal soignée s'infecte et devient rouge et douloureuse.
\item \textbf{Asepsie :} ensemble des mesures qui empêchent l'arrivée des microbes dans un milieu ou sur un tissu. \emph{Exemple :} se laver les mains avant un soin, porter des gants, couvrir les aliments.
\item \textbf{Antisepsie :} destruction des microbes présents sur la peau ou sur une plaie à l'aide d'un antiseptique. \emph{Exemple :} appliquer de la bétadine sur une éraflure.
\item \textbf{Antiseptique :} substance qui détruit les microbes sur les tissus vivants (peau, plaie, muqueuse). \emph{Exemple :} l'alcool à \SI{70}{\celsius}... non : l'alcool à 70°, la bétadine, l'eau oxygénée.
\end{itemize}
\textbf{Point de vigilance :} bien distinguer contamination (entrée des microbes) et infection (multiplication des microbes dans l'organisme) : la contamination précède l'infection.
\end{corrige}

\begin{corrige}[Exercice 4 — Classer des pratiques]
\begin{center}
\begin{tabularx}{\linewidth}{|X|X|}
\hline
\rowcolor{popblueL} \textbf{Asepsie (empêcher l'arrivée des microbes)} & \textbf{Antisepsie (détruire les microbes présents)} \\
\hline
Porter des gants propres pour faire un pansement (barrière entre les mains et la plaie) & Appliquer de la bétadine sur une plaie (destruction des microbes de la plaie) \\
\hline
Utiliser un préservatif (barrière empêchant le passage des microbes) & Nettoyer la peau à l'alcool avant une piqûre (destruction des microbes de la peau) \\
\hline
Conserver les aliments couverts, bien cuire la viande... & Faire bouillir l'eau de boisson (la chaleur détruit les microbes présents dans l'eau) \\
\hline
\end{tabularx}
\end{center}
\textbf{Classement détaillé et justifié :}
\begin{enumerate}
\item \textbf{Faire bouillir l'eau : antisepsie} (au sens large de destruction des microbes) : l'ébullition tue les microbes déjà présents dans l'eau.
\item \textbf{Bétadine sur une plaie : antisepsie} : destruction des microbes présents sur la plaie.
\item \textbf{Gants propres : asepsie} : ils empêchent les microbes des mains d'atteindre la plaie.
\item \textbf{Préservatif : asepsie} : barrière qui empêche l'arrivée des microbes, sans les détruire.
\item \textbf{Bien cuire la viande : antisepsie} (destruction) : la cuisson tue les microbes présents dans l'aliment.
\item \textbf{Alcool sur la peau : antisepsie} : destruction des microbes présents sur la peau avant la piqûre.
\end{enumerate}
\textbf{Point de vigilance :} le critère de classement est le suivant : si le geste \emph{tue} des microbes déjà présents, c'est de l'antisepsie ; s'il \emph{empêche leur arrivée}, c'est de l'asepsie. On accepte que la destruction des microbes dans l'eau ou les aliments soit rangée dans l'antisepsie au sens large (destruction des microbes présents dans un milieu).
\end{corrige}

\begin{corrige}[Exercice 5 — Le lavage des mains]
\begin{enumerate}
\item Ces gestes correspondent à l'\textbf{asepsie} : le lavage des mains au savon et le port de gants propres \textbf{empêchent l'arrivée des microbes} (présents sur les mains du soignant) sur le malade ou sur sa plaie. On ne détruit pas des microbes déjà présents sur le malade : on empêche leur apport.
\item Les mains touchent de nombreux objets souillés et transportent des microbes. En les lavant puis en portant des gants propres, l'infirmier supprime ou isole ces microbes : ils ne peuvent plus pénétrer dans les plaies des malades. Le risque de \textbf{contamination} (entrée des microbes) puis d'\textbf{infection} (multiplication des microbes) est donc fortement réduit.
\item Autres gestes d'asepsie dans un centre de santé : utiliser du matériel stérilisé (seringues à usage unique, instruments stériles) ; porter un masque et une blouse propre ; nettoyer et couvrir les plaies avec un pansement stérile ; jeter les aiguilles usagées dans un conteneur prévu à cet effet.
\end{enumerate}
\end{corrige}

\begin{corrige}[Exercice 6 — Compte rendu de TP : mains lavées et mains non lavées]
\begin{enumerate}
\item Schéma des deux boîtes après 48 h à l'étuve :
\begin{popfigure}
\begin{tikzpicture}
% Boîte A
\draw[thick,popdark] (0,0) circle (1.6);
\foreach \p in {(0.4,0.5),(-0.6,0.3),(0.9,-0.4),(-0.3,-0.8),(0.1,1.0),(-1.0,-0.3),(0.7,0.9),(-0.8,0.9),(1.1,0.3),(-0.2,-0.2),(0.5,-1.0),(-1.2,0.5)}
  \fill[poporange] \p circle (0.09);
\node[below,font=\small\bfseries] at (0,-1.8) {Boîte A : main non lavée};
% Boîte B
\draw[thick,popdark] (5,0) circle (1.6);
\fill[poporange] (4.6,0.4) circle (0.09);
\fill[poporange] (5.5,-0.5) circle (0.09);
\node[below,font=\small\bfseries] at (5,-1.8) {Boîte B : main lavée au savon};
\end{tikzpicture}
\legende{Boîtes de Pétri après 48 h à \SI{37}{\celsius}. Chaque point orange représente une colonie de microbes développée sur la gélose. Boîte A : nombreuses colonies ; boîte B : très peu de colonies.}
\end{popfigure}
\item \textbf{Hypothèse :} le lavage des mains au savon élimine les microbes présents sur les mains (ou : les mains non lavées portent des microbes que le lavage supprime).
\item \textbf{Interprétation :} la boîte A, ensemencée par la main non lavée, présente de nombreuses colonies : la main non lavée portait donc de nombreux microbes, qui se sont multipliés sur la gélose à \SI{37}{\celsius}. La boîte B, ensemencée par la main lavée, ne présente que très peu de colonies : le lavage au savon a éliminé la quasi-totalité des microbes des mains. La variable testée est le lavage ; toutes les autres conditions (même gélose, même température, même durée) sont identiques, donc la différence observée est bien due au lavage.
\item \textbf{Conclusion :} le lavage soigneux des mains au savon supprime les microbes présents sur les mains et empêche ainsi leur arrivée sur un milieu ou sur une plaie : c'est une pratique d'\textbf{asepsie}, efficace pour éviter la contamination.
\end{enumerate}
\textbf{Points de vigilance :} l'hypothèse doit être formulée avant l'interprétation ; l'interprétation doit comparer les deux boîtes et citer les conditions identiques (méthode comparative) ; la conclusion doit répondre à l'hypothèse et nommer la pratique.
\end{corrige}

\begin{corrige}[Exercice 7 — Une blessure au genou]
\begin{enumerate}
\item Conduite à tenir, étape par étape :
\begin{enumerate}
\item Se \textbf{laver les mains} à l'eau et au savon (ou porter des gants propres) avant de toucher la plaie.
\item \textbf{Nettoyer la plaie} à l'eau propre et au savon pour éliminer la terre et les saletés.
\item \textbf{Rincer} à l'eau propre et sécher en tamponnant avec une compresse propre, sans frotter.
\item \textbf{Appliquer un antiseptique} (bétadine, alcool à 70° ou eau oxygénée) sur la plaie.
\item \textbf{Couvrir la plaie} avec un pansement propre (stérile si possible).
\item Consulter un centre de santé, car le clou rouillé présente un risque particulier (plaie souillée, profondeur possible).
\end{enumerate}
\item Nature de chaque pratique :
\begin{itemize}
\item Lavage des mains, gants propres, pansement propre : \textbf{asepsie}, car ces gestes empêchent l'arrivée de nouveaux microbes sur la plaie.
\item Nettoyage de la plaie et application de l'antiseptique : \textbf{antisepsie}, car ils détruisent les microbes déjà présents dans la plaie souillée.
\end{itemize}
\item Laisser la plaie souillée sans soin est dangereux car les microbes apportés par la terre et le clou pénètrent dans l'organisme (\textbf{contamination}) puis s'y multiplient (\textbf{infection}) : la plaie devient rouge, douloureuse, avec du pus, et l'infection peut s'aggraver et se généraliser. Un clou rouillé et une plaie souillée de terre exposent en outre à des infections graves : il faut donc nettoyer, désinfecter et consulter rapidement.
\end{enumerate}
\end{corrige}

\begin{corrige}[Exercice 8 — L'eau de boisson au village]
\begin{enumerate}
\item \textbf{Faire bouillir l'eau} relève de l'\textbf{antisepsie} (au sens large) : l'ébullition détruit les microbes déjà présents dans l'eau du puits. \textbf{Conserver l'eau dans des récipients propres et fermés} relève de l'\textbf{asepsie} : le récipient propre et fermé empêche l'arrivée de nouveaux microbes dans l'eau après son traitement.
\item Après ébullition, l'eau ne contient pratiquement plus de microbes, mais l'air, la poussière, les mains ou un récipient sale peuvent en apporter de nouveaux. Si l'eau est conservée dans un récipient ouvert ou sale, elle est de nouveau \textbf{contaminée} et l'ébullition devient inutile. Le récipient propre et fermé maintient donc l'eau à l'abri des microbes : asepsie et antisepsie sont \textbf{complémentaires}.
\end{enumerate}
\textbf{Point de vigilance :} bien montrer la complémentarité des deux pratiques : détruire les microbes présents (antisepsie) ne suffit pas si l'on n'empêche pas ensuite leur retour (asepsie).
\end{corrige}

\begin{corrige}[Exercice 9 — Type BEPC : l'expérience de Semmelweis]
\begin{enumerate}
\item Calcul des taux de mortalité :
\begin{itemize}
\item Avant le lavage des mains : $\dfrac{400}{\num{4000}} \times 100 = \num{0,1} \times 100 = \textbf{10 \%}$.
\item Après le lavage des mains : $\dfrac{35}{\num{3500}} \times 100 = \num{0,01} \times 100 = \textbf{1 \%}$.
\end{itemize}
\item Variation du taux de mortalité : $10\,\% - 1\,\% = \textbf{9 points de pourcentage}$ de baisse. (Le taux a été divisé par 10.)
\item Le lavage des mains avec une solution \textbf{antiseptique} détruit les microbes présents sur les mains du personnel soignant : c'est une pratique d'\textbf{antisepsie} (destruction des microbes présents sur un tissu vivant, la peau). Ce geste a ensuite un effet d'asepsie vis-à-vis des patientes, puisqu'il empêche l'arrivée des microbes lors de l'accouchement.
\item \textbf{Conclusion :} avant l'introduction du lavage des mains antiseptique, une accouchée sur dix mourait de la fièvre puerpérale (taux de mortalité de 10 \%). Après l'imposition de ce geste par Semmelweis, le taux tombe à 1 \%, soit une baisse de 9 points : le nombre de décès est divisé par dix. Le lavage des mains à la solution antiseptique est donc une pratique très efficace qui, en détruisant les microbes transportés par les mains des soignants, empêche la contamination des patientes et sauve de nombreuses vies.
\end{enumerate}
\textbf{Barème indicatif (sur 10) :} calcul du taux avant (2 pts) ; calcul du taux après (2 pts) ; variation en points (2 pts) ; identification de l'antisepsie justifiée (2 pts) ; conclusion rédigée exploitant les résultats chiffrés (2 pts).

\textbf{Points de vigilance :} appliquer la formule du taux avec les bonnes valeurs et exprimer le résultat en \% ; la variation se calcule en \emph{points de pourcentage} (soustraction), pas en pourcentage relatif ; la conclusion doit citer les résultats calculés et nommer la pratique (antisepsie) ; ne pas oublier que le geste porte sur les mains du soignant (tissu vivant), d'où l'antisepsie.
\end{corrige}

\begin{corrige}[Exercice 10 — Type BEPC : situation-problème au centre de santé]
\begin{enumerate}
\item \textbf{Soins pour l'enfant blessé :}
\begin{itemize}
\item \emph{Gestes d'asepsie} (empêcher l'arrivée de nouveaux microbes) : l'infirmière se lave les mains au savon et enfile des gants propres ; elle utilise du matériel propre ou stérile (compresses, pansement) ; elle couvre la plaie d'un pansement propre.
\item \emph{Gestes d'antisepsie} (détruire les microbes présents) : elle nettoie la plaie à l'eau propre et au savon pour éliminer les saletés, puis applique un antiseptique (bétadine) sur la plaie.
\end{itemize}
Justification : les gestes d'asepsie créent une barrière entre les sources de microbes (mains, matériel sale) et la plaie ; les gestes d'antisepsie détruisent les microbes déjà entrés dans la plaie lors de la blessure.
\item Entre deux patients, les mains et les gants de l'infirmière peuvent transporter les microbes du premier malade. Se laver les mains et changer de gants \textbf{empêche le transport des microbes d'un patient à l'autre} : cela évite la contamination croisée (par exemple entre la plaie de l'enfant et la femme enceinte, plus fragile).
\item \textbf{Double rôle du préservatif :} il forme une barrière qui empêche le passage des \textbf{microbes} d'un partenaire à l'autre, donc il protège des \textbf{IST} (dont le VIH) ; il empêche aussi le passage des \textbf{spermatozoïdes}, donc il protège des \textbf{grossesses précoces ou non désirées}. Il joue ainsi un double rôle de prévention, à condition d'être correctement utilisé à chaque rapport.
\item \textbf{Message de sensibilisation (exemple) :} « Les infections sexuellement transmissibles, comme le VIH, la blennorragie ou la syphilis, passent d'une personne à l'autre lors de rapports sexuels non protégés. Pour te protéger et protéger les autres, utilise un préservatif à chaque rapport sexuel : c'est une barrière qui empêche le passage des microbes. Le préservatif protège aussi des grossesses précoces. En cas de doute ou de signe d'infection, va vite consulter dans un centre de santé : beaucoup d'IST se soignent si elles sont prises à temps. Protège ta santé et ton avenir ! »
\end{enumerate}
\textbf{Barème indicatif (sur 10) :} soins distinguant asepsie/antisepsie justifiés (3 pts) ; explication du lavage des mains et du changement de gants (2 pts) ; double rôle du préservatif (2 pts) ; message de sensibilisation clair et correct (3 pts).

\textbf{Points de vigilance :} exiger la justification de chaque geste par la définition (empêcher l'arrivée / détruire les microbes présents) ; le message doit citer le préservatif, les IST et la consultation ; rester dans un ton adapté et scientifiquement exact.
\end{corrige}

\begin{corrige}[Exercice 11 — Type BEPC : étude documentaire sur les IST au Cameroun]
\begin{enumerate}
\item \textbf{Contamination :} pénétration de microbes (agents pathogènes) dans l'organisme, par exemple lors d'un rapport sexuel non protégé avec un partenaire contaminé. \textbf{Infection :} multiplication des microbes dans l'organisme après la contamination, responsable de la maladie et de ses complications.
\item Lors d'un rapport sexuel, le préservatif recueille les sécrétions et constitue une \textbf{barrière physique} : les microbes présents chez un partenaire contaminé ne peuvent pas atteindre l'autre partenaire. Il \textbf{empêche l'arrivée des microbes} sans les détruire : son utilisation correspond donc à une pratique d'\textbf{asepsie}.
\item Deux IST mentionnées dans le document : la \textbf{blennorragie} et la \textbf{syphilis} (on peut aussi citer l'infection à VIH). Une complication possible d'une IST non soignée : la \textbf{stérilité}.
\item \textbf{Oui.} Le préservatif empêche également le passage des \textbf{spermatozoïdes} vers les organes génitaux de la partenaire : la fécondation ne peut pas avoir lieu. Correctement utilisé, il protège donc aussi des grossesses non désirées, en plus de la protection contre les IST.
\item \textbf{Conclusion :} les pratiques de prévention sont complémentaires et indispensables à la santé de chacun. L'\textbf{asepsie} (lavage des mains, gants, matériel propre, préservatif) empêche l'arrivée des microbes ; l'\textbf{antisepsie} (antiseptiques sur la peau et les plaies) détruit les microbes déjà présents ; le \textbf{préservatif}, pratique d'asepsie, protège à la fois des IST et des grossesses non désirées. En appliquant ces gestes simples au quotidien et en consultant un centre de santé en cas de besoin, chacun évite la contamination, l'infection et leurs complications parfois graves, comme la stérilité.
\end{enumerate}
\textbf{Barème indicatif (sur 10) :} définitions exactes (2 pts) ; rôle de barrière du préservatif et classement en asepsie justifié (3 pts) ; deux IST et une complication citées (2 pts) ; protection contre les grossesses justifiée (1 pt) ; conclusion dégageant la complémentarité des trois pratiques (2 pts).

\textbf{Points de vigilance :} la justification du classement du préservatif doit insister sur « barrière qui empêche le passage des microbes » (et non destruction) ; les réponses aux questions 3 doivent être tirées du document ; la conclusion doit mobiliser les trois pratiques du programme (asepsie, antisepsie, préservatif) et rester rédigée en phrases.
\end{corrige}

\newpage

\coursec{Fiche bilan}

\begin{popfigure}
\centering
\begin{tikzpicture}[
  mindmap,
  grow cyclic,
  every node/.style={concept, execute at begin node=\hskip0pt},
  concept color=popblue,
  level 1/.append style={level distance=4.5cm, sibling angle=60, font=\small\bfseries},
  level 2/.append style={level distance=3cm, sibling angle=45, font=\footnotesize},
  branch/.style={concept color=#1}
]
\node[concept, text width=3.2cm, align=center, font=\bfseries\normalsize, fill=popdark!90!popblue] {Prévention de la\\ contamination et\\ de l'infection}
  child[branch=popgreen] { node[concept, text width=2.5cm, align=center] {1. Rappels\\ Contamination $\neq$ Infection\\ Portes d'entrée\\ Maladies infectieuses}
    child { node[concept, text width=2cm, align=center] {Contamination :\\ arrivée microbes} }
    child { node[concept, text width=2cm, align=center] {Infection :\\ multiplication + signes} }
  }
  child[branch=poporange] { node[concept, text width=2.5cm, align=center] {2. Asepsie\\ Empêcher l'arrivée\\ Stérilisation (121°C, 20 min)\\ Champ opératoire\\ Gants/Champs stériles}
    child { node[concept, text width=2cm, align=center] {Chaleur humide/ sèche} }
    child { node[concept, text width=2cm, align=center] {Filtration, UV, Gaz} }
    child { node[concept, text width=2cm, align=center] {Règles bloc opératoire} }
  }
  child[branch=poppurple] { node[concept, text width=2.5cm, align=center] {3. Antisepsie\\ Détruire sur vivant\\ Peau / Muqueuses\\ Bétadine, Dakin, Alcool 70°\\ Chlorhexidine}
    child { node[concept, text width=2cm, align=center] {Antiseptiques chimiques} }
    child { node[concept, text width=2cm, align=center] {Concentration / Temps contact} }
    child { node[concept, text width=2cm, align=center] {Ne pas mélanger} }
  }
  child[branch=poppink] { node[concept, text width=2.5cm, align=center] {4. TP N°6\\ Pratique asepsie/antisepsie\\ Lavage mains simple/chirurgical\\ Désinfection matériel\\ Mise en place champ stérile}
    child { node[concept, text width=2cm, align=center] {Technique friction hydroalcoolique} }
    child { node[concept, text width=2cm, align=center] {Gantage stérile} }
    child { node[concept, text width=2cm, align=center] {Évaluation stérilité} }
  }
  child[branch=popgold] { node[concept, text width=2.5cm, align=center] {5. Préservatifs\\ Barrière physique\\ M (latex) / F (nitrile)\\ VIH / IST / Contraception\\ Conservation / Date / Norme CE}
    child { node[concept, text width=2cm, align=center] {Masculin : mise en place} }
    child { node[concept, text width=2cm, align=center] {Féminin : anneaux souples} }
    child { node[concept, text width=2cm, align=center] {Lubrifiant à base d'eau} }
  }
  child[branch=popblue] { node[concept, text width=2.5cm, align=center] {6. Vie courante\\ Hygiène eau/aliments\\ Vaccination (PEV)\\ Gestion déchets/ordures\\ Lutte vecteurs (moustiques)}
    child { node[concept, text width=2cm, align=center] {Eau traitée / bouillie} }
    child { node[concept, text width=2cm, align=center] {Cuisson complète viandes} }
    child { node[concept, text width=2cm, align=center] {Moustiquaires imprégnées} }
  };
\end{tikzpicture}
\legende{Carte mentale récapitulative de la leçon 5 : pratiques de prévention (asepsie, antisepsie, préservatifs, hygiène).}
\end{popfigure}

\begin{bilan}

\courssub{Définitions clés à connaître par cœur}

\begin{definition}[Contamination]
Arrivée et dépôt de micro-organismes (bactéries, virus, champignons, parasites) sur un milieu inerte (objets, instruments, surfaces) ou vivant (peau, muqueuses saines). Elle est \textbf{silencieuse} (pas de signes cliniques).
\end{definition}

\begin{definition}[Infection]
Pénétration, multiplication et développement d'un agent pathogène dans l'organisme hôte, entraînant une réaction de défense (inflammation, fièvre, pus) et des signes cliniques. Toute infection est précédée d'une contamination, l'inverse n'est pas vrai.
\end{definition}

\begin{definition}[Asepsie]
Ensemble des techniques visant à \textbf{empêcher l'arrivée} de tout micro-organisme sur un milieu stérile (blessure, matériel chirurgical, milieu de culture). Elle s'applique sur du \textbf{vivant sain} (peau intacte avant incision) et de l'\textbf{inerte}. Objectif : \textbf{stérilité totale} (0 germe).
\end{definition}

\begin{definition}[Antisepsie]
Ensemble des techniques visant à \textbf{détruire ou inhiber} les micro-organismes \textbf{sur un milieu vivant} (peau lésée, muqueuses, plaie) à l'aide de produits chimiques (antiseptiques) tolérés par les tissus. Objectif : \textbf{réduire la charge microbienne} (pas de stérilité absolue).
\end{definition}

\begin{definition}[Stérilisation]
Opération physique ou chimique qui détruit \textbf{tous} les micro-organismes (bactéries, spores, virus) sur un matériel inerte. Référence : \textbf{chaleur humide} (autoclave) \SI{121}{\celsius} pendant \SI{20}{\minute} (ou \SI{134}{\celsius} \SI{3}{\minute}).
\end{definition}

\begin{definition}[Désinfection]
Opération qui détruit les germes \textbf{végétatifs} (pas les spores obligatoirement) sur matériel inerte ou surfaces. Niveau inférieur à la stérilisation.
\end{definition}

\begin{definition}[Préservatif (Masculin / Féminin)]
Dispositif médical à usage unique constituant une \textbf{barrière physique étanche} (latex ou nitrile/polyuréthane) empêchant le contact des muqueuses et des liquides biologiques (sperme, sécrétions vaginales, sang). Protège contre les \textbf{IST (VIH, syphilis, HPV...)} et assure une \textbf{contraception} efficace (si bien utilisé).
\end{definition}

\courssub{Tableau comparatif : Asepsie vs Antisepsie (à reproduire)}

\begin{center}
\begin{tabularx}{\linewidth}{|l|>{\raggedright\arraybackslash}X|>{\raggedright\arraybackslash}X|}\hline
\rowcolor{popblueL}
\textbf{Critère} & \textbf{Asepsie} & \textbf{Antisepsie} \\ \hline
\textbf{But} & Empêcher l'arrivée (stérilité) & Détruire / inhiber (réduction flore) \\ \hline
\textbf{Support} & Inerte + Vivant sain (peau intacte) & Vivant lésé (plaie, muqueuse, peau saine pré-op) \\ \hline
\textbf{Moyens} & Physiques (chaleur, filtration, UV) + Mécaniques (gants, champs, masques) & Chimiques (Antiseptiques : Iode, Chlorhexidine, Alcool, Dakin) \\ \hline
\textbf{Exemples} & Autoclave, champ opératoire, gantage stérile, laminaire & Bétadine® cutanée, Dakin®, Alcool 70\%, Chlorhexidine 4\% \\ \hline
\textbf{Résultat visé} & 0 UFC (Unité Formant Colonie) & Diminution significative de la flore \\ \hline
\end{tabularx}
\end{center}

\courssub{Méthodes express (à savoir dérouler)}

\begin{methode}[Lavage des mains hygiénique (solution hydroalcoolique - SHA) -- \SI{20}{\second}--\SI{30}{\second}]
\begin{enumerate}
\item Appliquer \SI{3}{\milli\liter} SHA dans le creux de la main sèche.
\item Frotter paume contre paume, dos des mains, espaces interdigitaux, pouces, ongles, poignets.
\item Laisser sécher \textbf{sans essuyer} ni rincer. Mains sèches = stériles pour soin courant.
\end{enumerate}
\end{methode}

\begin{methode}[Lavage des mains chirurgical (SHA ou savon antiseptique) -- \SI{3}{\minute}--\SI{5}{\minute}]
\begin{enumerate}
\item Brosser ongles (premier lavage seulement).
\item Laver mains + avant-bras + coudes (mouillé $\rightarrow$ savon $\rightarrow$ rinçage $\rightarrow$ séchage stérile $\rightarrow$ SHA).
\item Gantage stérile immédiat (technique « fermée » ou « ouverte »).
\end{enumerate}
\end{methode}

\begin{methode}[Mise en place d'un champ opératoire stérile]
\begin{enumerate}
\item Ouvrir le champ par les coins (non stériles) sans toucher la face interne.
\item Poser sur table stérile / drap stérile. Ne pas tendre au-dessus de la zone stérile.
\item Disposer instruments stériles (ouverts par aide non stérile qui verse/présente).
\item Respecter la règle : \textbf{Stérile touche stérile ; Non stérile ne touche jamais stérile.}
\end{enumerate}
\end{methode}

\begin{methode}[Utilisation correcte du préservatif masculin]
\begin{enumerate}
\item Vérifier date péremption, norme CE/ISO, état emballage (coussin d'air).
\item Ouvrir délicatement (pas dents, pas ciseaux). Sens du décalottage : bague à l'extérieur.
\item Pincer le réservoir (chasser l'air) $\rightarrow$ décalotter sur pénis en érection.
\item Après éjaculation : retirer \textbf{avant} érection complète, tenir la bague, nouer, jeter (poubelle, pas WC).
\item \textbf{Jamais} deux préservatifs ensemble (frottement $\rightarrow$ rupture). Lubrifiant \textbf{uniquement à base d'eau}.
\end{enumerate}
\end{methode}

\courssub{Produits antiseptiques usuels (Concentrations efficaces)}

\begin{center}
\begin{tabularx}{\linewidth}{|l|X|X|}\hline
\rowcolor{poporangeL}
\textbf{Produit (DCI / Marque)} & \textbf{Indication principale} & \textbf{Précautions / Contre-indications} \\ \hline
\textbf{Povidone iodée} (Bétadine® 10\% dermique, 4\% muqueuse) & Peau saine, plaies, muqueuses, pré-op & Allergie iode, thyroïde, femme enceinte/allaitante, enfant $<$ 30 mois (muqueuse) \\ \hline
\textbf{Chlorhexidine} (4\% cutané, 0,5\% muqueux + alcool) & Peau saine, pré-op, cathéters, muqueuses & Œil, oreille (tympan percé), méninges, prématurés. Inactivé par savon/anioniques \\ \hline
\textbf{Alcool éthylique / isopropylique 60--70\%} (SHA, peau saine) & Antisepsie cutanée rapide, SHA mains & Peau lésée/muqueuses (brûlure), inflammable, dessèchement \\ \hline
\textbf{Hypochlorite de sodium} (Eau de Javel / Dakin® \SI{0,1}{\%}--\SI{0,5}{\%} Cl actif) & Plaies sales, nécrotiques, matériel inerte (dilution) & Instable (lumière/chaleur), inactivé par matière organique, irritant \\ \hline
\textbf{Hexamidine} (Hexomédine®) & Plaies superficielles, brûlures, ombilic & Moins large spectre (peu antiviral) \\ \hline
\end{tabularx}
\end{center}

\courssub{Points de vigilance (Pièges BEPC)}

\begin{attention}[Confusion Stérilisation / Désinfection / Antisepsie]
\begin{itemize}
\item \textbf{Stérilisation} = 0 germe (spores comprises) $\rightarrow$ \textbf{Inerte uniquement} (autoclave).
\item \textbf{Désinfection} = germes végétatifs $\rightarrow$ \textbf{Inerte / Surfaces}.
\item \textbf{Antisepsie} = produit chimique toléré $\rightarrow$ \textbf{Vivant (plaie, muqueuse)}.
\item \textbf{Asepsie} = \textbf{Démarche globale} (règles de comportement, tenue, champ) pour \textbf{garder stérile}.
\end{itemize}
\end{attention}

\begin{attention}[Antisepsie : Règles d'or]
\begin{itemize}
\item \textbf{Ne jamais mélanger} deux antiseptiques (inactivation, toxicité : ex Bétadine + Dakin = gaz toxiques / perte efficacité).
\item Respecter \textbf{concentration} et \textbf{temps de contact} (étiquette).
\item Nettoyer (eau + savon / sérum physio) \textbf{AVANT} d'antisepsier (la matière organique inactive l'antiseptique).
\item Sens de passage : du \textbf{propre vers le sale} (centre $\rightarrow$ périphérie).
\end{itemize}
\end{attention}

\begin{attention}[Préservatif : Erreurs fréquentes]
\begin{itemize}
\item Ouvrir avec les dents / ciseaux $\rightarrow$ micro-déchirures.
\item Ne pas pincer le réservoir $\rightarrow$ poche d'air $\rightarrow$ rupture à l'éjaculation.
\item Mettre après le début du rapport (pré-éjaculation = risque IST/grossesse).
\item Utiliser corps gras (vaseline, huile, beurre) $\rightarrow$ dégrade latex en \SI{30}{\second}.
\item Réutiliser / retourner / conserver au portefeuille/voiture (chaleur/frottement).
\end{itemize}
\end{attention}

\end{bilan}

\begin{aretenir}[Checklist avant le BEPC -- Leçon 5]
\begin{enumerate}
\item $\square$ Différencier clairement \textbf{contamination} (arrivée, silencieuse) et \textbf{infection} (multiplication, symptômes).
\item $\square$ Définir \textbf{asepsie} (empêcher l'arrivée, stérilité) et \textbf{antisepsie} (détruire sur vivant, produit chimique).
\item $\square$ Citer \textbf{3 méthodes de stérilisation} (chaleur humide 121°C/20min, chaleur sèche, filtration, rayonnements, gaz) et donner la référence autoclave.
\item $\square$ Citer \textbf{4 antiseptiques majeurs} (Povidone iodée, Chlorhexidine, Alcool, Dakin/Hypochlorite) avec \textbf{une indication} et \textbf{une contre-indication} chacun.
\item $\square$ Expliquer la règle \textbf{« Ne jamais mélanger deux antiseptiques »} et le sens de passage (propre $\rightarrow$ sale).
\item $\square$ Décrire les étapes du \textbf{lavage des mains chirurgical} (durée, brossage, séchage stérile, gantage) vs hygiénique (SHA, 30s).
\item $\square$ Schématiser / Décrire la \textbf{mise en place d'un champ opératoire stérile} (ouverture coins, non contact face interne, règle stérile/stérile).
\item $\square$ Lister les \textbf{étapes correctes de mise en place du préservatif masculin} (vérif, ouverture, réservoir, décalottage, retrait, élimination).
\item $\square$ Citer \textbf{3 erreurs à ne pas commettre} avec le préservatif (corps gras, double usage, ouverture dents, conservation chaleur).
\item $\square$ Citer \textbf{3 gestes de prévention quotidiens} (eau traitée, cuisson viandes, lavage mains avant manger/après WC, moustiquaire, vaccination PEV).
\item $\square$ Savoir exploiter un \textbf{protocole expérimental TP N°6} (identifier groupe témoin/témoin, variable testée, critère d'évaluation stérilité : absence de culture).
\end{enumerate}
\end{aretenir}

\findecours

\end{document}
